\documentclass[a4paper,onecolumn,superscriptaddress,10pt,unpublished]{compositionalityarticle}
\pdfoutput=1
\togglefalse{@accepted}
\makeatletter\def\@@addcommabetween#1{#1}\makeatother
\newcommand{\compositionalityVolume}{}\newcommand{\compositionalityIssue}{}
\usepackage[utf8]{inputenc}
\usepackage[english]{babel}
\usepackage[T1]{fontenc}
\usepackage{amsmath,amssymb,amsthm,amscd,booktabs,array,graphicx}
\usepackage{hyperref}
\DeclareUnicodeCharacter{0393}{\ensuremath{\Gamma}}
\emergencystretch=3em
\fancypagestyle{firststyle}{\fancyhf{}\fancyfoot[C]{\sffamily\thepage}}
\fancyhf{}
\fancyhead[L]{\sffamily\scriptsize \'A.~J.~Toranzo Portela}
\fancyhead[R]{\sffamily\scriptsize A Formal Audit of a Categorical Unification of Consciousness Theories}
\fancyfoot[C]{\sffamily\thepage}
\theoremstyle{plain}
\newtheorem{theorem}{Theorem}
\newtheorem{proposition}[theorem]{Proposition}
\newtheorem{lemma}[theorem]{Lemma}
\newtheorem{corollary}[theorem]{Corollary}
\theoremstyle{definition}
\newtheorem{definition}[theorem]{Definition}
\newtheorem{remark}[theorem]{Remark}
\newcommand{\stL}{\textsf{[L]}}
\newcommand{\stA}{\textsf{[A]}}
\newcommand{\stM}{\textsf{[M]}}
\newcommand{\stS}{\textsf{[S]}}
\newcommand{\stN}{\textsf{[N]}}

\begin{document}
\title{A Formal Audit of a Categorical Unification of Consciousness Theories: Structural Proxies, Formal Obstructions and a Positive Route}
\author{\'Angel Jos\'e Toranzo Portela}
\thanks{Preprint, version 1.1, 19 November 2026; not peer reviewed. The changes from version~1.0 are listed in \texttt{CHANGES.md} of the archive. The Lean~4 files, the SymPy and numerical scripts and the instructions to verify them are archived with this paper (Appendix~\ref{sec:archivos}) on Zenodo, under the concept DOI \href{https://doi.org/10.5281/zenodo.24349943}{10.5281/zenodo.24349943}, which resolves to the latest version; version~1.0 has the DOI \href{https://doi.org/10.5281/zenodo.24349944}{10.5281/zenodo.24349944}.}
\affiliation{Independent Researcher in Theoretical Physics and Cognitive Science}
\maketitle\hypersetup{pdfcreator={LaTeX with hyperref}}
\begin{abstract}
Several theories of consciousness, among them integrated information theory (IIT), global neuronal workspace theory (GNWT) and attention schema theory, together with general frameworks that have been brought to bear on consciousness, such as the free-energy principle (FEP) and quantum Darwinism, describe the same open physical systems through different invariants. A candidate theory, Q, formulated by the author and defined in this paper, proposes to unify them categorically. We audit this unification formally and structurally, to delimit with mathematical precision which parts of it can be established and which promises of unification fail, for explicit structural proxies of the theories, not for the theories themselves. The test case is the unity-sufficiency obligation $S_u$ of Q; the audit does not change the descriptor, the domain or the meaning of Q. Results are labelled by their evidence: proofs checked by the Lean~4 kernel, paper proofs, exact computation, or numerics. (I)~Negative delimitation. The weak categorical thesis holds, but cheaply: every isomorphism-invariant assignment is a functor out of a symmetric monoidal groupoid of open systems. For the presence proxies of IIT type (a Gaussian integration functional $J>0$, not the $\Phi$ of IIT~3.0) and GNWT type (a broadcast hub of the undirected coupling graph, a structural stand-in that models neither ignition nor dynamic broadcast) there is no translation between the projections, their bridge square does not commute, and neither is monoidal into the additive reals: both vanish on systems with disconnected support, in particular on tensor products of nonempty systems, and with the empty system sent to an adjoined unit they are strong monoidal into a zero-product monoid, whereas predictive information is additive. The two proxies disagree exactly on a band of connected systems without a hub, empty for $n\le3$ units and containing the path $P_n$ and the cycle $C_n$ for every $n\ge4$. A premise ``the degree depends only on $\kappa$'' implies $S_u$ exactly when the kernel of the descriptor lies in the equivalence it generates, and none of the principles examined closes $S_u$ without presupposing it, equating only isomorphic realizations, forcing a constant degree, changing the meaning of the degree, or changing the domain. A conditional theorem extends this criterion to the sufficiency obligation of every component of Q, to its order clause and to its descent condition: it states which premises would discharge each of these obligations, and justifies none of them. (II)~A Gaussian laboratory, exact in a two-unit family and numerical in four-unit networks, illustrates where the bridges commute and where they break, and shows exactly that the min-cut lower bound on $J$ is attained. (III)~Of the reformulations examined, only intervention monotonicity escapes the collapse; it is a declared weakening of the order clause of Q, not $S_u$, and we have not shown that it is the only one. Degrees are forced to be constant exactly on the strongly connected components of an intervention digraph. With $J$ as a stand-in for the descriptor of Q, on systems with at least two units, $J$ increases strictly under uniform scaling of any connected dynamics, or with a fixed scalar self-coupling (checked in Lean with Mathlib, as are the equivalence of $J>0$ with connected support and the min-cut bounds on $J$), but numerically it can decrease against a fixed non-scalar self-coupling. Q is neither demonstrated nor refuted. Seven Lean files without Mathlib and a Lean project with Mathlib of eight modules, all without \texttt{sorry}, and twelve scripts archived with the paper reproduce every result labelled as checked by Lean, by exact computation or numerically.
\end{abstract}

\section{Introduction}\label{sec:intro}
Several scientific theories of consciousness, and some general frameworks that have been brought to bear on it, describe the same open physical systems through different invariants. Integrated information theory (IIT) looks at how much a system's causal structure resists being cut~\cite{Oizumi2014}; global neuronal workspace theory (GNWT) looks at broadcast and global availability~\cite{DehaeneChangeux2011}; attention schema theory (AST) looks at an internal model of attention~\cite{GrazianoKastner2011}; the free-energy principle (FEP), an account of self-organizing systems, looks at Markov blankets and the statistics of the internal states that they shield~\cite{Friston2013}; and quantum Darwinism, an account of how classical records emerge, looks at redundant records in the environment~\cite{Zurek2009}. Category theory has already been used to study single theories and their relation to experience: the relation between consciousness and IIT~\cite{Northoff2019,Tsuchiya2016}; IIT itself, formulated in process theories, that is, in symmetric monoidal categories~\cite{TullKleiner2021}, on the basis of an axiomatic account of its mathematical structure~\cite{KleinerTull2021}; and graphical reasoning about conscious experience in process theories~\cite{SignorelliWangCoecke2021}. A further proposal is to unify several theories at once: to treat each one as a functor from a common category of open physical systems into its own category of descriptions, and to look for natural transformations, limits or bridge laws that connect them. This paper defines and audits a proposal of this kind, the candidate theory Q (Section~\ref{sec:obligacion}): a correspondence between open physical realizations and five components of phenomenal structure, $(c,q,u,\tau,\sigma)$. Q is the author's own proposal, so the audit is a self-audit, and it applies to Q the same standard that it applies to the proxies. Quantum-collapse proposals such as Orch-OR~\cite{HameroffPenrose2014}, which would enter Q only as additional postulates, are left aside: we define no proxy for them. The central open obligation of Q for the unity component $u$ is the unity sufficiency
\begin{equation}\label{eq:Su}
S_u:\qquad \Phi_Q(s)=\Phi_Q(t)\ \Longrightarrow\ u(s)=u(t),
\end{equation}
where $\Phi_Q$ is the physical descriptor of the theory and $s,t$ range over the original domain of realizations.

This paper is a formal and structural audit of that unification. Its aim is to delimit with mathematical precision which parts of it can be established, and which promises of unification fail, for explicit structural proxies of the theories, not for the theories themselves; the logical and combinatorial core of the argument is checked by computer in Lean~4, and its analytic core, three results on the Gaussian integration functional $J$ that serves as an IIT-type proxy, is checked in Lean~4 with the Mathlib library~\cite{Mathlib2020}. As a rule of method, the audit changes neither the descriptor, nor the domain, nor the components, nor the meaning of Q. We do not try to demonstrate Q, and we do not refute it, and the kind of evidence that supports each result is stated (Section~\ref{sec:status}). The paper pursues three aims, one per act.
\begin{itemize}
\item \emph{Act~I: negative delimitation} (Sections~\ref{sec:obligacion}--\ref{sec:debil}). First, the weak categorical thesis holds, but cheaply, and the naive strong unification fails: for the structural presence proxies of IIT type ($J>0$, with $J$ a Gaussian integration functional) and GNWT type (a broadcast hub), there is no translation between the two projections, their bridge square does not commute, and neither projection is monoidal into the additive reals. Second, the two proxies give opposite verdicts exactly on a band of connected systems without a hub, which is empty for $n\le3$ units and contains the path $P_n$ and the cycle $C_n$ for every $n\ge4$. Both vanish on systems with disconnected support, such as two disjoint edges, $2K_2$; since they vanish on tensor products of nonempty systems, they are strong monoidal into a zero-product monoid once the empty system is sent to an adjoined unit, whereas the extensive projections are additive. Third, we delimit the obligation $S_u$: we characterize exactly when a premise of the form ``the degree depends only on $\kappa$'' implies $S_u$, and we show that none of the candidate principles examined closes it without presupposing it, equating only isomorphic realizations, forcing the degree to be constant, changing the meaning of the degree, or changing the domain, as selecting the state $z=0$ of minimal surprisal does in a Gaussian family. With the degree a function of the coexperience structure, the banded segregation SEG that would be needed is non-degenerate $S_u$ itself, in testable form.
\item \emph{Act~II: validation in a Gaussian laboratory} (Section~\ref{sec:gauss}). In a two-unit linear-Gaussian family, exactly, and in four-unit networks, numerically, we confirm where the bridges commute (all systems with $n\le3$ units, and centralized systems such as the star $K_{1,3}$), where they break (the band, from $n\ge4$), and how sharp the min-cut bounds are.
\item \emph{Act~III: a positive route and a methodological guide} (Section~\ref{sec:via-futura}). Of the reformulations of $S_u$ examined, only intervention monotonicity escapes the collapse of Act~I. It is not $S_u$, nor an obligation of the original Q, but a declared weakening of its order clause B to a class of interventions on one bearer, and we have not shown that it is the only reformulation that escapes. We prove its exact frontier: on the intervention digraph, every degree that satisfies intervention monotonicity is constant on each strongly connected component, and degrees exist that vary across components. Taking $J$ as a stand-in descriptor, which is not $\Phi_Q$, we prove, and check in Lean with Mathlib, that, for at least two units, $J$ increases strictly under uniform scaling of a connected dynamics, and under scaling of the couplings with a fixed scalar self-coupling, and we show numerically that it can decrease if the fixed self-coupling is not scalar. From these results we draw a methodological guide for comparing theories by directed interventions within one system. The archive of the paper contains six Lean files without Mathlib and a Lean project with Mathlib of seven modules, all without \texttt{sorry}, and eleven symbolic and numerical scripts with their expected outputs.
\end{itemize}

\subsection{Main results, limitations and assumptions}\label{sec:summary}
\paragraph{Main results.} The status labels are explained in Section~\ref{sec:status}, and Table~\ref{tab:evidence} there lists the evidence for the principal results. Several of the formal results are elementary facts about kernels, factorizations and products: the entailment criterion (Theorem~\ref{thm:impl}), the translation criterion (\texttt{nat\_iff\_functional}), the pairing into the product (\texttt{pairing\_unique}) and the reconstruction criterion (Theorem~\ref{thm:reconstruccion}). We claim no novelty for them. What the paper adds is their application to the obligation $S_u$ of Q and to explicit proxies, the counterexamples, and the delimitation that follows.
\begin{enumerate}
\item[(R1)] \emph{Factorization premises.} Premises ``the degree depends only on $\kappa$'' imply $S_u$ if and only if $\ker\Phi_Q$ is contained in the equivalence relation they generate (Theorem~\ref{thm:impl}) \stL. If $\kappa$ is transversal to $\Phi_Q$ (Section~\ref{sec:intento}) and co-occurs connectedly with its levels, $S_u$ forces the degree to be constant (Theorem~\ref{thm:dicotomia}) \stL.
\item[(R2)] \emph{Candidate principles.} None of the principles examined in the available sources supplies an independent foundation excluding that two original realizations with the same $\Phi_Q$ differ in $u$ when $u$ is not constant (Sections~\ref{sec:principios}--\ref{sec:mejor}). The banded segregation SEG is not a new premise that would close $S_u$: with the degree a function of the coexperience structure, SEG with distinct values between groups is equivalent to non-degenerate $S_u$ (\texttt{segd\_iff}) \stL.
\item[(R3)] \emph{Categorical thesis.} The weak thesis holds as a scheme: every isomorphism-invariant assignment is a functor out of a symmetric monoidal groupoid of open systems (Proposition~\ref{prop:funtor}) \stA, and the projections pair uniquely into their product (\texttt{pairing\_unique}) \stL. It is cheap, because it holds for any family of invariants. The strong thesis fails for the structural presence bridges of IIT type ($J>0$, with $J$ a Gaussian integration functional) and GNWT type (a broadcast hub): their bridge square does not commute on the ring $C_4$ and on the path $P_4$, there is no translation from positive integration to broadcast, and joint projections have blind spots \stL; the reverse translation fails by the same criterion \stA. Lean without Mathlib encodes the IIT-type presence as connectivity of the support, and that this is $J>0$ is Lemma~\ref{lem:J}, \stA\ \stM\ for at least two units. The two presence projections vanish on tensor products of nonempty systems (Proposition~\ref{prop:monoidal} \stA; \texttt{tensor\_no\_universal} \stL), so they are not monoidal into the additive monoid $([0,+\infty],+,0)$ (Section~\ref{sec:proyecciones}). With the convention that the empty system goes to an adjoined unit and no nonempty system does, they are strong monoidal into a zero-product monoid (\texttt{strong\_monoidal\_zero}, in abstract form) \stL; without it they are not. The extensive projections (predictive information, and the mutual information of the FEP-type projection in the Gaussian case) are additive \stA.
\item[(R4)] \emph{The band.} The two presence bridges disagree exactly on connected systems without a broadcast hub: the Boolean form is compiled (\texttt{noncomm\_locus}, \texttt{commute\_off\_band}, \texttt{universal\_connected}) \stL, and the identification of $J>0$ with connectivity is Lemma~\ref{lem:J}, \stA\ \stM\ for at least two units. There is no such system with $n\le3$ units (\texttt{n2\_equiv}, \texttt{n3\_equiv}) \stL. The path $P_n$ is one for every $n\ge4$ (\texttt{path\_reach\_from\_zero}, \texttt{path\_no\_universal}) \stL, and so is the cycle $C_n$, whose maximum degree $2$ is smaller than $n-1$ \stA. For $n=4$ the band has $15$ labeled graphs (\texttt{n4\_count}) \stL, and for $n=5$ it has $472$ (exhaustive enumeration, \texttt{pistas\_numerico.py}) \stS.
\item[(R5)] \emph{Bounds.} For systems with at least two units, $J>0$ if and only if the undirected support is connected (Lemma~\ref{lem:J}, \texttt{QAudit.\allowbreak lemmaJ\_stationary}), and $J$ is sandwiched by the minimum cut $w_{\min}$, for each representative, since $w_{\min}$ is not an isomorphism invariant (Theorem~\ref{thm:cheeger}, \texttt{QAudit.\allowbreak sandwich}) \stA\ \stM. A broadcast hub of strength $c$ gives an explicit lower bound \stA. The lower bound of the sandwich is attained in a subfamily, so its constant cannot be improved \stS.
\item[(R6)] \emph{Gaussian laboratory.} In the two-unit families $F_0$ and $F_{\rm mem}$ of Section~\ref{sec:gauss} everything is exact ($156$ symbolic checks) \stS: $F_0$ and $F_0\otimes F_0$ realize the two levels on which the bridges commute; predictive information equals $J$ and the cut information vanishes; selecting the state of minimal surprisal, a free-energy-type selection, is trivial at the level of laws and discards almost every realization at the level of states and histories, since its minimizer, the mode, has probability zero; and on $F_0$ the premise IPS$_K$ (Definition~\ref{def:IPS}), that the degree depends only on the parameter $k=c^2s_2/s_1$ of the family, is equivalent to the sufficiency $S_J$ that it would be used to support. The four-unit networks are numerical \stN.
\item[(R7)] \emph{Positive route.} Exact covariance under perturbations and the order clause on all pairs collapse \stL. Intervention monotonicity does not: it forces the degree to be constant exactly on the strongly connected components of an intervention digraph (Theorem~\ref{thm:frontier}) \stL; in Lean the degrees take integer values, and the transcription to real values is not formalized (Section~\ref{sec:via-futura}). It is not $S_u$, nor an obligation of the original Q, but a declared weakening of the order clause B to a class of interventions on one bearer, and we have not shown that it is the only reformulation that escapes the collapse. In the band it can be instantiated through Proposition~\ref{prop:scaling} if $J$ is taken in place of $\Phi_Q$, which is not justified here: for systems with at least two units, $J(tA)$ increases strictly under uniform scaling for every connected support, and so does $J(aI+cW)$ in $c$ for a connected $W$ with a fixed scalar self-coupling $a$ (\texttt{QAuditScaling.\allowbreak prop\_scaling\_a}, \texttt{QAuditScaling.\allowbreak prop\_scaling\_b}) \stA\ \stM. Strengthening only the couplings while a non-scalar diagonal of self-couplings is held fixed can decrease $J$ \stN.
\end{enumerate}

\paragraph{Limitations.}
\begin{itemize}
\item \emph{Proxies.} $J$ is the Gaussian residual functional of geometric integrated information~\cite{Oizumi2016PNAS}, not the $\Phi$ of IIT~3.0~\cite{Oizumi2014}, and the broadcast hub is a structural stand-in for global availability, not the nonlinear ignition of GNWT. The disagreement on the band is a disagreement between these proxies computed on the same model and at the same grain, not between the theories as such.
\item \emph{Formalization.} Six Lean files use Lean~4 without Mathlib and check logical skeletons and finite combinatorics \stL. A Lean project with Mathlib checks three analytic statements about $J$ over real matrices \stM: Lemma~\ref{lem:J} and the min-cut sandwich (Theorem~\ref{thm:cheeger}) for $n\ge2$ units, the sandwich exactly as stated, with the extreme eigenvalues, and both parts of Proposition~\ref{prop:scaling}, the latter slightly more generally ((a) without assuming $\operatorname{spr}(A)<1$, (b) without assuming that $W$ has zero diagonal). In Lean the inner minimum over $\mathfrak B_\pi$ (Lemma~\ref{lem:J}) is an infimum, whose attainment is a separate checked theorem, and the stationary covariance is either given as data ($P\succ0$ with $P=APA^\top+E$, Section~\ref{sec:base}) or constructed as the series $\sum_{j\ge0}A^jE(A^j)^\top$, whose convergence is proved from $\operatorname{spr}(A)<1$, where $\operatorname{spr}$ is Mathlib's spectral radius of the complexified matrix. Not formalized: the conventions for fewer than two units (Remark~\ref{rem:small}), such as $J=+\infty$ for one unit; the derivative identity $Q'=AQ'A^\top+S$ for the matrix $Q(c)=c^2P(c)$ of the proof of Proposition~\ref{prop:scaling}(b) and identity~\eqref{eq:S-identity}, in whose place the Lean proof of Proposition~\ref{prop:scaling}(b) uses the inequality $Q(c_2)-Q(c_1)\succeq(c_2-c_1)^2E$, so that~\eqref{eq:S-identity} remains \stS\ for $3\times3$ matrices and \stA\ for general $n$; the uniqueness of the stationary covariance, which the Lean statements do not need; the vanishing of $J$ on tensor products (Proposition~\ref{prop:monoidal}); and Corollary~\ref{cor:band-scaling} as a single statement, which combines graph facts \stL\ with Proposition~\ref{prop:scaling}. The files without Mathlib and the Mathlib project are separate developments. The Lean statements of Section~\ref{sec:via-futura} on exact covariance, the order clause and the intervention digraph take degrees with integer values; their transcription to real values is not formalized. The hub bound (Theorem~\ref{thm:laxa}(d)), the additivity of the predictive information $I_{\rm pred}$ (Proposition~\ref{prop:monoidal}), Proposition~\ref{prop:funtor} and the other analytic statements remain paper proofs \stA, the Gaussian computations are symbolic checks in SymPy \stS, and the numerical results \stN\ are not certificates.
\item \emph{Self-audit.} Q is the author's own proposal. Every definition, hypothesis and result that the audit uses is stated in this paper (Section~\ref{sec:obligacion}), so the audit can be checked from the paper and its archive alone.
\item \emph{Q itself.} The original descriptor $\Phi_Q$, the original domain and the comparison relation for the degree are decisions of the author that have not been taken; the original obligations of Q remain open, and nothing here measures $u$.
\item \emph{Review.} All reviews of this work were carried out by AI agents (see the \hyperref[sec:ai]{declaration on the use of generative AI}); there has been no independent human review.
\end{itemize}

\paragraph{Assumptions.} The linear-Gaussian instance $\mathbf{Open}_G$ of the category of open systems (Section~\ref{sec:base}); the undirected-support reading of the broadcast hub; and the definitions of Section~\ref{sec:obligacion} as the framework audited.

\subsection{Status labels and verification}\label{sec:status}
The results of Section~\ref{sec:summary} and of Sections~\ref{sec:obligacion}--\ref{sec:via-futura} carry one or more of five labels; in Sections~\ref{sec:obligacion}--\ref{sec:debil} a label may be implicit, as explained below.
\begin{itemize}
\item \stL: checked by the kernel of Lean~4 (version~4.24.0, without Mathlib): logic and finite combinatorics. Each of the six Lean files compiles with exit code~0, contains no \texttt{sorry}, \texttt{admit}, \texttt{native\_decide}, \texttt{axiom} or \texttt{opaque} declarations, and depends at most on the standard axioms \texttt{propext}, \texttt{Classical.choice} and \texttt{Quot.sound}; the axioms of each theorem are listed in Appendix~\ref{sec:archivos}.
\item \stM: checked by the kernel of Lean~4 (version~4.24.0) with Mathlib~\cite{Mathlib2020} (v4.24.0) over real matrices. Each of the seven modules of the Lean project compiles with exit code~0 and contains no \texttt{sorry}, \texttt{admit}, \texttt{native\_decide}, \texttt{axiom} or \texttt{opaque} declarations, and every one of its $73$ \texttt{\#print axioms} reports is \texttt{[propext, Classical.choice, Quot.sound]} (Appendix~\ref{sec:archivos}). Where the Lean statement differs in scope from the statement of the paper, the difference is stated where the label is used and in Section~\ref{sec:summary} (\emph{Limitations}).
\item \stA: proved on paper in this paper. Alone, the label means that the result is not formalized; \stA\ \stM\ marks a result whose paper proof is kept and whose statement is also checked with Mathlib.
\item \stS: verified exactly by computer: symbolic identities and exact rational arithmetic in SymPy~1.14~\cite{SymPy2017} for the Gaussian model, or exact integer enumeration in Python.
\item \stN: numerical, without a certificate (local optimization with restarts, or high-precision root finding).
\end{itemize}
Where a result in Sections~\ref{sec:obligacion}--\ref{sec:debil} carries no explicit label, the label is implicit: a statement that cites a Lean name in \texttt{typewriter} font is checked by Lean in the abstract form stated next to it; Lean names whose namespace begins with \texttt{QAudit} belong to the Mathlib project and are \stM; the others belong to the files without Mathlib and are \stL. When a statement concerns $J$, the link between Lean's Boolean predicates and $J$ is Lemma~\ref{lem:J}, which is \stA\ \stM\ for at least two units; the files without Mathlib and the Mathlib project are separate developments, and this link is not itself a compiled statement. Lemma~\ref{lem:J} and Theorem~\ref{thm:cheeger} are \stA\ \stM\ for at least two units, Theorem~\ref{thm:laxa}(d) and the other analytic statements are \stA, and the values of $J$ in Table~\ref{tab:C-1} are \stN. Lean checks the logical form of a statement on abstract types and finite graphs and, with Mathlib, on real matrices; whether a physical or phenomenal premise holds is never decided by a compilation.

Table~\ref{tab:evidence} gives, for the principal results of Section~\ref{sec:summary}, the label and the Lean theorems or scripts that check each one; the other statements of Section~\ref{sec:summary} carry their labels and checks in the text. Appendix~\ref{sec:correspondencia} gives the complete correspondence, for every numbered result of the paper.

\begin{table}[htbp]\centering\footnotesize
\begin{tabular}{>{\raggedright\arraybackslash}p{0.06\textwidth}>{\raggedright\arraybackslash}p{0.42\textwidth}>{\raggedright\arraybackslash}p{0.08\textwidth}>{\raggedright\arraybackslash}p{0.32\textwidth}}
\toprule
Item & Main result & Status & Checked by\\
\midrule
(R1) & Entailment criterion (Theorem~\ref{thm:impl}) and co-occurrence dichotomy (Theorem~\ref{thm:dicotomia}) & \stL & \texttt{entails\_Su\_iff}, \texttt{connected\_forces\_constant}\\
(R2) & SEG with distinct values between groups is non-degenerate $S_u$ & \stL & \texttt{segd\_iff}\\
(R3) & Weak thesis: every isomorphism-invariant assignment is a functor (Proposition~\ref{prop:funtor}) & \stA & paper proof\\
(R3) & The bridge square fails on $C_4$ and $P_4$; no translation from positive integration to broadcast & \stL & \texttt{square\_fails}, \texttt{values}, \texttt{no\_translation}\\
(R3) & Strong monoidality into a zero-product monoid, in abstract form & \stL & \texttt{strong\_monoidal\_zero}\\
(R4) & Exact locus of disagreement of the two presence bridges (the band), in Boolean form; its link to $J>0$ is Lemma~\ref{lem:J} & \stL & \texttt{noncomm\_locus}, \texttt{commute\_off\_band}, \texttt{universal\_connected}\\
(R4) & The band is empty for $n\le3$, contains $P_n$ for every $n\ge4$, and has $15$ labeled graphs for $n=4$ & \stL & \texttt{n2\_equiv}, \texttt{n3\_equiv}, \texttt{path\_reach\_from\_zero}, \texttt{path\_no\_universal}, \texttt{n4\_count}\\
(R4) & The band has $472$ labeled graphs for $n=5$ & \stS & \path{pistas_numerico.py}\\
(R5) & For $n\ge2$ units, $J>0$ if and only if the support is connected (Lemma~\ref{lem:J}) & \stA\ \stM & \texttt{QAudit.\allowbreak lemmaJ\_stationary}\\
(R5) & For $n\ge2$ units, the min-cut sandwich (Theorem~\ref{thm:cheeger}) & \stA\ \stM & \texttt{QAudit.\allowbreak sandwich}\\
(R5) & The lower bound of the sandwich is attained in a subfamily & \stS & \path{gauss_conexiones.py}\\
(R6) & Gaussian laboratory: the $156$ exact checks on $F_0$, $F_{\rm mem}$ and $F_0\otimes F_0$ & \stS & \path{v12_familia_escala.py}, \path{v3_modelos_apuntados.py}, \path{v4_historias_ac.py}, \path{v5_rkp_cadenas.py}, \path{gauss_conexiones.py}\\
(R6) & The four-unit networks (Table~\ref{tab:networks}) & \stN & \path{pistas_numerico.py}, \path{J_conectividad.py}\\
(R7) & Collapse of exact covariance and of the order clause on all pairs & \stL & \texttt{exact\_covariance\_gives\_Su}, \texttt{order\_all\_pairs\_collapse}\\
(R7) & Exact frontier of intervention monotonicity (Theorem~\ref{thm:frontier}) & \stL & \texttt{im\_iff\_edges}, \texttt{constant\_on\_scc}, \texttt{nonconstant\_if\_not\_scc}\\
(R7) & For $n\ge2$ units and connected support, strict increase of $J(tA)$ in $t$ and of $J(aI+cW)$ in $c$ (Proposition~\ref{prop:scaling}) & \stA\ \stM & \texttt{QAuditScaling.\allowbreak prop\_scaling\_a}, \texttt{QAuditScaling.\allowbreak prop\_scaling\_b}\\
(R7) & $J$ can decrease when a non-scalar self-coupling is held fixed (Remark~\ref{rem:fixed-diagonal}) & \stN & \path{contraejemplo_diagonal_fija.py}\\
\bottomrule
\end{tabular}
\caption{Evidence for the main results of Section~\ref{sec:summary}: status label (Section~\ref{sec:status}) and the Lean theorems or scripts that check each one. Lean names written without their namespace are theorems of the six files without Mathlib; names written in the namespaces \texttt{QAudit} and \texttt{QAuditScaling} are theorems of the Lean project with Mathlib, with the scope stated in Section~\ref{sec:summary} (\emph{Limitations}); the Lean statements in the two \stL\ rows of (R7) take degrees with integer values. Appendix~\ref{sec:correspondencia} gives the complete correspondence.}\label{tab:evidence}
\end{table}


\part*{Act I.\quad Negative delimitation}
\section{The framework Q and the unity obligation}\label{sec:obligacion}
This section fixes the object of study. Definitions~\ref{def:pcm}--\ref{def:orderB} state the candidate theory Q as far as it is specified: its realizations, its admissible descriptors, its descriptor $\Phi_Q$ and its domain $D_0$, its five components with their sufficiency obligations, the unity degree and the order clause. Definitions~\ref{hyp:LSS}--\ref{def:TL} state the hypotheses that the paper examines. The definitions state what Q asserts; they do not demonstrate it, and nothing in this paper refutes it. The results on an unselected descriptor below are proved here on paper \stA, and those of Section~\ref{sec:intento} are checked by the Lean kernel \stL; Theorem~\ref{thm:conditional-discharge} of Section~\ref{sec:conditional} is a paper proof whose abstract form is also checked by the Lean kernel \stA\ \stL.

\subsection{Statement of the framework}\label{sec:enunciado}
For a function $f$ on a set $D$, its \emph{kernel} is the equivalence relation $\ker f=\{(s,t)\in D\times D: f(s)=f(t)\}$, and its \emph{fibres} are the classes of $\ker f$.

\begin{definition}[Realizations, pointed causal models and isomorphism]\label{def:pcm}
The physical objects of Q and their organizational comparison are the following.
\begin{enumerate}
\item[(a)] A \emph{realization} is a bearer, that is, a physical system with a declared boundary and declared interfaces with its environment, together with its causal dynamics and its realized history over a time window of length $\Delta>0$ that ends at the instant of evaluation.
\item[(b)] A \emph{pointed causal model} is a tuple $M=(N,(X_i)_{i\in N},(p_i)_{i\in N},x)$: a finite set $N$ of units; for each unit $i$ a set $X_i$ of states; for each unit a local update kernel $p_i$, which gives the law of the next state of unit $i$ as a function of the current state of all the units and is read with its interventional semantics (an intervention that sets some units replaces their kernels by constants and leaves the others unchanged); and a current state $x\in\prod_{i\in N}X_i$.
\item[(c)] A \emph{grain} $\gamma$ is a declared rule that assigns to each realization $s$ of a given set a pointed causal model $M_\gamma(s)$: the bearer coarse-grained into units at a declared spatial and temporal resolution, with its current state. The grain is \emph{intrinsic} if $M_\gamma(s)$ depends, up to isomorphism, only on the intrinsic physical history of the bearer over the window.
\item[(d)] An \emph{organizational isomorphism} $M\to M'$ is a bijection $\beta:N\to N'$ together with bijections $r_i:X_i\to X'_{\beta(i)}$ (local recodings) that transport every kernel and the current state exactly: the kernel of unit $\beta(i)$ in $M'$ is the image of the kernel of unit $i$ under the recodings (for finite state sets, $p'_{\beta(i)}(r_i(y)\mid r(z))=p_i(y\mid z)$ for every $y\in X_i$ and $z\in\prod_jX_j$, where $r(z)_{\beta(j)}=r_j(z_j)$; for continuous state sets, the $X_i$ are standard Borel spaces, the $r_i$ are Borel isomorphisms, and $p'_{\beta(i)}(\cdot\mid r(z))=(r_i)_*\,p_i(\cdot\mid z)$ for every $z$), and $x'=r(x)$. Since the kernels are transported, so is their interventional semantics. Two realizations $s,t$ are \emph{isomorphic at grain $\gamma$}, written $s\cong_\gamma t$, if $M_\gamma(s)$ and $M_\gamma(t)$ are isomorphic. Identities, inverses and composites of organizational isomorphisms are organizational isomorphisms, so $\cong_\gamma$ is an equivalence relation; its classes are the \emph{isomorphism classes} at grain $\gamma$.
\item[(e)] Fix finite units and finite state sets. The kernels are then given by a vector $\theta$ of transition probabilities, the \emph{kernel parameters}, in a finite-dimensional real vector space; the kernels of full support form a relatively open convex set, of dimension at least two as soon as some unit has at least two states. The pairs $g=(\beta,(r_i))$ of a permutation of the units and recodings compatible with it form a finite group $G$, which acts on pairs $(\theta,x)$ by permuting the coordinates of $\theta$ and the states; two pointed models with these units and state sets are isomorphic if and only if their pairs $(\theta,x)$ lie in one $G$-orbit.
\end{enumerate}
\end{definition}

\begin{definition}[Descriptors, admissibility and dependence base]\label{def:adm}
Let $D$ be a set of realizations with a grain $\gamma$.
\begin{enumerate}
\item[(a)] A \emph{descriptor} on $D$ is a function $\Phi:D\to V$ into a set of values. It is a descriptor \emph{at grain $\gamma$} if it is computed from the pointed model, $\Phi(s)=\varphi(M_\gamma(s))$ for a function $\varphi$ on pointed causal models.
\item[(b)] A descriptor at grain $\gamma$ is \emph{admissible} (condition \emph{Adm}) if it depends only on the isomorphism class of the pointed model: $s\cong_\gamma t\Rightarrow\Phi(s)=\Phi(t)$, that is, ${\cong_\gamma}\subseteq\ker\Phi$.
\item[(c)] The \emph{dependence base} $B_\Phi$ of a descriptor is the data that its specification reads: the region of the bearer, its interfaces and the dynamical information that the computation needs, including the law of the inputs that the bearer receives over the window if the descriptor is \emph{law-level}, that is, computed from the law of the process rather than from the realized history (\emph{history-level}). A law-level descriptor also needs a declared rule that assigns a law to each realization. Since $\Phi$ is a function of $B_\Phi$, two realizations whose data in $B_\Phi$ coincide, under the same laws of nature, have the same value of $\Phi$.
\end{enumerate}
\end{definition}

\begin{definition}[The candidate theory Q]\label{def:phiQ}
The \emph{candidate theory Q} consists of the following data and claims.
\begin{enumerate}
\item[(a)] \emph{Domain.} A set $D_0$ of realizations, the \emph{original domain}, on which the claims of Q are made. It is meant to contain episodes of conscious subjects; Q does not fix which realizations it contains, in particular whether it contains realizations without experience or artificial systems.
\item[(b)] \emph{Descriptor.} A grain $\gamma$ on $D_0$ and a real-valued admissible descriptor $\Phi_Q:D_0\to\mathbb R$ at grain $\gamma$, the \emph{descriptor of Q}, meant as a measure of the integration of the bearer's causal structure. A specification of $\Phi_Q$ must fix the system boundary, the units, the state, the causal dynamics, the time scale, the partitions, the distance and the selection conventions; Q specifies neither $\Phi_Q$ nor $\gamma$. No functional used in this paper, in particular not the Gaussian functional $J$ of Lemma~\ref{lem:J}, is identified with $\Phi_Q$.
\item[(c)] \emph{Components.} For each $s\in D_0$, a phenomenal tuple $\mathbf e(s)=(c(s),q(s),u(s),\tau(s),\sigma(s))$: the content, the quality, the unity, the temporality and the subjectivity (perspective) of the experience of $s$. The components are features of the total experience of $s$: realizations with the same total experience, under the same declaration of aspects (Definition~\ref{def:unity}), have the same components. They keep their phenomenal meaning; they are not defined from their physical inputs.
\item[(d)] \emph{Reduced inputs.} $d_c=(W_{\rm GW},R_{\rm QD})$, the structured content of a global workspace~\cite{DehaeneChangeux2011} and the redundancy of the records of a pointer variable in the bearer's environment, in the sense of quantum Darwinism~\cite{Zurek2009}; $d_q=(\Phi_Q,\mathcal C)$, with $\mathcal C$ a description of the causal architecture at grain $\gamma$; $d_u=\Phi_Q$; $d_\tau=\mathcal D^\Delta$, a declared collection of finite-resolution dynamical descriptors over the window, such as the finite difference of $\Phi_Q$ over the window divided by $\Delta$ and declared relaxation and decoherence time scales, so that the temporality component is defined on the state extended by $\mathcal D^\Delta$; and $d_\sigma=(A,R)$, attentional and self-model descriptors and homeostatic and interoceptive descriptors. (This $R$ is unrelated to the relation $R$ of Section~\ref{sec:intento}.)
\item[(e)] \emph{Sufficiency obligations.} For $x\in\{c,q,u,\tau,\sigma\}$,
\[
S_x:\qquad d_x(s)=d_x(t)\ \Longrightarrow\ x(s)=x(t)\qquad(s,t\in D_0).
\]
$S_u$ is the obligation~\eqref{eq:Su} of the Introduction. $S_x$ holds if and only if $x=h_x\circ d_x$ for a function $h_x$ on $d_x(D_0)$, and $h_x$ is then unique: set $h_x(d_x(s))=x(s)$, which is well defined exactly when $S_x$ holds. $S_u$ does not identify $u$ with $\Phi_Q$: equal descriptors require equal degrees, unequal descriptors do not require unequal degrees, and $S_u$ says nothing about the nonconstancy, monotonicity, continuity or injectivity of $h_u$. None of the five obligations is established.
\end{enumerate}
Q has further obligations that this paper does not treat, apart from the conditional criterion for descent in Section~\ref{sec:conditional}: descent of the assignment along a physical equivalence $\sim_{\rm phys}$ constructed independently of Q (condition O: $s\sim_{\rm phys}t$ implies $\mathbf e(s)=\mathbf e(t)$), joint admissibility of the five components on one realization, and images of physical processes that respect identities and composition, which rest on premises about the type of the realizations and of the arrows (P$_{\rm branch}$: experience supervenes on the bearer's realized configuration, not on its reduced density operator; P$_{\rm arr}$: the arrows are streams of total experiential states).
We call Q as defined here the \emph{original} Q, and its domain $D_0$, descriptor $\Phi_Q$, degree $u$ and obligations the original ones, to distinguish them from the test bench of Section~\ref{sec:gauss} and from the weakening of Section~\ref{sec:via-futura}.
\end{definition}

\begin{definition}[Episodes, coexperience and the unity degree]\label{def:unity}
The phenomenal side of the unity component is described as follows.
\begin{enumerate}
\item[(a)] An \emph{episode} is a tuple $\mathcal E=(e,\omega,E,\mathsf I)$: an identifier $e$, a delimited time window $\omega$, a finite set $E$ of candidate experiential aspects (unrelated to the noise covariance $E$ of $\mathbf{Open}_G$) and a declared individuation scheme $\mathsf I$ that fixes how aspects are counted. Each realization $s\in D_0$ comes with a declared episode, whose objects carry the index $s$.
\item[(b)] The \emph{coexperience family} of an episode is the set $H_e\subseteq\mathcal P(E)\setminus\{\varnothing\}$ of the sets of aspects that are experienced jointly within $\omega$, that is, subsumed by a single experience, in the subsumptive sense of phenomenal unity, which is distinct from access unity~\cite{BayneChalmers2003}. It is downward closed: if $A\in H_e$ and $\varnothing\ne B\subseteq A$, then $B\in H_e$; this is a modelling commitment. The aspects \emph{present} in the episode are $P_e=\{x\in E:\{x\}\in H_e\}$, and $H_e\subseteq\mathcal P(P_e)$. A realization $s$ is \emph{conscious} if at least one of its declared aspects is present, $P_s\ne\varnothing$, and \emph{non-conscious} otherwise; $D_{\rm cons}$ and $D_\bot$ are the conscious and the non-conscious realizations of $D_0$.
\item[(c)] For a declared bijection $f:E\to E'$ (a \emph{correspondence}) between the aspects of two episodes, $\mathcal E\equiv_f\mathcal E'$ means that $A\in H_e\Leftrightarrow f(A)\in H_{e'}$ for every nonempty $A\subseteq E$. Equality of structures is relative to the correspondence, and labelled equality and equality up to a correspondence are not interchanged.
\item[(d)] The \emph{coexperience structure} of $s\in D_0$ is $\mathfrak U(s)=(E_s,H_s)$, and $\mathfrak U:D_0\to\mathfrak H$, with $\mathfrak H$ the set of such pairs, is the \emph{structural assignment}.
\item[(e)] The \emph{unity degree} $u:D_0\to\mathbb R$, also written $u_{\rm orig}$, is the component $u$ of Definition~\ref{def:phiQ}. Its provisional target is synchronous coexperience: aspects belong to one joint experience within a delimited window. It is distinct from $\mathfrak U$. A degree needs a comparison relation and a statement of which differences it ignores. A \emph{comparison} is a function $\psi:\mathfrak U(D_0)\to\mathbb R$, and the premise
\[
B_\psi:\qquad u=\psi\circ\mathfrak U
\]
states that the degree is a function of the coexperience structure. $B_\psi$ is a hypothesis examined in this paper, not a claim of Q, and Q fixes no $\psi$ and no other comparison relation for its degree.
\end{enumerate}
\end{definition}

\begin{definition}[Order clause; stipulations A9--U and A14--U]\label{def:orderB}
Q constrains the degree by an order clause and two stipulations.
\begin{enumerate}
\item[(a)] The \emph{order clause}, with a strict antecedent, is
\[
\mathrm B:\qquad \Phi_Q(s)<\Phi_Q(t)\ \Longrightarrow\ u(s)\le u(t)\qquad(s,t\in D_0).
\]
$S_u\wedge\mathrm B$ is equivalent to the non-strict implication $\Phi_Q(s)\le\Phi_Q(t)\Rightarrow u(s)\le u(t)$: given the latter, equal descriptors give $u(s)\le u(t)$ and $u(t)\le u(s)$, hence $S_u$, and $\mathrm B$ is a special case; conversely, $S_u$ covers the case of equality and $\mathrm B$ the strict case.
\item[(b)] \emph{A9--U} is the stipulation of Q that the degree satisfies $S_u$ and $\mathrm B$, that is, $u=h_u\circ\Phi_Q$ with $h_u$ nondecreasing on $\Phi_Q(D_0)$. It allows plateaus; it requires neither variation of $u$, nor injectivity, nor strict monotonicity.
\item[(c)] \emph{A14--U} is the stipulation of Q that fixes a rule assigning observables to the degree. Its adoption establishes neither a calibration of a measurement channel for $u$ nor a faithful identification of the degree. An observable that is a function of $\Phi_Q$ is constant on the fibres of $\Phi_Q$ by construction, so agreement with it cannot test $S_u$.
\end{enumerate}
A9--U and A14--U are stipulations of Q; their labels are only names, and no other stipulation of Q is used in this paper. They are not derived from other premises of Q or of this paper, and stipulating them does not justify them: stipulating a restriction on the degree and deriving it from independently justified premises are different achievements. The obligations $S_u$ and $\mathrm B$ that A9--U stipulates remain to be justified on $D_0$.
\end{definition}

\paragraph{Hypotheses examined.} The following hypotheses are not claims of Q. None of them is adopted or justified in this paper; they are examined, and every result that uses one of them says so.

\begin{definition}[Hypothesis LSS$_\Delta$: local supervenience with a finite window]\label{hyp:LSS}
The experience occurring at the end of the window supervenes on the intrinsic physical history of the bearer over the window of length $\Delta$: two realizations whose bearers have duplicate intrinsic histories over their windows, at whatever place, time or environment, have the same total experience.
\end{definition}
If the grain $\gamma$ is intrinsic, realizations with duplicate intrinsic histories are isomorphic at grain $\gamma$, so the equalities of $u$ that LSS$_\Delta$ yields all hold between realizations isomorphic at grain $\gamma$, and, under Adm, between realizations with equal $\Phi_Q$.

\begin{definition}[Hypothesis OI$_\gamma$: organizational invariance]\label{hyp:OI}
If two realizations, one of which is conscious, are isomorphic at a grain $\gamma$ fine enough to capture their functional organization, they have the same total experience.
\end{definition}
This is Chalmers' principle of organizational invariance~\cite{Chalmers1995}, stated for the pointed causal models of Definition~\ref{def:pcm}. Since the components of Q are features of the total experience, OI$_\gamma$ gives $u(s)=u(t)$ whenever $s\cong_\gamma t$ and one of them is conscious (with the declaration of aspects transported along the isomorphism), and it makes each isomorphism class either wholly conscious or wholly non-conscious. Hence, under OI$_\gamma$, $u$ is constant on each isomorphism class at grain $\gamma$ contained in $D_{\rm cons}$, and on every isomorphism class if $u$ takes a single value on $D_\bot$, as it does under B$_\psi^{\rm pres}$ (Definition~\ref{hyp:pres}), where every non-conscious realization has the degree $u_\bot$. When we say below that OI$_\gamma$ makes the degree constant on isomorphism classes, the statement is meant in this sense.

\begin{definition}[Hypothesis Rob: robustness of the domain]\label{hyp:Rob}
There is a conscious realization $s_0\in D_0$ whose pointed model at grain $\gamma$ has finite units and finite state sets, with kernel parameters $\theta_0$ and current state $x_0$, and a convex set $P$ of kernel parameters for these units and state sets (Definition~\ref{def:pcm}(e)) that contains $\theta_0$, is relatively open in its affine hull and has dimension at least two, such that for every $\theta\in P$ the domain $D_0$ contains a conscious realization $s_\theta$ whose pointed model at grain $\gamma$ has these units and state sets, kernel parameters $\theta$ and current state $x_0$.
\end{definition}
The motivation of Rob is that small fluctuations of synaptic parameters do not abolish consciousness.

\begin{definition}[Hypothesis UT: unity thesis]\label{hyp:UT}
For every conscious realization $s\in D_0$, whose window is taken within one specious present, all the aspects present in $s$ are subsumed by one total experience: $P_s\in H_s$, hence $H_s=\mathcal P(P_s)\setminus\{\varnothing\}$.
\end{definition}
This is the unity thesis of Bayne and Chalmers~\cite{BayneChalmers2003} and of Bayne~\cite{Bayne2010}, stated for the episodes of Definition~\ref{def:unity}.

\begin{definition}[Hypothesis B$_\psi^{\rm pres}$: degree of the present structure]\label{hyp:pres}
$B_\psi$ holds, and $\psi(E,H)$ depends only on $H$: the degree is a function of the coexperience among the present aspects, so that the absence of a candidate aspect is not counted as disunity. We write $\psi(H)$ for $\psi(E,H)$ and $u_\bot=\psi(\varnothing)$ for the degree of a non-conscious realization.
\end{definition}

\begin{definition}[Hypothesis N: normalization]\label{hyp:N}
Total coexperience has maximal degree, whatever the number of aspects: there is a value $u_{\max}$ with $u\le u_{\max}$ on $D_0$ such that $u(s)=u_{\max}$ for every conscious $s\in D_0$ with $H_s=\mathcal P(P_s)\setminus\{\varnothing\}$.
\end{definition}

\begin{definition}[Hypothesis: absence anchor]\label{hyp:anchor}
$u(t)<u(s)$ for every non-conscious $t\in D_0$ and every conscious $s\in D_0$.
\end{definition}

\begin{definition}[Threshold law TL]\label{def:TL}
Consciousness is a threshold property of $\Phi_Q$ on $D_0$: $\Phi_Q(t)<\Phi_Q(s)$ for every non-conscious $t\in D_0$ and every conscious $s\in D_0$.
\end{definition}

\paragraph{What Q leaves unspecified, and what is not adopted.} Q does not specify the descriptor $\Phi_Q$ (beyond Adm and real values), the original domain $D_0$, the grain $\gamma$, or the comparison relation of the degree (a $\psi$ as in Definition~\ref{def:unity}, or a substitute), nor the correspondences $f$ and the reading of $\mathcal C$ (unpointed wiring or pointed isomorphism class). These are decisions that Q has not taken, and this paper does not take them: its results hold for every choice that satisfies their stated hypotheses. None of LSS$_\Delta$, OI$_\gamma$, Rob, UT, $B_\psi$, B$_\psi^{\rm pres}$, N or the absence anchor is adopted or justified here; they are hypotheses examined. TL is not assumed either: Theorem~\ref{thm:threshold} shows under which hypotheses it is the content of A9--U.

\subsection{What can be proved about an unselected descriptor}
The results of this subsection are paper proofs \stA. The first says which equalities of the descriptor can be established while $\Phi_Q$ is not selected.

\begin{proposition}[Provable descriptor equalities]\label{prop:provable}
\stA\ Let $D$ be a set of realizations with a grain $\gamma$. For $s,t\in D$, $\Phi(s)=\Phi(t)$ holds for every admissible descriptor $\Phi$ on $D$ if and only if $s\cong_\gamma t$. The same holds when the pointed models have finite units and state sets and the admissible descriptors are required, in addition, to be real-valued and continuous in the kernel parameters.
\end{proposition}
\begin{proof}
If $s\cong_\gamma t$, every admissible descriptor takes equal values at $s$ and $t$. If $s\not\cong_\gamma t$, the indicator of the isomorphism class of $M_\gamma(s)$ is admissible and separates $s$ from $t$. For the continuous case, let $M_1=M_\gamma(s)$, with parameters $\theta_1$ and current state $x_1$. Call two pointed models of the same \emph{shape} if there are bijections between their units and between the state sets of corresponding units; a pointed model of another shape than $M_1$ is not isomorphic to it. For a pointed model $M$ with parameters $\theta$ and current state $x$, put $f(M)=1$ if $M$ is not of the shape of $M_1$, and otherwise
\[
f(M)=\min_{\phi}\bigl(\lVert\phi_*\theta-\theta_1\rVert+\delta(\phi_*x,x_1)\bigr),
\]
where $\phi$ ranges over the finitely many bijections of units and state sets onto those of $M_1$, $\phi_*$ transports parameters and states, $\lVert\cdot\rVert$ is the Euclidean norm and $\delta$ is $0$ on equal states and $1$ otherwise. Composing with an organizational isomorphism permutes the set of the $\phi$, so $f$ is invariant under isomorphism; it is continuous in $\theta$ as a minimum of finitely many continuous functions; and $f(M)=0$ exactly when $\phi_*\theta=\theta_1$ and $\phi_*x=x_1$ for some $\phi$, that is, when $M\cong M_1$. Then $\Phi=f\circ M_\gamma$ is admissible, continuous in the kernel parameters, $\Phi(s)=0$ and $\Phi(t)>0$.
\end{proof}
Hence, as long as $\Phi_Q$ is not selected, the antecedent $\Phi_Q(s)=\Phi_Q(t)$ of $S_u$ can be established, from Adm alone, only for pairs isomorphic at grain $\gamma$; for a non-isomorphic pair it requires selecting and computing $\Phi_Q$, an obligation on the physical side.

\begin{theorem}[Relational coverage]\label{thm:coverage}
\stA\ Let $\Phi:D\to V$ be a function and $\approx$ an equivalence relation on $D$ with ${\approx}\subseteq\ker\Phi$. Every function $u$ on $D$ that is constant on the $\approx$-classes satisfies $S_\Phi:\ \Phi(s)=\Phi(t)\Rightarrow u(s)=u(t)$ if and only if ${\approx}=\ker\Phi$.
\end{theorem}
\begin{proof}
If ${\approx}=\ker\Phi$, then $\Phi(s)=\Phi(t)$ gives $s\approx t$ and hence $u(s)=u(t)$. Conversely, if $\Phi(x)=\Phi(y)$ and $x\not\approx y$, the indicator of the $\approx$-class of $x$ is constant on the $\approx$-classes and takes the value $1$ at $x$ and $0$ at $y$, so it violates $S_\Phi$. Hence, if every such $u$ satisfies $S_\Phi$, then $\ker\Phi\subseteq{\approx}$, and equality follows.
\end{proof}
With ${\approx}={\cong_\gamma}$, which is contained in $\ker\Phi_Q$ by Adm, the theorem applies to the isomorphism groupoid used by LSS$_\Delta$ and OI$_\gamma$: these premises entail $S_u$ for every degree compatible with them if and only if every fibre of $\Phi_Q$ is a single isomorphism class. It is also the case of Theorem~\ref{thm:impl} below in which the relation $R$ is an equivalence relation contained in $\ker\Phi$.

\paragraph{The order clause does not reach the residual \stA.} Theorem~\ref{thm:coverage} remains true when the compatible degrees are also required to satisfy the order clause B (Definition~\ref{def:orderB}(a), with $\Phi$ in place of $\Phi_Q$). Let $\Phi$ be admissible, with values in a strict order, and let a fibre of $\Phi$, of value $c$, contain non-isomorphic realizations $s$ and $t$. The degree equal to $0$ where $\Phi<c$, to $2$ on the class of $t$, to $1$ on the rest of the fibre and to $2$ elsewhere is constant on the isomorphism classes and satisfies B, and $u(s)=1\ne2=u(t)$. Hence admissibility, the constancy of $u$ on isomorphism classes and B together entail $S_u$ for every compatible degree if and only if every fibre of $\Phi_Q$ is a single isomorphism class. Two finite cases show the scope. (A)~Two non-isomorphic realizations with $\Phi_Q=0,0$ and $u=0,1$: B holds vacuously, as in CM1 of Section~\ref{sec:intento}. (B)~Three pairwise non-isomorphic realizations with $\Phi_Q=0,0,1$ and $u=0,1,1$: $\Phi_Q$ is not constant, and B holds on the two pairs whose descriptor increases strictly. In (B) one can take $u=\psi\circ\mathfrak U$ with downward-closed coexperience families on two aspects: the first realization has one aspect present and the other two have both present and coexperienced, so that UT and B$_\psi^{\rm pres}$ hold, with $\psi$ one less than the number of present aspects; N then fails. These models violate $S_u$, and with it A9--U, so they are neither models of Q nor empirical counterexamples to it: they are countermodels to a derivation of $S_u$ from admissibility, the isomorphism invariance of the degree, the order clause and a nonconstant descriptor.

\begin{theorem}[Uncountably many classes in a fibre]\label{thm:continuum}
\stA\ Let $P$ be a relatively open convex subset, of dimension at least two, of a finite-dimensional real affine space, let $\Phi:P\to\mathbb R$ be continuous, and let $c$ be an interior point of the interval $\Phi(P)$. Then the fibre $\Phi^{-1}(c)$ is uncountable. Consequently, if $\sim$ is an equivalence relation on $P$ whose classes have at most $m$ elements for a fixed finite $m$, for example the relation of lying in one orbit of a finite group $G$ acting on a set that contains $P$, with $m=|G|$, then $\Phi^{-1}(c)$ meets uncountably many classes of $\sim$; if $\Phi$ is constant on the classes, it contains uncountably many of them.
\end{theorem}
\begin{proof}
Since $c$ is interior to $\Phi(P)$, there are $\theta_-,\theta_+\in P$ with $\Phi(\theta_-)<c<\Phi(\theta_+)$. Put $w=\theta_+-\theta_-\ne0$ and choose a vector $v$ of the direction space of the affine hull of $P$ that is not parallel to $w$, which exists because the dimension is at least two. The segment $\{\theta_-+\lambda w:\lambda\in[0,1]\}$ lies in $P$ by convexity and is compact, and $P$ is relatively open, so some neighbourhood of the segment in the affine hull lies in $P$. Since the curve $\gamma_\sigma(\lambda)=\theta_-+\lambda w+\sigma\lambda(1-\lambda)v$, $\lambda\in[0,1]$, is within $|\sigma|\,\lVert v\rVert/4$ of the segment, there is $\sigma_0>0$ such that $\gamma_\sigma$ lies in $P$ for $|\sigma|<\sigma_0$. By the intermediate value theorem, $\Phi\circ\gamma_\sigma$ takes the value $c$ at some $\lambda_\sigma\in(0,1)$; put $p_\sigma=\gamma_\sigma(\lambda_\sigma)$. If $p_\sigma=p_{\sigma'}$, then $(\lambda_\sigma-\lambda_{\sigma'})w+(\sigma\lambda_\sigma(1-\lambda_\sigma)-\sigma'\lambda_{\sigma'}(1-\lambda_{\sigma'}))v=0$; since $w$ and $v$ are linearly independent, $\lambda_\sigma=\lambda_{\sigma'}$ and then $\sigma=\sigma'$, because $\lambda_\sigma(1-\lambda_\sigma)>0$. Thus $\sigma\mapsto p_\sigma$ is injective on $(-\sigma_0,\sigma_0)$, and the fibre is uncountable. A countable union of classes with at most $m$ elements each is countable, so the fibre meets uncountably many classes; if $\Phi$ is constant on the classes, each class that meets the fibre is contained in it.
\end{proof}

\begin{proposition}[Form of any extension]\label{prop:extension}
\stA\ Assume Rob, with $s_0$, $\theta_0$, $x_0$, $P$ and the realizations $s_\theta$ as there, and let $G$ be the finite group of Definition~\ref{def:pcm}(e) for the units and state sets of $s_0$. Suppose that $\theta\mapsto\Phi_Q(s_\theta)$ is continuous on $P$ and that $\Phi_Q(s_0)$ is an interior point of its range. Let $F=\{s\in D_0:\Phi_Q(s)=\Phi_Q(s_0)\}$ be the fibre through $s_0$.
\begin{enumerate}
\item[(i)] $F$ contains uncountably many pairwise non-isomorphic conscious realizations.
\item[(ii)] $S_u$ on $F$ follows neither from the equalities of $u$ that OI$_\gamma$ yields nor from those that LSS$_\Delta$ yields at an intrinsic grain: some function on $D_0$ that is constant on the isomorphism classes at grain $\gamma$ violates $S_u$ on $F$.
\item[(iii)] If a family $\Pi$ of premises of the form $u(a)=u(b)$ entails, together with the constancy of $u$ on isomorphism classes, $S_u$ on $F$, in the sense that every function on $D_0$ that satisfies these premises satisfies $S_u$ on $F$, then $\Pi$ contains uncountably many pairs $(a,b)$ of non-isomorphic realizations. In particular no countable family of case-by-case equalities suffices: a uniform principle equating the degree across non-isomorphic organizations is required.
\end{enumerate}
\end{proposition}
\begin{proof}
(i) By Rob the realizations $s_\theta$, $\theta\in P$, are conscious and lie in $D_0$. By Definition~\ref{def:pcm}(e), $s_\theta\cong_\gamma s_{\theta'}$ only if $(\theta',x_0)=g\cdot(\theta,x_0)$ for some $g\in G$, so the relation $s_\theta\cong_\gamma s_{\theta'}$ on $P$ has classes with at most $|G|$ elements. Theorem~\ref{thm:continuum}, applied to $\theta\mapsto\Phi_Q(s_\theta)$ and $c=\Phi_Q(s_0)$, gives uncountably many classes in the fibre.
(ii) Let $u$ be the indicator of the isomorphism class of $s_0$. It is constant on the isomorphism classes, hence satisfies every equality that OI$_\gamma$ yields and, at an intrinsic grain, every equality that LSS$_\Delta$ yields (both hold only between isomorphic realizations). By (i) there is $t\in F$ with $t\not\cong_\gamma s_0$, and $u(t)=0\ne1=u(s_0)$.
(iii) Suppose that only countably many pairs of $\Pi$ are non-isomorphic. Their realizations lie in countably many isomorphism classes, so by (i) there are two non-isomorphic realizations $s,t\in F$ whose classes contain no realization of such a pair. The indicator $u$ of the class of $s$ is constant on the isomorphism classes; it satisfies every premise of $\Pi$, since a pair of isomorphic realizations lies in one class and a non-isomorphic pair lies outside the class of $s$; and $u(s)=1\ne0=u(t)$ with $\Phi_Q(s)=\Phi_Q(t)$. So the premises do not entail $S_u$ on $F$, a contradiction.
\end{proof}

\begin{theorem}[Threshold reduction]\label{thm:threshold}
\stA\ Let $D=D_{\rm cons}\sqcup D_\bot$ be a set of realizations divided into conscious and non-conscious ones, let $\Phi:D\to V$ take values in a strict total order $(V,<)$, and write $S$ and $\mathrm B$ for $S_u$ and the order clause with $\Phi$ in place of $\Phi_Q$.
\begin{enumerate}
\item[(i)] Let $u=u_{\max}$ on $D_{\rm cons}$ and $u=u_\bot$ on $D_\bot$, with $u_\bot<u_{\max}$. Then $S$ holds if and only if no fibre of $\Phi$ meets both $D_{\rm cons}$ and $D_\bot$.
\item[(ii)] For $u$ as in (i), $\mathrm B$ holds if and only if $\Phi(t)\le\Phi(s)$ for every $s\in D_{\rm cons}$ and $t\in D_\bot$.
\item[(iii)] For $u$ as in (i), $S\wedge\mathrm B$ holds if and only if $\Phi(t)<\Phi(s)$ for every $s\in D_{\rm cons}$ and $t\in D_\bot$, that is, if and only if consciousness is a threshold property of $\Phi$ on $D$ (for $\Phi=\Phi_Q$ and $D=D_0$, the threshold law TL).
\item[(iv)] For $u$ as in (i), $\mathrm B$ alone allows at most one fibre of $\Phi$ that meets both $D_{\rm cons}$ and $D_\bot$, and the bound is attained.
\item[(v)] For any $u:D\to\mathbb R$ that satisfies the absence anchor, $S\wedge\mathrm B$ implies $\Phi(t)<\Phi(s)$ for every $s\in D_{\rm cons}$ and $t\in D_\bot$. The converse fails in general.
\end{enumerate}
\end{theorem}
\begin{proof}
(i) $u$ takes two distinct values, so $S$ fails exactly when some $s\in D_{\rm cons}$ and $t\in D_\bot$ have $\Phi(s)=\Phi(t)$. (ii) Since $u_\bot<u_{\max}$, $u(s)>u(t)$ holds exactly when $s\in D_{\rm cons}$ and $t\in D_\bot$, so $\mathrm B$ fails exactly when some such pair has $\Phi(s)<\Phi(t)$; in a strict total order the negation of $\Phi(s)<\Phi(t)$ is $\Phi(t)\le\Phi(s)$. (iii) By (i) and (ii), $S\wedge\mathrm B$ holds exactly when $\Phi(t)\le\Phi(s)$ and $\Phi(t)\ne\Phi(s)$ for every such pair. (iv) If two such fibres had values $a<a'$, the first would contain some $s\in D_{\rm cons}$ and the second some $t\in D_\bot$, with $\Phi(s)<\Phi(t)$, against (ii). With $D=\{s,t\}$, $s$ conscious, $t$ non-conscious and $\Phi(s)=\Phi(t)$, $\mathrm B$ holds and there is exactly one such fibre. (v) If $\Phi(t)=\Phi(s)$, $S$ gives $u(t)=u(s)$; if $\Phi(s)<\Phi(t)$, $\mathrm B$ gives $u(s)\le u(t)$; both contradict the anchor, and totality leaves $\Phi(t)<\Phi(s)$. For the converse, two conscious realizations with equal $\Phi$ and different $u$, and a non-conscious one with smaller $\Phi$ and smaller $u$, satisfy the anchor and the separation and violate $S$.
\end{proof}

\paragraph{The degree under UT, B$_\psi^{\rm pres}$ and N \stA.} Under UT, every conscious $s\in D_0$ has $H_s=\mathcal P(P_s)\setminus\{\varnothing\}$, so N gives $u(s)=u_{\max}$; under B$_\psi^{\rm pres}$, every non-conscious $s$ has $H_s=\varnothing$ and $u(s)=\psi(\varnothing)=u_\bot$, and N gives $u_\bot\le u_{\max}$. So under these three hypotheses the degree on $D_0$ has the form of Theorem~\ref{thm:threshold}(i), except for the strict inequality $u_\bot<u_{\max}$, which they do not supply; when $D_0$ contains both conscious and non-conscious realizations, that inequality is the absence anchor. If $u_\bot=u_{\max}$, the degree is constant on $D_0$, and $S_u$ and $\mathrm B$ hold for every descriptor. Hence, under UT, B$_\psi^{\rm pres}$, N and the anchor, the non-degenerate content of A9--U on $D_0$ is exactly TL, by Theorem~\ref{thm:threshold}(iii); this reading needs a two-valued degree, that is, N together with the anchor. The refuting direction does not need N: by Theorem~\ref{thm:threshold}(v), $S_u$, $\mathrm B$ and the anchor already imply TL. Applying either statement to particular realizations would also require an independently justified attribution of absence of experience to the non-conscious ones, which is not available.

\subsection{Dependency graph}\label{sec:grafo}
\begin{figure}[htbp]
\centering
{\small\resizebox{\textwidth}{!}{$\begin{array}{ccccccccc}
\text{\bf Definitions} && \text{\bf Physical results} && \text{\bf Psychophysical premises} && \text{\bf Sufficiencies} && \text{\bf Objects/processes}\\[2pt]
\begin{array}{c}D_0,\ \Phi_Q,\ \gamma,\ \mathcal C,\\ f,\ \psi\ \text{(comparison of }u)\end{array}
&\to&
\begin{array}{c}\text{Adm, base }B_\Phi,\\ \text{continuum},\\ \text{provable}\\ \text{equalities}\end{array}
&\to&
\begin{array}{c}\text{LSS}_\Delta,\ \text{OI}_\gamma,\ \text{Rob},\\ \text{UT},\ B_\psi,\ B_\psi^{\rm pres},\ \text{N},\\ \text{anchor},\ \text{TL}\end{array}
&\to&
\begin{array}{c}S_c,S_q,\mathbf{S_u}\ (\text{SEG}^{\ast}),\\ S_\tau,S_\sigma,\ \mathrm B\end{array}
&\to&
\begin{array}{c}\sim_{\rm phys},\ \mathrm O,\\ \text{P}_{\rm arr},\ \text{arrows},\\ \text{descent}\end{array}
\end{array}$}}
\caption{Dependency graph: definitions, physical results, psychophysical premises, sufficiencies, and objects and processes. The obligation $S_u$ is in bold. ${}^{\ast}$SEG is equivalent only to non-degenerate $S_u$, and only under $B_\psi$ (Definition~\ref{def:SEG} in Section~\ref{sec:mejor}).}
\label{fig:graph}
\end{figure}
The items in the first column of Figure~\ref{fig:graph} are left unspecified by Q (Definitions~\ref{def:phiQ} and~\ref{def:unity}), so no proposition in the later columns that mentions them can be evaluated for $D_0$. The remaining dependencies are implications: factorization and functorial assembly are conditional results (Definition~\ref{def:phiQ}(e) and Section~\ref{sec:intento}). The last column lists the further obligations of Q that this paper does not treat, except for the conditional criterion for descent in Section~\ref{sec:conditional} (Definition~\ref{def:phiQ}). SEG is not a premise of the third column: it is non-degenerate $S_u$, under $B_\psi$, rewritten in testable form (Section~\ref{sec:mejor}). That is why it appears next to $S_u$, and identifying it does not close $S_u$.

\subsection{The exact gap}\label{sec:laguna}
\paragraph{(L0) A specification gap precedes the justification gap.} $\Phi_Q$, $D_0$, $\gamma$ and the comparison relation of the degree (a $\psi$, or a substitute) are not specified by Q (Definitions~\ref{def:phiQ} and~\ref{def:unity}). Hence $S_u$ is a schema, not a determinate proposition. Moreover, by Proposition~\ref{prop:provable}, as long as $\Phi_Q$ is not selected, the antecedent $\Phi_Q(s)=\Phi_Q(t)$ can be established only for pairs isomorphic at grain $\gamma$. For any non-isomorphic pair, $S_u$ cannot even be put to the test.

\paragraph{(L1) The cross-class residual.} Under Adm and OI$_\gamma$ (not adopted), the degree is constant on each isomorphism class at grain $\gamma$ (in the sense stated after Definition~\ref{hyp:OI}). Then
\begin{equation}
S_u\iff u_{\rm orig}(s)=u_{\rm orig}(t)\ \text{for every non-isomorphic pair with } \Phi_Q(s)=\Phi_Q(t)
\end{equation}
(Theorem~\ref{thm:residuo}). With Rob and $\Phi_Q$ continuous in the kernel parameters, every fibre of $\Phi_Q$ whose value is interior to the range of $\Phi_Q$ on the neighbourhood of Rob contains uncountably many non-isomorphic classes (Theorem~\ref{thm:continuum} and Proposition~\ref{prop:extension}(i)). No isomorphism-type principle covers that residual (Theorem~\ref{thm:coverage}): a uniform principle equating the degree across non-isomorphic organizations is required (Proposition~\ref{prop:extension}(iii)).

\paragraph{Why the available arguments do not close it.}
\begin{itemize}
\item Factorization ($u=h\circ\Phi_Q$) is equivalent to $S_u$ (Definition~\ref{def:phiQ}(e)), so it presupposes $S_u$.
\item The route through an equivalence ${\approx}$ with $\ker\Phi_Q\subseteq{\approx}\subseteq\ker u$ and the route through transformations that cover the fibres and preserve $u$ are conditional. The least such equivalence is $\ker\Phi_Q$ itself, and preserving $u$ on it is $S_u$. A relation $R\subseteq\ker\Phi_Q$ whose generated equivalence contains $\ker\Phi_Q$, such as a family of $\Phi_Q$-preserving transformations that connects each fibre, can have far fewer pairs, but the premises ``$u$ agrees on $R$-related pairs'' are then jointly equivalent to $S_u$ (Theorem~\ref{thm:impl}). Such a route would discharge $S_u$ only if each of its premises had a defense of its own: equivalence is not circularity, but no such defense is available for the original descriptor.
\item Supervenience (LSS$_\Delta$ at an intrinsic grain) and invariance under a change of measurement units (a recoding of states, hence an organizational isomorphism) equate the degree only within isomorphism classes at grain $\gamma$; so do stability, bisimulation and equality of uncalibrated observations when they are read as exact correspondences between unit-level models with their interventions. None of them goes beyond the isomorphism classes, so, by Theorem~\ref{thm:coverage}, none reaches the residual (L1).
\item An observational route would use an observable $L$ with $L_s=K(u(s))$ for an injective $K$, together with the premise $J_0(s)=J_0(t)\Rightarrow L_s=L_t$, where $J_0(s)=J(\ell(s))$ is the Gaussian functional of Lemma~\ref{lem:J} evaluated at a linear-Gaussian law $\ell(s)$ assigned to $s$ by a declared rule, a candidate descriptor that is not $\Phi_Q$. Given the first two premises, the third is equivalent to the sufficiency $S_{J_0}$: $J_0(s)=J_0(t)\Rightarrow u(s)=u(t)$, since, $K$ being injective, $L_s=L_t$ holds exactly when $u(s)=u(t)$. Justifying it requires measuring $L$ across non-isomorphic realizations with equal $J_0$, that is, testing $S_{J_0}$; that does not prove it.
\end{itemize}

\subsection{Substantive attempt: what factorization premises can deliver}\label{sec:intento}
All the independent defenses of a unity bridge that this paper examines have the form ``the degree depends only on $\kappa$'', for some feature $\kappa$ with a defense of its own: duplication (LSS$_\Delta$), pointed organization (OI$_\gamma$), coexperience structure ($B_\psi$), presence class (UT with B$_\psi^{\rm pres}$) or membership in a complex (the exclusion postulate of IIT~\cite{Oizumi2014}). Let $R$ be the relation ``having the same $\kappa$'' for the defended features, and let $\langle R\rangle$ be the equivalence it generates. We call a feature $\kappa$ \emph{aligned} with $\Phi_Q$ if equal $\kappa$ implies equal $\Phi_Q$, that is, if each class of $\kappa$ lies inside a fibre of $\Phi_Q$, and \emph{transversal} to $\Phi_Q$ otherwise, that is, if one and the same value of $\kappa$ occurs at different values of $\Phi_Q$; the pair $(\Phi_Q,\kappa)$ has \emph{product image} if every realized value of $\Phi_Q$ occurs together with every realized value of $\kappa$. All of the following statements are formalized and compiled in \texttt{IntentoSu.lean} \stL\ (Appendix~\ref{sec:archivos}); the short proofs given here follow the Lean proofs. In Lean, Theorems~\ref{thm:impl} and~\ref{thm:holgura} quantify over degrees with values in an arbitrary type, and their proofs use $\mathrm{Prop}$-valued class indicators; composing these with the map that sends true to $1$ and false to $0$ gives the statements for real-valued degrees, a step that is not compiled.

\begin{theorem}[Entailment criterion; \texttt{entails\_Su\_iff}]\label{thm:impl}
The premises ``$u$ agrees on $R$-related pairs'' entail $S_u$ for every compatible $u$ if and only if $\ker\Phi_Q\subseteq\langle R\rangle$.
\end{theorem}
\begin{proof}
($\Leftarrow$) Every $u$ that respects $R$ is constant on the classes of $\langle R\rangle$, by induction on the chain. ($\Rightarrow$) The indicator function $x\mapsto[\langle R\rangle(t,x)]$ respects $R$. Applying $S_u$ to it on a pair with equal $\Phi_Q$ yields $\langle R\rangle(s,t)$.
\end{proof}
The criterion is abstract: it says which relations would suffice, not that one is available. For the original descriptor, the open task is to find a relation $R$ whose pairs each have a defense of their own and whose generated equivalence contains $\ker\Phi_Q$, that is, joins every two realizations of each fibre of $\Phi_Q$, non-isomorphic ones included. A compiled criterion does not show that such a relation exists physically, and for none of the relations examined in this paper that has a defense of its own is the inclusion established (Section~\ref{sec:mejor}); for $\ker\Phi_Q$ itself it is trivial, and the premise is then $S_u$.

\begin{theorem}[Slack and saturation; \texttt{nonconstant\_iff}, \texttt{saturated}]\label{thm:holgura}
These premises admit a nonconstant $u$ if and only if $\langle R\rangle$ is not total. If they entail $S_u$, every class of $\langle R\rangle$ is a union of fibres of $\Phi_Q$.
\end{theorem}
\begin{proof}
If $\langle R\rangle$ is total, every $u$ that respects $R$ is constant, by the induction of Theorem~\ref{thm:impl}; if $\langle R\rangle(a,b)$ fails, the indicator of the class of $a$ respects $R$ and separates $a$ from $b$. If the premises entail $S_u$, Theorem~\ref{thm:impl} gives $\ker\Phi_Q\subseteq\langle R\rangle$, so $\Phi_Q(s')=\Phi_Q(s)$ and $\langle R\rangle(s,t)$ give $\langle R\rangle(s',t)$.
\end{proof}

\begin{theorem}[Co-occurrence dichotomy; \texttt{connected\_forces\_constant}, \texttt{constant\_suff}]\label{thm:dicotomia}
Let $u=f\circ\kappa$. If $S_u$ holds and the relation ``equal $\kappa$ or equal $\Phi_Q$'' connects all of $D_0$, then $u$ is constant on $D_0$. Conversely, a constant $u$ satisfies $S_u$.
\end{theorem}
\begin{proof}
Equal $\kappa$ gives equal $f\circ\kappa$, and equal $\Phi_Q$ gives equal $u$ by $S_u$; so $u$ is constant along every chain of the relation, and by hypothesis every two realizations of $D_0$ are joined by such a chain. The converse is immediate.
\end{proof}
This is an elementary fact about generated equivalences: a function that is constant along two kernels is constant along every chain of the equivalence they generate. A product joint image of $(\Phi_Q,\kappa)$ is a sufficient condition for connectivity, since any $a,b$ are then joined through a realization $c$ with $\Phi_Q(c)=\Phi_Q(a)$ and $\kappa(c)=\kappa(b)$; connectivity is weaker.

\begin{theorem}[Residual under OI; \texttt{oi\_residual}]\label{thm:residuo}
If isomorphic pairs have equal $u$, then $S_u$ is equivalent to the equality of $u$ on the non-isomorphic pairs with equal $\Phi_Q$.
\end{theorem}
\begin{proof}
$S_u$ gives the equality on all pairs with equal $\Phi_Q$. Conversely, a pair with equal $\Phi_Q$ is either isomorphic, and then has equal $u$ by hypothesis, or not, and then has equal $u$ by assumption.
\end{proof}

\begin{theorem}[Robust form; \texttt{robust\_iff\_relational}]\label{thm:robusta}
Suppose that there exists an injective representation of the coexperience structures, which is the case, for example, if they are countable. Then $S_u$ holds for every $u=\psi\circ\mathfrak U$ if and only if $\Phi_Q$ determines the coexperience structure, that is, the relational sufficiency $\ker\Phi_Q\subseteq\ker\mathfrak U$.
\end{theorem}
\begin{proof}
If $\Phi_Q$ determines $\mathfrak U$, equal $\Phi_Q$ gives equal $\mathfrak U$ and hence equal $\psi\circ\mathfrak U$ for every $\psi$. Conversely, $S_u$ for $u=\psi_0\circ\mathfrak U$ with $\psi_0$ injective gives $\psi_0(\mathfrak U(s))=\psi_0(\mathfrak U(t))$, hence $\mathfrak U(s)=\mathfrak U(t)$, whenever $\Phi_Q(s)=\Phi_Q(t)$.
\end{proof}
Since Q fixes no $\psi$ (Definition~\ref{def:unity}), a proof of $S_u$ that uses no property of a specific $\psi$ is equivalent to proving that $\Phi_Q$ determines the coexperience structure. For one specified $\psi$ the condition is weaker: $S_u$ for $u=\psi\circ\mathfrak U$ is $\ker\Phi_Q\subseteq\ker(\psi\circ\mathfrak U)$, and it can hold, with $u$ not constant, although $\Phi_Q$ does not determine $\mathfrak U$. An example is given in the proof of Theorem~\ref{thm:conditional-discharge}(ii). The inclusion $\ker\Phi_Q\subseteq\ker\mathfrak U$ is necessary for a proof that holds for every $\psi$, and for a proof for one $\psi$ when that $\psi$ takes different values on any two distinct structures that occur in one fibre of $\Phi_Q$, as an injective $\psi$ does (CM4 below). \stA

\paragraph{Countermodels (checked by the Lean kernel).}
\begin{itemize}
\item \textbf{CM1.} OI and Adm hold and $S_u$ fails: two classes with equal $\Phi_Q$ and different $u$.
\item \textbf{CM2.} Transversal $\kappa$ with product image: $u=f\circ\kappa$ is nonconstant and $S_u$ fails.
\item \textbf{CM3.} Segregation: if $\kappa$ is a function of $\Phi_Q$, then a nonconstant $u=\kappa$ satisfies $S_u$, with disconnected co-occurrence. This shows that the connectivity hypothesis of Theorem~\ref{thm:dicotomia} is necessary.
\item \textbf{CM4.} Equal $\Phi_Q$ and different structures: with $\psi$ injective, $S_u$ fails.
\end{itemize}

\paragraph{Application to the defended features.}
\begin{itemize}
\item \emph{LSS$_\Delta$ and OI$_\gamma$ are aligned with $\Phi_Q$.} By Adm, their classes lie inside the fibres (for LSS$_\Delta$, at an intrinsic grain), so they give equality only within each class. The residual (L1) remains intact (Theorem~\ref{thm:residuo}, CM1).
\item \emph{$B_\psi$ is transversal.} One and the same coexperience structure can occur with different values of $\Phi_Q$. Under OI, isomorphic realizations have the same total experience and hence the same structure, so $\langle\cong\cup\ker\mathfrak U\rangle=\ker\mathfrak U$, and by Theorem~\ref{thm:impl} the entailment of $S_u$ for every compatible degree, that is, for every $\psi$, requires $\Phi_Q$ to determine $\mathfrak U$ (Theorem~\ref{thm:robusta}); for one specified $\psi$ it requires only $\ker\Phi_Q\subseteq\ker(\psi\circ\mathfrak U)$, which is $S_u$ for that degree (the remark after Theorem~\ref{thm:robusta}). If the structures co-occur in a connected way across the levels of $\Phi_Q$, $S_u$ forces $u$ to be constant (Theorem~\ref{thm:dicotomia}).
\item \emph{With UT and B$_\psi^{\rm pres}$}, the degree of a conscious realization is $\psi(\mathcal P(P_s)\setminus\{\varnothing\})$, a function of the set $P_s$ of present aspects. If the degree is moreover invariant under the admitted correspondences of aspects ($\psi$ takes equal values on structures related by $\equiv_f$ for an admitted $f$, the admitted correspondences forming a groupoid), it is a function of the presence class, the class of $P_s$ under the admitted correspondences: a correspondence $f$ with $f(P_s)=P_t$ gives $\mathcal E_s\equiv_f\mathcal E_t$, since $\varnothing\ne A\subseteq P_s$ if and only if $\varnothing\ne f(A)\subseteq P_t$. Then $\kappa$ is the presence class; without that invariance, $\kappa$ is the set of present aspects, and what follows holds likewise by Theorem~\ref{thm:dicotomia}. Under connected co-occurrence, $S_u$ is equivalent to the constancy of $u$ on $D_0$; in a domain of conscious episodes, that is the content of N restricted to $D_0$. A9--U admits plateaus and does not require variation (Definition~\ref{def:orderB}), and not adopting N is not denying it. Under these hypotheses, a defender of A9--U has to accept a constant $u$ on $D_0$ or, if a non-degenerate $S_u$ is sought, reject UT, reject B$_\psi^{\rm pres}$ or reject connectivity.
\end{itemize}

\subsection{Candidate principles and primary sources}\label{sec:principios}
An AI agent (see \hyperref[sec:ai]{the declaration on the use of generative AI}) checked the available sources word for word against local copies: Chalmers (1995)~\cite{Chalmers1995}, Bayne (2010)~\cite{Bayne2010}, Schechter (2018)~\cite{Schechter2018}, three technical papers on integrated information~\cite{Kitazono2018,Oizumi2016PNAS,TokerSommer2019} and Mendola (2025)~\cite{Mendola2025}. The page numbers were taken from the running headers of plain-text copies and have not been checked against the printed editions. IIT~3.0~\cite{Oizumi2014}, whose postulates three rows of Table~\ref{tab:B-1} use, was not among these copies: those rows describe it without quotation, and the description has not been checked word for word against the original. Our answer, on the basis of these sources only, to the decisive question of Section~\ref{sec:mejor} (whether any available material rules out, with an independent foundation, that two original realizations with the same $\Phi_Q$ have different $u_{\rm orig}$) is \textbf{no}.

Table~\ref{tab:B-1} lists the candidate principles. The hypotheses OI$_\gamma$ and UT are stated in Definitions~\ref{hyp:OI} and~\ref{hyp:UT}, and Rob, B$_\psi^{\rm pres}$, N and the absence anchor in Definitions~\ref{hyp:Rob}, \ref{hyp:pres}, \ref{hyp:N} and~\ref{hyp:anchor}; the degree $u$ and the coexperience structure $\mathfrak U$ are those of Definition~\ref{def:unity}, and $B_\psi$ is the premise $u=\psi\circ\mathfrak U$.

\begin{table}[htbp]\centering\footnotesize
\begin{tabular}{>{\raggedright\arraybackslash}p{0.19\textwidth}>{\raggedright\arraybackslash}p{0.15\textwidth}>{\raggedright\arraybackslash}p{0.57\textwidth}}
\toprule
Principle & Scope & Foundation and limit\\
\midrule
OI$_\gamma$~\cite{Chalmers1995} & classes only & ``any system that has the same functional organization at a fine enough grain will have qualitatively identical conscious experiences'' (Chalmers, §1). It is independent, but it equates only isomorphic realizations, so it does not reach the cross-class residual (Theorem~\ref{thm:coverage}), which under Rob, with $\Phi_Q$ continuous, is a continuum in each interior fibre (Theorem~\ref{thm:continuum}, Proposition~\ref{prop:extension}).\\
Extension by fading & fibres, but (a) & A fading-type argument in the style of~\cite{Chalmers1995} infers that experience stays the same along a path of realizations from the fact that the subject's judgments and reports stay the same along it. Inside an isomorphism class the reports stay the same because the organization does. To extend the argument across the non-isomorphic realizations of a fibre of $\Phi_Q$, the reports about unity must stay the same across the fibre, that is, equal $\Phi_Q$ must give equal reports, which is $S_u$ at the level of reports. The extension therefore presupposes the bridge at the level of reports.\\
UT~\cite{Bayne2010,BayneChalmers2003} & fibres, but (b) & ``the simultaneous conscious states that S has at t will be subsumed by a single conscious state'' (Bayne, p.~16). Its domain is human beings (p.~17). With B$_\psi^{\rm pres}$ and N it gives $S_u$ in degenerate form: UT makes the coexperience structure total on the present aspects, B$_\psi^{\rm pres}$ makes $u$ a function of that structure, and N gives it the maximal value, so $u$ is constant on conscious realizations and $S_u$ holds there for every descriptor. Without N it is silent across different sets of present aspects: $u$ is then a function of the set of present aspects, and UT does not constrain that function.\\
IIT integration and exclusion~\cite{Oizumi2014} & (b) or (c) & Applied per complex, integration and exclusion make the experience of a complex a single one, so all its aspects are coexperienced, as UT says of a subject; with B$_\psi^{\rm pres}$ and N, $u$ then collapses as under UT. Applied per bearer, which may host several complexes, they make $u$ depend on how many complexes the bearer hosts, a count rather than a degree of coexperience, which changes the meaning of $u$. None of the sources checked argues that the degree of coexperience within a complex varies with $\Phi$.\\
IIT identity~\cite{Oizumi2014} & classes only & An experience is identical with the conceptual (cause--effect) structure that its complex specifies. The thesis equates $u$ only between realizations whose structures are isomorphic, and it entails $S_u$ for every compatible degree if and only if $\Phi_Q$ determines that structure up to isomorphism (Theorem~\ref{thm:impl}). Used to derive $S_u$, it thus needs a further premise that, with the thesis, already implies $S_u$: that $\Phi_Q$ determines the structure, or that $u$ depends on the structure only through $\Phi_Q$, which without a defense of its own is $S_u$ under another name.\\
IIT quantity; GNWT access & (c) & In IIT~3.0~\cite{Oizumi2014}, $\Phi^{\max}$ measures the quantity, or level, of consciousness of a complex; GNWT~\cite{DehaeneChangeux2011} concerns the global availability of information across brain systems. Coexperience is a relation among aspects within one experience (Definition~\ref{def:unity}). A level and an access relation are different targets, and taking either as $u$ changes the meaning of $u$.\\
Two unities (split brain)~\cite{Schechter2018}; Mendola~\cite{Mendola2025} & countermodels & ``not disunity or even simply partial unity but two unities'' (Schechter, p.~20). ``different sorts of phenomenal consciousness that cannot be organized along a single scale from least to greatest'' (Mendola, p.~82), said of forms and degrees of phenomenal consciousness, which Mendola carries over to unity in split-brain cases: ``a lesser degree or attenuated form of phenomenal consciousness and so unity that bridges the whole'' (p.~113). They count against scalar sufficiency, not for it.\\
\bottomrule
\end{tabular}
\caption{Candidate principles, with their scope, foundation and limit.}\label{tab:B-1}
\end{table}

In Table~\ref{tab:B-1}, (a) presupposes the bridge; (b) collapses $u$ into a constant or into a presence indicator; (c) changes the meaning of $u$. Together with the principles that reach only isomorphism classes, these are the ways out that Theorems~\ref{thm:impl}--\ref{thm:robusta} leave for the principles examined; changing the domain, as selection in the fibres does, is a further one (Section~\ref{sec:l3}).

\subsection{Best result and the premises that remain free}\label{sec:mejor}
\paragraph{Answer to the decisive question.} No available material offers an independent foundation for ruling out that two original realizations with the same $\Phi_Q$ have different $u_{\rm orig}$ when $u_{\rm orig}$ is not constant. By the theorems of Section~\ref{sec:intento}, the reason is structural:
\begin{enumerate}
\item The principles with a defense of their own that are aligned with $\Phi_Q$ (LSS$_\Delta$ and OI$_\gamma$, Definitions~\ref{hyp:LSS} and~\ref{hyp:OI}) equate only within isomorphism classes. The cross-class residual of each fibre remains intact (Theorems~\ref{thm:coverage} and~\ref{thm:residuo}), and under Rob, with $\Phi_Q$ continuous, it is a continuum (Theorem~\ref{thm:continuum}, Proposition~\ref{prop:extension}).
\item The transversal principles (coexperience structure, presence class, membership in a complex) imply $S_u$ for every degree compatible with them, that is, whatever function of the structure the degree is, only if $\Phi_Q$ determines that structure (Theorems~\ref{thm:impl} and~\ref{thm:robusta}). For one specified function $\psi$, $S_u$ requires only that structures that share a value of $\Phi_Q$ have equal values of $\psi$, which is condition~(ii) of SEG below with the level sets of $\psi$ as groups; it requires $\Phi_Q$ to separate two structures only when $\psi$ takes different values on them, and to determine the structures when $\psi$ takes different values on any two distinct structures, as an injective $\psi$ does. No $\psi$ has been specified (Definition~\ref{def:unity}), and for no nonconstant $\psi$ is that condition justified. If their co-occurrence with the levels of $\Phi_Q$ is connected, $S_u$ forces constancy (Theorem~\ref{thm:dicotomia}): the degree becomes the presence indicator or the normalization N, which are not adopted (Definitions~\ref{hyp:N} and~\ref{hyp:anchor}).
\end{enumerate}

\paragraph{SEG: non-degenerate $S_u$ in testable form.} Suppose that the degree is a function of the coexperience structure, $u=\psi\circ\mathfrak U$ ($B_\psi$, Definition~\ref{def:unity}). Then non-degenerate $S_u$, that is, $S_u$ with $u$ not constant on $D_0$, is equivalent to the following condition.

\begin{definition}[SEG]\label{def:SEG}
Assume $B_\psi$. The \emph{banded segregation} SEG holds on $D_0$ if the coexperience structures present in $D_0$, that is, the values $\mathfrak U(s)$ with $s\in D_0$, split into at least two groups such that:
\begin{enumerate}
\item[(i)] $\psi$ is constant on each group and takes different values on distinct groups;
\item[(ii)] no value of $\Phi_Q$ hosts structures from two distinct groups: if $s,t\in D_0$ and $\Phi_Q(s)=\Phi_Q(t)$, then $\mathfrak U(s)$ and $\mathfrak U(t)$ lie in the same group.
\end{enumerate}
\end{definition}

\begin{proposition}[\texttt{seg\_implies\_Su}, \texttt{segd\_iff}]\label{prop:seg}
Assume $B_\psi$. \textup{(a)} A split of the structures present in $D_0$ into at least two groups, with $\psi$ constant on each group and with condition~(ii) of Definition~\ref{def:SEG}, implies $S_u$. \textup{(b)} SEG holds if and only if $S_u$ holds and $u=\psi\circ\mathfrak U$ is not constant on $D_0$. \stL
\end{proposition}
\begin{proof}
(a) If $\Phi_Q(s)=\Phi_Q(t)$, condition~(ii) puts $\mathfrak U(s)$ and $\mathfrak U(t)$ in one group, on which $\psi$ is constant, so $u(s)=u(t)$. (b) If SEG holds, part~(a) gives $S_u$, and two realizations whose structures lie in distinct groups have different values of $u$, so $u$ is not constant. Conversely, if $S_u$ holds and $u$ is not constant on $D_0$, group the present structures by their value of $\psi$: condition~(i) holds by construction, there are at least two groups because $u$ is not constant, and condition~(ii) is $S_u$.
\end{proof}
In \texttt{FiltracionCociente.lean} (Appendix~\ref{sec:archivos}), the predicate \texttt{SEG} is the split of part~(a), without the requirement that $\psi$ differ between groups, and \texttt{SEGd} is Definition~\ref{def:SEG}.

SEG is not a new independent premise: it is non-degenerate $S_u$ written in the form in which it could be tested. Being equivalent is not being circular: a defense of SEG of its own, theoretical or empirical, would be a defense of non-degenerate $S_u$, and what is missing is such a defense, not another formulation. Its informative content is that the defended principles about coexperience ($B_\psi$, the presence class, membership in a complex) are transversal to $\Phi_Q$ (Section~\ref{sec:intento}), so $S_u$ requires $\Phi_Q$ to \emph{segregate} the coexperience structures into bands. The threshold law TL is the two-group case, presence versus absence. Under UT, B$_\psi^{\rm pres}$, N and the absence anchor, the degree takes one value on the conscious realizations and, by B$_\psi^{\rm pres}$ and the anchor, a smaller one on the non-conscious ones; SEG with these two groups says that no fibre of $\Phi_Q$ contains both a conscious and a non-conscious realization, and together with the order condition B (Definition~\ref{def:orderB}) it is TL (Definition~\ref{def:TL}, Theorem~\ref{thm:threshold}). Refuting SEG through a counterexample to $S_u$ would require three things:
\begin{itemize}
\item two original realizations with the same computed value of $\Phi_Q$;
\item coexperience structures determined by a calibrated observation channel (Q's observable-assignment rule A14--U is a stipulation, not a calibration);
\item a specified $\psi$ that distinguishes them.
\end{itemize}
None of the three is available: $\Phi_Q$ has not been selected (Definition~\ref{def:phiQ}), so none of its values has been computed; no calibrated channel is available; and no $\psi$ has been adopted (Definition~\ref{def:unity}). (By Proposition~\ref{prop:seg}(b), SEG also fails if $u=\psi\circ\mathfrak U$ is constant on $D_0$; that would not refute $S_u$.)

\paragraph{Free premises and decisions.} The following are open decisions in the specification of Q, not results to be proved: the selection of $\Phi_Q$, of the domain $D_0$ and of the grain $\gamma$ (Section~\ref{sec:enunciado}); and the comparison relation for the degree, a function $\psi$ of the coexperience structure or a substitute, together with the correspondence of aspects that it uses (Definition~\ref{def:unity}). The following remain unjustified for the original domain: LSS$_\Delta$, OI$_\gamma$, Rob, UT, $B_\psi$, B$_\psi^{\rm pres}$, N, the absence anchor, TL and the selection law SEL (Definition~\ref{def:SEL}). A9--U and A14--U are stipulated by Q, not derived (Definition~\ref{def:orderB}), and no result of this paper takes one of the decisions above or justifies one of these premises. The obligations of Q themselves remain open: $S_u$, or SEG in its testable form (Definition~\ref{def:SEG}), and the order condition B (Definition~\ref{def:orderB}). Table~\ref{tab:kinds} separates the kinds of item on which a derivation of the obligations of Q would rest.

\begin{table}[htbp]\centering\small
\begin{tabular}{>{\raggedright\arraybackslash}p{0.21\textwidth}>{\raggedright\arraybackslash}p{0.45\textwidth}>{\raggedright\arraybackslash}p{0.24\textwidth}}
\toprule
Kind & Items & Status in this paper\\
\midrule
Specification decisions (the author's) & $D_0$; the grain $\gamma$; $\Phi_Q$, with its boundary, units, state, dynamics, time scale, partitions, distance and selection conventions; the comparison relation of $u$ (a $\psi$, or a substitute) with its correspondences $f$; the reading of $\mathcal C$ & not taken; each result holds for every choice that satisfies its stated hypotheses\\
Stipulations of Q & A9--U ($S_u$ and B) and A14--U (Definition~\ref{def:orderB}) & stipulated, neither derived nor justified\\
Substantive premises & $B_\psi$, LSS$_\Delta$, OI$_\gamma$, Rob, UT, B$_\psi^{\rm pres}$, N, the absence anchor, TL and SEL on $D_0$; IPS$_K$ on the test bench of Section~\ref{sec:gauss} & examined; none adopted or justified\\
Results & theorems and propositions labelled \stL, \stM, \stA\ or \stS; numerical illustrations \stN & conditional on the hypotheses that each one states\\
\bottomrule
\end{tabular}
\caption{The kinds of item on which a derivation of the obligations of Q would rest; the reading of $\mathcal C$ enters only $S_q$, and IPS$_K$ only the test bench. No result of the last row takes a decision or justifies a premise of the rows above it.}\label{tab:kinds}
\end{table}

\paragraph{The other obligations of Q.} Q also carries sufficiency obligations for the components $c$, $q$, $\tau$ and $\sigma$, the order clause $\mathrm B$, the descent condition O, joint admissibility and conditions on processes (Definition~\ref{def:phiQ}). Section~\ref{sec:conditional} states which premises would discharge the five sufficiencies, $\mathrm B$ and O, and establishes none of them; apart from it, this paper treats only $S_u$. All these obligations remain open.

\subsection{What premises would discharge the obligations}\label{sec:conditional}
The following theorem applies the entailment criterion of Theorem~\ref{thm:impl} to each component, extends it to premises of order type, and adds the descent condition O. It changes neither the domain, nor the descriptor, nor the components, nor the obligations of Q, and it adopts none of the hypotheses of Section~\ref{sec:enunciado}.

\begin{theorem}[Conditional discharge of the sufficiency, order and descent obligations]\label{thm:conditional-discharge}
\stA\ Let $D_0$, $\Phi_Q$, the components $x\in\{c,q,u,\tau,\sigma\}$, their reduced inputs $d_x$ and the assignment $\mathbf e$ be as in Definition~\ref{def:phiQ}. For a relation $R$ on $D_0$, let $R^\ast$ be its reflexive and transitive closure and $\langle R\rangle$ the equivalence relation that it generates. A function $y$ on $D_0$ \emph{respects} $R$ if $s\,R\,t$ implies $y(s)=y(t)$, and a real-valued $y$ is \emph{nondecreasing along} $R$ if $s\,R\,t$ implies $y(s)\le y(t)$.
\begin{enumerate}
\item[(i)] \emph{Sufficiency, component by component.} Let $R_x$ be a relation on $D_0$. If $x$ respects $R_x$ and $\ker d_x\subseteq\langle R_x\rangle$, then $S_x$ holds. Conversely, let $X$ be a set with at least two elements; if every function $y:D_0\to X$ that respects $R_x$ satisfies $d_x(s)=d_x(t)\Rightarrow y(s)=y(t)$, then $\ker d_x\subseteq\langle R_x\rangle$. In particular, if $R_x$ is an equivalence relation contained in $\ker d_x$, every such $y$ satisfies that implication if and only if $R_x=\ker d_x$.
\item[(ii)] \emph{Unity through the coexperience structure.} Assume $B_\psi$ (Definition~\ref{def:unity}). Then $S_u$ is the inclusion $\ker\Phi_Q\subseteq\ker(\psi\circ\mathfrak U)$: realizations with equal $\Phi_Q$ may have different coexperience structures, provided that $\psi$ takes equal values on them. The relational sufficiency $\ker\Phi_Q\subseteq\ker\mathfrak U$ implies $S_u$ for every comparison $\psi$; it is not necessary in general for $S_u$ with one given $\psi$, even when $u$ is not constant; it is necessary for $S_u$ with one $\psi$ that takes different values on any two distinct structures occurring in one fibre of $\Phi_Q$, for instance an injective $\psi$ (for such a $\psi$, $S_u$ is equivalent to it), and for $S_u$ with every $\psi$ (Theorems~\ref{thm:impl} and~\ref{thm:robusta}). If $\mathfrak U$ respects a relation $R$ and $\ker\Phi_Q\subseteq\langle R\rangle$, then $\ker\Phi_Q\subseteq\ker\mathfrak U$.
\item[(iii)] \emph{Order.} Let $P$ be a relation on $D_0$, and put each real-valued function on $D_0$ that is nondecreasing along $P$ in place of $u$. (a) All these functions satisfy $S_u$ if and only if $\ker\Phi_Q\subseteq P^\ast\cap(P^\ast)^{-1}$. (b) All of them satisfy $\mathrm B$ if and only if $\Phi_Q(s)<\Phi_Q(t)$ implies $(s,t)\in P^\ast$. (c) All of them satisfy $S_u\wedge\mathrm B$ if and only if $\Phi_Q(s)\le\Phi_Q(t)$ implies $(s,t)\in P^\ast$. In particular, if $u$ is nondecreasing along $P$, each of these conditions implies the corresponding obligation. For $P=R\cup R^{-1}$ one has $P^\ast=\langle R\rangle$, a real-valued function is nondecreasing along $P$ exactly when it respects $R$, and (a) is Theorem~\ref{thm:impl}. If $S_u$ holds and $h_u$ is the function on $\Phi_Q(D_0)$ with $u=h_u\circ\Phi_Q$, then $\mathrm B$ holds if and only if $h_u$ is nondecreasing on $\Phi_Q(D_0)$.
\item[(iv)] \emph{Interventions.} Let $\mathcal T$ be a class of pairs of realizations of $D_0$ and $P_{\mathcal T}=\{(s,t)\in\mathcal T:\Phi_Q(s)<\Phi_Q(t)\}$, so that $\mathrm{IM}_{\mathcal T}$~\eqref{eq:im} says that $u$ is nondecreasing along $P_{\mathcal T}$. The real-valued functions on $D_0$ that are nondecreasing along $P_{\mathcal T}$ all satisfy $S_u$ if and only if $\Phi_Q$ is injective on $D_0$, and they all satisfy $\mathrm B$ if and only if any $s,t\in D_0$ with $\Phi_Q(s)<\Phi_Q(t)$ are joined by a chain of pairs of $P_{\mathcal T}$.
\item[(v)] \emph{Descent.} Let $\sim$ be a relation on $D_0$. If, for every $x$, $S_x$ holds and $s\sim t$ implies $d_x(s)=d_x(t)$, then $s\sim t$ implies $\mathbf e(s)=\mathbf e(t)$, which is condition O for $\sim$; when $\sim$ is an equivalence relation, this says that $\mathbf e$ factors through $D_0/{\sim}$. Conversely, if $s\sim t$ and $d_x(s)\ne d_x(t)$ for some $x$, then for every set $X$ with at least two elements some $y:D_0\to X$ satisfies $d_x(r)=d_x(r')\Rightarrow y(r)=y(r')$ for all $r,r'\in D_0$ and $y(s)\ne y(t)$.
\end{enumerate}
\end{theorem}
\begin{proof}
(i) If $d_x(s)=d_x(t)$ and $\ker d_x\subseteq\langle R_x\rangle$, there is a chain $s=r_0,\dots,r_n=t$ in which $r_{k-1}\,R_x\,r_k$ or $r_k\,R_x\,r_{k-1}$ for each $k$; $x$ takes one value along it, so $x(s)=x(t)$. Conversely, let $a\ne b$ in $X$ and $d_x(s)=d_x(t)$. The function $y$ with $y(r)=a$ if $r$ lies in the class of $t$ under $\langle R_x\rangle$ and $y(r)=b$ otherwise respects $R_x$, since $R_x$-related realizations lie in one class; hence $y(s)=y(t)=a$, and $s$ lies in the class of $t$. If $R_x$ is an equivalence relation contained in $\ker d_x$, then $\langle R_x\rangle=R_x\subseteq\ker d_x$, and the inclusion $\ker d_x\subseteq\langle R_x\rangle$ is the equality.

(ii) The first statement is $S_u$ for $u=\psi\circ\mathfrak U$, written out. If $\ker\Phi_Q\subseteq\ker\mathfrak U$, equal descriptors give equal structures, hence equal values of every $\psi$. For the failure of the converse in general, take a finite model, which is not a specification of Q: three realizations $s,t,r$ with descriptor values $0,0,1$ and three distinct structures on which $\psi=0,0,1$: for instance $E=\{a,b\}$ for the three episodes, $H_s=\{\{a\}\}$, $H_t=\{\{b\}\}$, $H_r=\{\{a\},\{b\},\{a,b\}\}$, and $\psi=1$ exactly when two aspects are coexperienced. In this model the sufficiency of the descriptor for $u=\psi\circ\mathfrak U$ holds, $u$ is not constant, and the descriptor does not determine $\mathfrak U$. The necessity for every $\psi$ is Theorem~\ref{thm:robusta}, or Theorem~\ref{thm:impl} with $R=\ker\mathfrak U$, whose indicators of single structures need no injective comparison; for one $\psi$ that separates the structures of a fibre, $\psi(\mathfrak U(s))=\psi(\mathfrak U(t))$ already gives $\mathfrak U(s)=\mathfrak U(t)$; and the last statement is the first part of~(i) with $\mathfrak U$ in place of $x$ and $\Phi_Q$ in place of $d_x$.

(iii) If $y$ is nondecreasing along $P$, induction on the length of a chain gives $y(s)\le y(t)$ whenever $(s,t)\in P^\ast$; this gives the ``if'' directions, in~(a) by antisymmetry. For the ``only if'' directions, fix $s$ and let $y_s$ be the indicator of $\{r:(s,r)\in P^\ast\}$. It is nondecreasing along $P$, since $(s,r)\in P^\ast$ and $r\,P\,r'$ give $(s,r')\in P^\ast$, and $y_s(s)=1$. In~(b), $\Phi_Q(s)<\Phi_Q(t)$ and $\mathrm B$ for $y_s$ give $1\le y_s(t)$, so $(s,t)\in P^\ast$; (c) is the same with $\le$ in place of $<$; in~(a), $S_u$ for $y_s$ gives $y_s(t)=1$, and $S_u$ for $y_t$ gives $(t,s)\in P^\ast$. If $P=R\cup R^{-1}$, then $P^\ast$ is symmetric, hence the smallest equivalence relation that contains $R$, and $y(s)\le y(t)\le y(s)$ on $R$-related pairs means $y(s)=y(t)$. Finally, if $S_u$ holds, $u=h_u\circ\Phi_Q$ with $h_u$ unique on $\Phi_Q(D_0)$ (Definition~\ref{def:phiQ}(e)); if $a<b$ lie in $\Phi_Q(D_0)$, with $a=\Phi_Q(s)$ and $b=\Phi_Q(t)$, $\mathrm B$ gives $h_u(a)=u(s)\le u(t)=h_u(b)$, and conversely a nondecreasing $h_u$ gives $\mathrm B$. This last equivalence is the reading of A9--U in Definition~\ref{def:orderB}(b): it restates $\mathrm B$ under $S_u$ and does not derive it.

(iv) Since $<$ is transitive, induction on chains shows that $(s,t)\in P_{\mathcal T}^\ast$ implies $s=t$ or $\Phi_Q(s)<\Phi_Q(t)$. Hence two distinct realizations with equal $\Phi_Q$ are never joined by a chain of pairs of $P_{\mathcal T}$, and the condition of~(iii)(a) holds exactly when every fibre of $\Phi_Q$ on $D_0$ has one element. The statement about $\mathrm B$ is~(iii)(b) with $P=P_{\mathcal T}$.

(v) If $s\sim t$, then $d_x(s)=d_x(t)$ for every $x$, so $S_x$ gives $x(s)=x(t)$ for every $x$, that is, $\mathbf e(s)=\mathbf e(t)$; a function that is constant on the classes of an equivalence relation factors through the quotient. Conversely, with $a\ne b$ in $X$, the function $y$ with $y(r)=a$ if $d_x(r)=d_x(s)$ and $y(r)=b$ otherwise is a function of $d_x$, and $y(s)=a\ne b=y(t)$.
\end{proof}
\emph{Formal check} \stL. The abstract form of Theorem~\ref{thm:conditional-discharge} is checked in Lean without Mathlib, in \path{CondicionalQ.lean} (namespace \texttt{CondicionalQ}); Table~\ref{tab:corr-1} lists the theorems that check each part. The realizations, the values of the components and of their reduced inputs are arbitrary types; the order of the descriptor values is a relation, irreflexive and transitive in~(iv), and in~(iii)(c) $\Phi_Q(s)\le\Phi_Q(t)$ is read as $\Phi_Q(s)<\Phi_Q(t)$ or $\Phi_Q(s)=\Phi_Q(t)$; and in~(iii) and~(iv) the degrees take values in a scale, a partial order with two values $\mathrm{lo}<\mathrm{hi}$ (\texttt{Scale}). The integers form a scale (\texttt{intScale}), and so do the real numbers, but that instance is not compiled, since Lean without Mathlib has no real numbers. In~(ii), \texttt{single\_psi\_example} takes the three structures as abstract codes, not as families of coexperienced sets, and the necessity for one $\psi$ that separates the structures of a fibre is not a separate compiled statement. Parts~(i), (iii)(b) and~(v) combine into one statement, \texttt{conditional\_discharge}; in a model on three realizations its bridge, coverage and order premises hold with $u$ not constant (\texttt{conditional\_discharge\_nonvacuous}).

\paragraph{What the theorem gives and what it does not.} Theorem~\ref{thm:conditional-discharge} is a conditional result. It separates premises of two kinds from its conclusions. The premises that $x$ respects $R_x$ and that $u$ is nondecreasing along $P$ are psychophysical hypotheses, and the construction of $\sim$ independently of Q is a physical task; the theorem justifies none of them. The inclusions of~(i), (iii) and~(v) are the second kind: when $R_x$, $P$ and $\sim$ are defined physically, they are conditions on $D_0$ and on the reduced inputs that become determinate propositions once $D_0$, $\gamma$ and the $d_x$ are selected, which Q has not done (Section~\ref{sec:laguna}, (L0)). Neither kind is an empirical validation, which would need a calibrated measurement channel for the components (Section~\ref{sec:mejor}). For $x=u$ and premises of the form $u(s)=u(t)$, part~(i) is Theorem~\ref{thm:impl} and adds nothing to Theorems~\ref{thm:coverage}--\ref{thm:robusta}: under Rob, with $\theta\mapsto\Phi_Q(s_\theta)$ continuous and $\Phi_Q(s_0)$ interior to its range, no relation contained in ${\cong_\gamma}$ covers $\ker\Phi_Q$ (Theorem~\ref{thm:coverage} and Proposition~\ref{prop:extension}), and with $R_u=\ker\Phi_Q$ the premise is $S_u$ itself. Part~(iv) makes the status of $\mathrm{IM}_{\mathcal T}$ (Section~\ref{sec:im}) exact: as a premise it entails $S_u$ for every compatible degree only when $\Phi_Q$ is injective on $D_0$, and it entails $\mathrm B$ only between realizations joined by interventions that raise $\Phi_Q$, so not between realizations of different bearers when $\mathcal T$ relates states of one bearer. The theorem does not treat joint admissibility or the process obligations, which Definition~\ref{def:phiQ} names without a formal statement; no conditional theorem about them can be stated before they are specified.

\paragraph{The quality component under one reading of $\mathcal C$ \stA\ \stL.} Q leaves open whether $\mathcal C$ in $d_q=(\Phi_Q,\mathcal C)$ is the unpointed wiring or the pointed isomorphism class at grain $\gamma$ (Section~\ref{sec:enunciado}). Under the second reading, Adm gives $\ker d_q={\cong_\gamma}$: $S_q$ then says exactly that $q$ is constant on the isomorphism classes at grain $\gamma$, the coverage condition of part~(i) holds with $R_q={\cong_\gamma}$, and $\Phi_Q$ plays no role in $S_q$ (\texttt{CondicionalQ.\allowbreak pointed\_reading\_kernel}, \texttt{CondicionalQ.\allowbreak Sq\_pointed}). OI$_\gamma$, which is not adopted here, yields that constancy on the classes contained in $D_{\rm cons}$ when $\gamma$ is fine enough in the sense of Definition~\ref{hyp:OI}, as stated after that definition for $u$. Under the first reading, $\ker d_q$ can contain pairs that are not isomorphic, and $S_q$ meets the coverage problem of $S_u$. Which reading Q intends is a decision of the author.

\section{The categorical unification thesis}\label{sec:categorica}
\subsection{Base category}\label{sec:base}
\begin{definition}[Finite open system with a blanket]
An object is $X=(V,\rho,(\Sigma_v)_{v\in V},K)$, where:
\begin{itemize}
\item $V$ is a finite set of units, possibly empty;
\item $\rho:V\to\{\mu,s,a,\eta\}$ assigns the internal, sensory, active and external roles;
\item $\Sigma=\prod_v\Sigma_v$ is finite;
\item $K\in\mathbf{FinStoch}(\Sigma,\Sigma)$ is a Markov kernel factorized over units, $K(x'\mid x)=\prod_vK_v(x'_v\mid x_{\mathrm{pa}(v)})$, with the blanket condition: no internal unit has external parents and no external unit has internal parents~\cite{Friston2013}.
\end{itemize}
The morphisms are the isomorphisms: bijections of units that preserve the roles, with recodings of the local states that transport $K$ exactly.
\end{definition}
$X\otimes Y$ is the disjoint union of units with the product kernel; the unit of the product is the empty system, which has no units, and the symmetry is the swap. $\mathbf{Open}$ is thus a \emph{symmetric monoidal groupoid} that projects onto $\mathbf{FinStoch}$, the Markov category of finite sets and stochastic matrices~\cite{Fritz2020}. The internal block $(\mu,a)$, with input $s$, is a parametrized kernel $\Sigma_\mu\times\Sigma_s\to\mathrm{Dist}(\Sigma_\mu\times\Sigma_a)$, that is, an open system with interface $(\Sigma_s,\Sigma_a)$ in the sense of lenses and of $\mathbf{Poly}$~\cite{Capucci2022,NiuSpivak2023}. Coupling it with the environment amounts to closing the loop.
\begin{itemize}
\item \emph{Quantum version:} the kernels are replaced by CPTP channels and the roles by a decomposition into system and environment fragments.
\item \emph{Linear-Gaussian instance} $\mathbf{Open}_G$: an object with $n\ge0$ units, each with a role in $\{\mu,s,a,\eta\}$ as in $\mathbf{Open}$, is a pair $(A,E)$ of real $n\times n$ matrices, with dynamics $X_{t+1}=AX_t+\varepsilon_{t+1}$, independent noise $\varepsilon_{t+1}\sim N(0,E)$, spectral radius $\operatorname{spr}(A)<1$ and $E\succ0$. Its stationary covariance $P$ is the unique solution of $P=APA^\top+E$, so $P\succeq E$. The isomorphisms are permutations within each role, positive rescalings and sign changes, and $\otimes$ is the block-diagonal sum; its unit is the empty system ($n=0$). $J$ and the broadcast hubs do not depend on the roles; the roles enter the isomorphisms and the FEP-type projection, through $\rho$ and $I(\mu;\eta)$ (Proposition~\ref{prop:monoidal}).
\end{itemize}
\paragraph{Four-unit examples.} Sections~\ref{sec:categorica} and~\ref{sec:debil} use four systems of $\mathbf{Open}_G$ on four units, one for each graph $\Gamma$ among the ring (the cycle) $C_4$, the star $K_{1,3}$ (one unit joined to the other three), the path $P_4$, and the split system $2K_2$ (two disjoint edges): $A=0.2\,I+0.3\,M_\Gamma$, with $M_\Gamma$ the adjacency matrix of $\Gamma$, and $E=I$. Section~\ref{sec:networks} studies them numerically. In general, $P_n$ and $C_n$ denote the path and the cycle on $n$ units.

\subsection{The theories as projection functors}\label{sec:proyecciones}
\begin{proposition}\label{prop:funtor}
An assignment $F:\mathrm{Ob}(\mathbf{Open})\to D$ defines a functor $\mathbf{Open}\to\mathrm{Disc}(D)$ if and only if it is invariant under isomorphism. The same holds for $\mathbf{Open}_G$.
\end{proposition}
\begin{proof}
A functor into the discrete category $\mathrm{Disc}(D)$ sends every morphism $X\to Y$ to an identity, so $F(X)=F(Y)$ whenever $X\cong Y$. Conversely, if $F$ is invariant under isomorphism, sending every morphism $X\to Y$ to the identity of $F(X)=F(Y)$ preserves identities and composition.
\end{proof}
Invariance under isomorphism is the analog, for objects without a distinguished state, of the admissibility condition Adm of Definition~\ref{def:adm}, which is formulated for pointed causal models.

Computable instances:
\begin{itemize}
\item $F_{\rm IIT}^J(X)=J(X)$, the Gaussian residual functional of Lemma~\ref{lem:J}, defined on $\mathbf{Open}_G$ with the conventions of Remark~\ref{rem:small} for systems with fewer than two units; $F_{\rm IIT}$ denotes an IIT-type projection in general, and $F_{\rm IIT}^J$ is its instance here. It is \emph{not} the $\Phi$ of IIT~3.0, which equals $0$ on strictly unidirectional systems~\cite{Oizumi2014}; $J>0$ on every system with at least two units and connected support (Lemma~\ref{lem:J}).
\item $F_{\rm GNWT}(X)=\lvert\mathrm{Hub}(X)\rvert\in\mathbb N$, the number of broadcast hubs of $X$, where $\mathrm{Hub}(X)$ is the set of units adjacent to all the others in the undirected graph of the support of $A$. An isomorphism of $\mathbf{Open}_G$ permutes the units within each role, by a bijection $\sigma$, and rescales them by nonzero factors $d_v$ (Section~\ref{sec:base}); the coefficient of the image of $A$ at $(\sigma(v),\sigma(w))$ is $d_vA_{vw}/d_w$, so the isomorphism carries the support graph of $X$ onto that of its image along $\sigma$ and restricts to a bijection between the two sets of hubs. Thus $\mathrm{Hub}$ is a functor into the groupoid of finite sets and bijections; by Proposition~\ref{prop:funtor} it is not a functor into a discrete category, since the set of hubs is not invariant under isomorphism (exchanging the centre of the star $K_{1,3}$ with a leaf of the same role moves it). The number of hubs, $F_{\rm GNWT}$, and their presence, $G_{\rm GNWT}F_{\rm GNWT}(X)=[\,\mathrm{Hub}(X)\ne\varnothing\,]$ (Section~\ref{sec:puentes}), are isomorphism invariants. Lean formalizes this reading only through the presence of a hub, as a Boolean (\texttt{values}, \texttt{n3\_equiv}), and the absence of hubs (\texttt{path\_no\_universal}, \texttt{tensor\_no\_universal}); no Lean statement takes the set of hubs as a value or counts hubs. The hub is a structural stand-in for global availability, a property of the support of $A$ and not of the dynamics: neither the nonlinear ignition~\cite{DehaeneChangeux2011} nor the dynamics of broadcast is modeled. In the systems of Table~\ref{tab:C-1} the couplings are symmetric, and requiring coupling in both directions would give the same result. With directed couplings, the undirected reading counts as a hub a unit that receives from all the others and sends to none, the opposite of a global broadcast; as a stand-in for broadcast it can be adequate only on symmetric supports, where it coincides with the readings that require a hub to send to every other unit or to be coupled with it in both directions. With directed couplings both demanding readings differ from the undirected one: requiring coupling in both directions, the two-unit family of Section~\ref{sec:gauss} has no hub (\texttt{f0\_hub\_readings}) and lies in the band at $n=2$; requiring a hub to send to every other unit, its driver is a hub, but the directed path on three units has none (Section~\ref{sec:summary}, \emph{Limitations}).
\item $F_{\rm FEP}(X)=(\rho,\ p_X,\ I_{p_X}(\mu;\eta))$, where $p_X$ is a stationary distribution of $K$ (unique if $K$ is ergodic; Proposition~\ref{prop:monoidal}) and $I_{p_X}(\mu;\eta)$ is the mutual information between the internal and the external units under $p_X$. Like the set of hubs, the pair $(\rho,p_X)$ is labelled by the units and their states, so it is not an isomorphism invariant; when $p_X$ is unique, the FEP-type assignment into a discrete category is the isomorphism class of $(\rho,p_X)$, with $I_{p_X}(\mu;\eta)$, or $I_{p_X}(\mu;\eta)$ alone (Section~\ref{sec:tesis}). The bare predicate ``has a blanket'' is constant on $\mathbf{Open}$, because every object has a blanket by construction.
\item $F_{\rm QD}(X)$: the redundancy profile $\delta\mapsto R_\delta$ of a pointer variable over a partition of $\eta$ into fragments~\cite{Zurek2009}, where $R_\delta$ is the number of disjoint fragments each of which carries at least a fraction $1-\delta$ of the information about the pointer variable.
\item $F_{\rm AST}(X)$: the existence of an internal subsystem whose state carries information about an attentional gain. AST deals with the attribution of consciousness~\cite{GrazianoKastner2011}.
\end{itemize}

\begin{lemma}[Gaussian integration and connectivity]\label{lem:J}
Let $X=(A,E)$ be a system of $\mathbf{Open}_G$ with $n\ge2$ units and stationary covariance $P$, and let
\[
J(A,E)=\min_\pi\min_{B\in\mathfrak B_\pi}\tfrac12\log\frac{\det\bigl(E+(A-B)P(A-B)^\top\bigr)}{\det E},
\]
where $\pi$ ranges over the bipartitions of the units (partitions into two nonempty blocks) and $\mathfrak B_\pi$ is the subspace of $n\times n$ matrices with no coefficients crossing $\pi$, that is, with $B_{vw}=0$ whenever $v$ and $w$ lie in different blocks. The minima are attained and $0\le J<+\infty$. Moreover, $J>0$ if and only if the undirected graph of the off-diagonal support of $A$ is connected.
\end{lemma}
\begin{proof}
Let $M=(A-B)P(A-B)^\top\succeq0$. Then $\det(E+M)/\det E=\det(I+E^{-1/2}ME^{-1/2})\ge1+\mathrm{tr}(E^{-1}M)\ge1$, because $\det(I+Y)\ge1+\operatorname{tr}Y$ for $Y\succeq0$; hence $J\ge0$. Moreover, $\mathrm{tr}(E^{-1}M)\ge\lambda_{\min}(P)\lambda_{\max}(E)^{-1}\lVert A-B\rVert_F^2$, with $\lambda_{\min}(P)\ge\lambda_{\min}(E)>0$ because $P\succeq E$. So the objective is continuous and coercive in $B$, its infimum over $\mathfrak B_\pi$ is attained, and, since there are finitely many bipartitions, $J$ is attained and finite. Fix a bipartition $\pi$. If $A\notin\mathfrak B_\pi$, then for every $B\in\mathfrak B_\pi$ the norm $\lVert A-B\rVert_F$ is at least the norm of the crossing coefficients of $A$, which is positive, so the minimum over $B$ is positive. If $A\in\mathfrak B_\pi$, taking $B=A$ gives $0$. Hence $J=0$ if and only if $A\in\mathfrak B_\pi$ for some $\pi$, that is, if and only if some bipartition is crossed by no off-diagonal coefficient of $A$, which means that the undirected graph of the off-diagonal support is disconnected.
\end{proof}
\emph{Status.} \stA\ \stM. Besides the paper proof above, the lemma is checked in Lean with Mathlib for $n\ge2$. \texttt{QAudit.lemmaJ\_stationary} states it with the hypotheses above: $E\succ0$, and $P\succ0$ with $P=APA^\top+E$, taken as data. \texttt{QAudit.lemmaJ} proves the same equivalence for every positive definite $P$, without the stationarity equation, and \texttt{QAudit.J\_eq\_zero\_iff} states that $J=0$ exactly when the support is disconnected, so that the two together give $J\ge0$. The stationary covariance can also be constructed: if $\operatorname{spr}(A)<1$, the series $\sum_{j\ge0}A^jE(A^j)^\top$ converges (\texttt{QAuditScaling.serSummable\_of\_spr\_lt\_one}), and its sum is positive definite and solves $P=APA^\top+E$ (\texttt{QAuditBridge.stationary\_of\_spr}), so the lemma holds with it (\texttt{QAuditBridge.lemmaJ\_spr}); here $\operatorname{spr}$ is Mathlib's spectral radius of $A$ regarded as a complex matrix. The uniqueness of the stationary covariance is not formalized, and the theorems do not need it: they hold for every positive definite solution. In Lean the inner minimum is an infimum over $\mathfrak B_\pi$, whose attainment is a separate checked theorem (\texttt{QAudit.exists\_minimizer}), and \texttt{QAudit.exists\_J\_minimizer} gives a bipartition and a matrix $B$ at which $J$ itself is attained. The convention $J=+\infty$ for one unit (Remark~\ref{rem:small}) is not formalized: in Lean, $J=0$ when there is no bipartition, and the lemma is checked only for $n\ge2$.

\begin{remark}[Systems with fewer than two units]\label{rem:small}
Bipartitions exist only when there are at least two units, so Lemma~\ref{lem:J}, the min-cut sandwich (Theorem~\ref{thm:cheeger}) and the decomposition $I_{\rm pred}=I_{\rm cut}+J$ of Section~\ref{sec:via1} are stated for systems with $n\ge2$ units, where $J$ and the minimum cut $w_{\min}$ are finite. For the two remaining cases we fix the following conventions.
\begin{itemize}
\item \emph{One unit} ($n=1$). There is no bipartition, and we set $J=w_{\min}=\min\varnothing=+\infty$. The graph of the support, a single vertex, is connected, and the unit is vacuously adjacent to all the others, so it is a broadcast hub. Both presence predicates, $J>0$ and the existence of a hub, take the value $1$, and the equivalence of Lemma~\ref{lem:J} also holds for $n=1$.
\item \emph{The empty system} ($n=0$), the unit of $\otimes$. It has no units and no bipartitions. The convention $\min\varnothing=+\infty$ is not applied to it: $J$ and $w_{\min}$ are not defined on the empty system, and the IIT-type projection sends it to an adjoined unit $e\notin[0,+\infty]$, the unit of the zero-product monoid of Section~\ref{sec:l2}. Its set of broadcast hubs is empty; Section~\ref{sec:l2}, where monoidality is stated, also sends it to $e$ under the GNWT-type projection. The empty system belongs to none of the classes defined below by conditions on $J$, $w_{\min}$ or hubs, and the bridge square and the filtration of Sections~\ref{sec:puentes} and~\ref{sec:l1} are considered on the full subcategory $\mathbf{Open}_G^{+}$ of nonempty systems.
\end{itemize}
Thus $F_{\rm IIT}^J$ takes values in $[0,+\infty]\sqcup\{e\}$: $J(X)\in[0,+\infty)$ if $X$ has at least two units, $+\infty$ if it has one, and $e$ if $X$ is empty. It is invariant under isomorphism, hence a functor by Proposition~\ref{prop:funtor}: an isomorphism acts by $A\mapsto TAT^{-1}$ and $E\mapsto TET^\top$, with $T$ a role-preserving permutation matrix with nonzero scalings, so it preserves the number of units, sends $P$ to $TPT^\top$ and $\mathfrak B_\pi$ to $\mathfrak B_{\pi'}$ for the permuted bipartition $\pi'$, and multiplies both determinants in the definition of $J$ by $\det(T)^2$.
\end{remark}

\begin{proposition}[Monoidal behavior]\label{prop:monoidal}
For nonempty $X,Y$ one has $J(X\otimes Y)=0$, and $X\otimes Y$ has no broadcast hub (\texttt{tensor\_no\_universal}). The predictive information $I_{\rm pred}=\tfrac12\log\det P-\tfrac12\log\det E$ and the mutual information $I(\mu;\eta)$ are additive under $\otimes$ in $\mathbf{Open}_G$. In $\mathbf{Open}$, $I(\mu;\eta)$ is additive if $p_X$ is the unique stationary distribution of an ergodic (irreducible and aperiodic) kernel, a property that is preserved under $\otimes$, or if one stipulates $p_{X\otimes Y}=p_X\otimes p_Y$. Without aperiodicity it fails: two swap kernels $x\mapsto1-x$, one internal and the other external, each have a unique stationary distribution, but their product admits the stationary distribution $\tfrac12(\delta_{00}+\delta_{11})$, with $I(\mu;\eta)=\log2$ versus $0+0$.
\end{proposition}
\begin{proof}
If $X$ and $Y$ are nonempty, $X\otimes Y$ has at least two units and its matrix $\mathrm{diag}(A_X,A_Y)$ has no coefficient crossing the bipartition into the units of $X$ and those of $Y$; taking $B=\mathrm{diag}(A_X,A_Y)$ gives $J(X\otimes Y)=0$. No unit of $X$ is adjacent to a unit of $Y$, so no unit is adjacent to all the others; \texttt{tensor\_no\_universal} checks this for the Boolean adjacency of the block sum of two nonempty graphs. In $\mathbf{Open}_G$ the stationary covariance and the noise covariance of $X\otimes Y$ are $\mathrm{diag}(P_X,P_Y)$ and $\mathrm{diag}(E_X,E_Y)$, so $\log\det P$ and $\log\det E$ are additive. The stationary law of $X\otimes Y$ is the product of those of $X$ and $Y$, so the internal and external units of $X$ are independent of those of $Y$, and $I(\mu;\eta)$ is additive, as mutual information is for independent pairs. In $\mathbf{Open}$, if $K_X$ and $K_Y$ are irreducible and aperiodic, then for all large $m$ the matrices $K_X^m$ and $K_Y^m$ have all entries positive, and so does $(K_X\otimes K_Y)^m=K_X^m\otimes K_Y^m$; hence $K_X\otimes K_Y$ is irreducible and aperiodic, its unique stationary distribution is $p_X\otimes p_Y$, and additivity follows as before. In the counterexample each swap fixes only the uniform distribution, while the product kernel exchanges the states $00$ and $11$ and therefore fixes $\tfrac12(\delta_{00}+\delta_{11})$, under which the internal and the external unit are equal and uniform.
\end{proof}
Thus the IIT- and GNWT-type projections, $J$ and the presence (or the number) of broadcast hubs, are not monoidal functors into the additive monoid $([0,+\infty],+,0)$, whatever value they take on the empty system: they vanish on the tensor ideal of Section~\ref{sec:via1}, which contains every product of nonempty systems, and they do not vanish on $K_{1,3}$, so $F(K_{1,3}\otimes K_{1,3})=0\ne F(K_{1,3})+F(K_{1,3})$. With the unit convention of Section~\ref{sec:l2} (the empty system goes to the unit and no nonempty system does), they are strong monoidal functors into the zero-product monoid of that section, and without it they are not. The extensive projections are monoidal functors into $(\mathbb R_{\ge0},+,0)$ in $\mathbf{Open}_G$, and in $\mathbf{Open}$ under the conditions of Proposition~\ref{prop:monoidal}. Every functor out of the groupoid preserves sequential composition; into an additive monoid, parallel composition is preserved by the extensive projections and not by the IIT- and GNWT-type projections. On processes, additive trajectory functionals, such as entropy production, are functors from the path category into $\mathbf B(\mathbb R,+)$, and state functions give coboundaries.

\subsection{Translation between theories}\label{sec:traduccion}
Since $F_i$ and $F_j$ land in different categories, a translation is a functor $\tau:D_i\to D_j$ (unrelated to the temporality component $\tau$ of Q) with $\tau\circ F_i=F_j$; with discrete codomains, the natural transformation is an equality.
\begin{theorem}[\texttt{nat\_iff\_functional}]\label{thm:translation}
Let $D_j$ be nonempty. Such a $\tau$ exists if and only if $F_j$ is constant on the fibres of $F_i$.
\end{theorem}
There is no translation from positive integration to broadcast (\texttt{no\_translation}: the ring and the star both have $J>0$ and differ in broadcast). Nor is there a translation in the reverse direction: the ring and the split system both lack a broadcast hub and differ in $J$ (by \texttt{values} and \texttt{nat\_iff\_functional}, with $J>0$ on the ring and $J=0$ on the split system by Lemma~\ref{lem:J}; this direction is not a separate compiled statement). The general translation is the relation $R_{ij}=\mathrm{im}\langle F_i,F_j\rangle$, a span $D_i\leftarrow\mathbf{Open}\to D_j$; the relation is single-valued exactly when $\tau$ exists.

\subsection{Reconstruction: the limit}\label{sec:limite}
For the discrete diagram $\mathfrak D=(D_i)_i$ one has $\lim\mathfrak D=\prod_iD_i$, and the cone $(\mathbf{Open},F_i)$ factors uniquely through it (\texttt{pairing\_unique}):
\begin{equation}
\begin{CD}
D_{\rm IIT} @<{\pi_1}<< D_{\rm IIT}\times D_{\rm GNWT}\times D_{\rm FEP} @>{\pi_2}>> D_{\rm GNWT}\\
@A{F_{\rm IIT}}AA @AA{F_\times=\langle F_{\rm IIT},F_{\rm GNWT},F_{\rm FEP}\rangle}A @AA{F_{\rm GNWT}}A\\
\mathbf{Open} @= \mathbf{Open} @= \mathbf{Open}
\end{CD}
\end{equation}
\begin{theorem}[Reconstruction; \texttt{reconstruction\_iff\_injective}]\label{thm:reconstruccion}
The isomorphism classes of objects of $\mathbf{Open}$ can be recovered from the joint projection, that is, there is a map $R:\prod_iD_i\to\mathrm{Ob}(\mathbf{Open})/{\cong}$ with $R(F_\times(X))=[X]$ for every object $X$, if and only if $F_\times$ is injective on isomorphism classes. There are blind spots:
\begin{itemize}
\item[(i)] with Boolean presence projections, the ring and the path have the same projection (\texttt{blind\_spot}) and are not isomorphic, because their degree multisets differ: $\{2,2,2,2\}$ versus $\{1,1,2,2\}$ (\texttt{blind\_spot\_multiset}; that the degree multiset is an isomorphism invariant is standard and is not formalized);
\item[(ii)] with a finite number $m$ of continuous scalar projections on an open family of dimension $d>m$, some fibre has positive dimension, by Hurewicz's dimension-lowering theorem applied to a compact ball~\cite[Theorem~VI\,7]{HurewiczWallman1941}. That fibre is compact and of positive dimension, hence it contains a nondegenerate continuum. If the isomorphisms act with finite orbits, as in $\mathbf{Open}$, where the recodings of finite states form a finite group, the fibre therefore contains a continuum of classes. In $\mathbf{Open}_G$ the positive rescalings give orbits of dimension $n$, and $d$ must be counted in the space of classes. For $m=1$, Theorem~\ref{thm:continuum} gives more: on a relatively open convex set of dimension at least two, acted on by a finite group of relabelings, the fibre of a continuous invariant projection at \emph{every} interior value contains uncountably many classes.
\end{itemize}
\end{theorem}
The criterion recovers objects up to isomorphism, not the groupoid. A functor into a discrete category sends every morphism to an identity, so $F_\times$ sends every automorphism to an identity, even when it is injective on classes, and it is an equivalence onto its image only if every automorphism group is trivial. A groupoid with one object and a nontrivial automorphism group has a single isomorphism class, so every projection is injective on its classes, but it is not equivalent to the terminal category; in $\mathbf{Open}_G$ the sign change of all units, $T=-I$, fixes every $(A,E)$ and is a nonidentity automorphism of every nonempty system. Recovering $\mathbf{Open}$ up to equivalence would also require the automorphism group of each class, up to isomorphism, since a groupoid is equivalent to the disjoint union, over its isomorphism classes, of the one-object groupoids of these groups; none of the projections considered here records these groups. \stA\ Thus $\mathbf{Open}$ is not the limit: it is the apex of a cone whose comparison with the limit loses the automorphisms, and loses the distinction between classes at the blind spots.

\subsection{Bridge layer: a concrete bridge square}\label{sec:puentes}
Let $G_{\rm IIT}(j)=[\,j>0\,]$ for $j\in[0,+\infty]$ and $G_{\rm GNWT}(m)=[\,m>0\,]$ for a number $m\in\mathbb N$ of hubs, with values in the poset $\mathbf 2=\{0<1\}$, so that $G_{\rm GNWT}F_{\rm GNWT}(X)=[\,\mathrm{Hub}(X)\ne\varnothing\,]$ is the presence of a hub. On the full subcategory $\mathbf{Open}_G^{+}$ of nonempty systems (Remark~\ref{rem:small}), the bridge square is
\begin{equation}
\begin{CD}
\mathbf{Open}_G^{+} @>{F_{\rm IIT}^J}>> [0,+\infty]\\
@V{F_{\rm GNWT}}VV @VV{G_{\rm IIT}}V\\
\mathbb N @>{G_{\rm GNWT}}>> \mathbf 2
\end{CD}
\end{equation}
On one-unit systems both paths give $1$ (Remark~\ref{rem:small}). The square commutes on the full subcategory $\{K_{1,3},2K_2\}$ (\texttt{square\_commutes\_on\_sub}) and does not commute on $C_4$ or on $P_4$ (\texttt{values}; \texttt{square\_fails} states the failure, with $C_4$ as witness). The values of $J$ are numerical (BFGS with restarts) and not a certificate \stN: each is the least value found, an upper bound for $J$. Their positivity and their vanishing follow from Lemma~\ref{lem:J}, not from these values: an upper bound cannot show that $J>0$, and the value $0$ of the split system is exact because its support is disconnected.

\begin{table}[htbp]\centering
\begin{tabular}{lcccc}
\toprule
system ($c=0.3$, $a=0.2$, $E=I$) & $J$ & $G_{\rm IIT}$ & $G_{\rm GNWT}$ & commutes?\\
\midrule
ring $C_4$ & $0.2213$ & 1 & 0 & no\\
star $K_{1,3}$ & $0.1135$ & 1 & 1 & yes\\
path $P_4$ & $0.1059$ & 1 & 0 & no\\
split $2K_2$ & $0$ & 0 & 0 & yes\\
\bottomrule
\end{tabular}
\caption{The bridge square on the four-unit systems $C_4$, $K_{1,3}$, $P_4$ and $2K_2$ of Section~\ref{sec:base}: values of $J$ (numerical upper bounds \stN\ for $C_4$, $K_{1,3}$ and $P_4$; the value $0$ of $2K_2$ is exact; all signs are given by Lemma~\ref{lem:J}) and of the presence bridges $G_{\rm IIT}$ and $G_{\rm GNWT}$, and whether the square commutes.}\label{tab:C-1}
\end{table}

\begin{theorem}[Bridge consistency; \texttt{bridge\_iff\_agree}]\label{thm:bridge}
There exists a phenomenal assignment that factors through both bridges if and only if the square commutes on the whole domain. If it exists, it is unique, since it is $G_{\rm IIT}\circ F_{\rm IIT}$; the Lean statement does not include uniqueness, which is immediate.
\end{theorem}
If each theory deals with a different component, $\Psi=\langle G_j\circ d_j\rangle_j$ always exists, by the universal property of the product: this is the architecture of Q, in which each component has its own reduced input $d_j$ (Definition~\ref{def:phiQ}). What has content are the sufficiencies, joint admissibility and the processes (Definition~\ref{def:phiQ}).

\subsection{What the thesis can claim}\label{sec:tesis}
\begin{itemize}
\item \textbf{Weak thesis (proved as a scheme):} every family of isomorphism-invariant assignments gives functors out of one and the same symmetric monoidal groupoid of open systems, and their joint projection exists and is unique (Proposition~\ref{prop:funtor}, \texttt{pairing\_unique}). It applies to stand-ins for the theories, not to their own invariants. $J$, the presence of a broadcast hub and the number of hubs are isomorphism invariants; the set of hubs itself is transported along isomorphisms, as a functor into finite sets and bijections. $F_{\rm QD}$ and $F_{\rm AST}$ are such assignments only when the pointer variable and the partition into fragments, and the attentional gain, are fixed in an invariant way. $F_{\rm FEP}$ is one if one takes the isomorphism class of $(\rho,p_X)$, or only $I_{p_X}(\mu;\eta)$. The weak thesis is cheap, because it holds for any family of admissible invariants, and it settles no empirical dispute.
\item \textbf{Strong thesis (not proved):} the bridge laws are mutually consistent and sufficient for phenomenal structure. It amounts to two things: that the overlapping squares commute on the domain, which fails on $C_4$ and $P_4$ for the IIT- and GNWT-type presence bridges; and that the joint projection be sufficient, which is the content of the obligations of Q.
\end{itemize}
For the proxies studied here, the framework does not dissolve the conflict: it locates their disagreement in squares that do not commute, and whether the theories themselves disagree there is an empirical question. This is consistent with the IIT--GNWT adversarial collaboration, whose results ``align with some predictions of IIT and GNWT, while substantially challenging key tenets of both theories''~\cite[abstract]{Cogitate2025}.

\section{Weak unification: four routes and four lines}\label{sec:debil}
\subsection{Route 1: the monoidal defect as integration}\label{sec:via1}
Let $X=(A,E)$ be a system of $\mathbf{Open}_G$ with stationary covariance $P$. For an $n\times n$ matrix $B$, used as a linear predictor of $X_{t+1}$ from $X_t$ with residual covariance $E+(A-B)P(A-B)^\top$, let $I_B=\tfrac12\log\det P-\tfrac12\log\det\bigl(E+(A-B)P(A-B)^\top\bigr)$ be the predictive information that $B$ captures. With $B=A$ one has $I_A=I_{\rm pred}=I(X_t;X_{t+1})$.
\begin{proposition}[Extensive decomposition plus defect]\label{prop:decomposition}
For systems with at least two units, $I_{\rm pred}=I_{\rm cut}+J$, with $I_{\rm cut}=\max_\pi\max_{B\in\mathfrak B_\pi}I_B$. (With one unit there is no bipartition, $J=+\infty$ by the convention of Remark~\ref{rem:small} and $I_{\rm cut}=\max\varnothing=-\infty$, so the identity is undefined; on the empty system neither $J$ nor $I_{\rm cut}$ is defined.) Moreover:
\begin{itemize}
\item $I_{\rm pred}$ is additive under $\otimes$, and it vanishes on the empty system;
\item $J$ vanishes exactly on the systems with at least two units and disconnected support (Lemma~\ref{lem:J}), which form a tensor ideal $\mathcal Z$: $X\in\mathcal Z$ implies $X\otimes Y\in\mathcal Z$ for every $Y$. $\mathcal Z$ contains the decomposable objects $X\otimes Y$ ($X,Y$ nonempty) and coincides with them only when $E$ is block diagonal along a cut without crossing coefficients; with noise correlated between the blocks, $J=0$ without the system being a tensor product (Section~\ref{sec:l2});
\item for nonempty $X,Y$, $I_{\rm cut}(X\otimes Y)=I_{\rm pred}(X)+I_{\rm pred}(Y)$; if $X$ and $Y$ have at least two units, this equals $I_{\rm cut}(X)+I_{\rm cut}(Y)+J(X)+J(Y)\ge I_{\rm cut}(X)+I_{\rm cut}(Y)$, with defect $J(X)+J(Y)$.
\end{itemize}
\end{proposition}
\begin{proof}
By the definition of $J$, the minimum of $\tfrac12\log\det\bigl(E+(A-B)P(A-B)^\top\bigr)$ over the bipartitions $\pi$ and the predictors $B\in\mathfrak B_\pi$ is $\tfrac12\log\det E+J$, so $I_{\rm cut}=\tfrac12\log\det P-\tfrac12\log\det E-J=I_{\rm pred}-J$. The additivity of $I_{\rm pred}$ is Proposition~\ref{prop:monoidal}; on the empty system both determinants are empty and equal to $1$. The vanishing set of $J$ is given by Lemma~\ref{lem:J}. If a bipartition $(V_1,V_2)$ of the units of $X$ is crossed by no coefficient, then so is the bipartition $(V_1\cup V_Y,V_2)$ of the units of $X\otimes Y$, where $V_Y$ are the units of $Y$; and for nonempty $X,Y$ the bipartition into the units of $X$ and those of $Y$ is crossed by none. Finally, for nonempty $X,Y$, $J(X\otimes Y)=0$ (Proposition~\ref{prop:monoidal}), so $I_{\rm cut}(X\otimes Y)=I_{\rm pred}(X\otimes Y)=I_{\rm pred}(X)+I_{\rm pred}(Y)$, and $I_{\rm pred}=I_{\rm cut}+J$ for $X$ and $Y$ with at least two units gives the rest.
\end{proof}
$J$ is the minimum divergence, over the cuts, from the manifold of disconnected predictors, which is the information-geometric reading~\cite{Oizumi2016PNAS}. The projections of IIT and GNWT type factor through the operation that sends $\mathcal Z$ to zero. The expression ``exact sequence'' can be used here only as an additive splitting in the monoid $(\mathbb R_{\ge0},+)$: there is no abelian structure. Calling $I_{\rm pred}$ the ``FEP base term'' is an interpretation: what is proved is extensivity.

\begin{theorem}[Min-cut sandwich]\label{thm:cheeger}
Let $X=(A,E)$ have at least two units, let $w_{\min}=\min_\pi\lVert A_{\rm cross}(\pi)\rVert_F^2$, where $A_{\rm cross}(\pi)$ keeps only the coefficients of $A$ that cross $\pi$, and let $\kappa_E=\lambda_{\max}(E)/\lambda_{\min}(E)$. Then
\begin{equation}
\tfrac12\log\bigl(1+w_{\min}/\kappa_E\bigr)\ \le\ J\ \le\ \tfrac12\,\frac{\lambda_{\max}(P)}{\lambda_{\min}(E)}\,w_{\min}.
\end{equation}
\end{theorem}
\begin{proof}
\emph{Lower bound.} For every bipartition $\pi$ and every $B\in\mathfrak B_\pi$, with $M=(A-B)P(A-B)^\top$, the proof of Lemma~\ref{lem:J} gives $\det(E+M)/\det E\ge1+\lambda_{\min}(P)\lambda_{\max}(E)^{-1}\lVert A-B\rVert_F^2$. Moreover, $\lambda_{\min}(P)\ge\lambda_{\min}(E)$, and $\lVert A-B\rVert_F\ge\lVert A_{\rm cross}(\pi)\rVert_F$ because $B$ vanishes at the crossing positions. Hence the ratio is at least $1+\lVert A_{\rm cross}(\pi)\rVert_F^2/\kappa_E\ge1+w_{\min}/\kappa_E$.
\emph{Upper bound.} Take a cut $\pi^*$ that attains $w_{\min}$ and $B=A-A_{\rm cross}(\pi^*)\in\mathfrak B_{\pi^*}$, so that $M=A_{\rm cross}(\pi^*)PA_{\rm cross}(\pi^*)^\top$. Since $\log\det(I+Y)\le\operatorname{tr}Y$ for $Y\succeq0$, one gets $J\le\tfrac12\operatorname{tr}(E^{-1}M)\le\tfrac12\lambda_{\min}(E)^{-1}\lambda_{\max}(P)\lVert A_{\rm cross}(\pi^*)\rVert_F^2$.
\end{proof}
It was checked numerically on the four systems and on twelve random systems with random $E$, of which eleven have connected support; the remaining one, the first drawn, has no off-diagonal coupling, so both bounds and $J$ vanish (\texttt{pistas\_numerico.py}). It resembles the Cheeger inequality in that it bounds a global quantity by a cut, but no spectral gap or algebraic connectivity is involved. Efficient searches for the minimum-information bipartition of geometric integrated information exist: Queyranne's algorithm for symmetric submodular functions, applied although the measure is not submodular~\cite{Kitazono2018}, and spectral clustering of the correlation matrix~\cite{TokerSommer2019}.

\emph{Status.} \stA\ \stM. Besides the paper proof above, the theorem is checked in Lean with Mathlib exactly as stated, with $\lambda_{\min}$ and $\lambda_{\max}$ the extreme eigenvalues, for $n\ge2$, $E\succ0$ and $P\succ0$ given as data with $P=APA^\top+E$ (\texttt{QAudit.sandwich}). A version with arbitrary Loewner bounds in place of the eigenvalues, $e_1I\preceq E\preceq e_2I$ with $e_1>0$ and $P\preceq p_2I$, is also checked (\texttt{QAudit.sandwich\_lower}, \texttt{QAudit.sandwich\_upper}); the upper bound needs only $P\succeq0$, not the stationarity equation. With the stationary covariance constructed as the series $\sum_{j\ge0}A^jE(A^j)^\top$ when $\operatorname{spr}(A)<1$, the sandwich is \texttt{QAuditBridge.sandwich\_spr}. As in Lemma~\ref{lem:J}, the inner minimum is an infimum whose attainment is checked separately, and the case of one unit, with $w_{\min}=+\infty$ (Remark~\ref{rem:small}), is not formalized.

\subsection{Route 2: a lax 2-cell between broadcast and integration}\label{sec:via2}
The target is enriched over the poset $\mathbf 2$: on $\mathbf{Open}_G^{+}$, a 2-cell $\alpha:G_{\rm GNWT}F_{\rm GNWT}\Rightarrow G_{\rm IIT}F_{\rm IIT}$ is the pointwise inequality.
\begin{theorem}[Lax 2-cell]\label{thm:laxa}
\begin{itemize}
\item[(a)] A broadcast hub implies connectivity of the undirected graph of the support (\texttt{universal\_connected}, for a symmetric adjacency) for every $n\ge1$; by Lemma~\ref{lem:J} for $n\ge2$, and by Remark~\ref{rem:small} for $n=1$, it implies $J>0$.
\item[(b)] For $1\le n\le3$, broadcast and connectivity coincide: \texttt{n3\_equiv} for $n=3$, \texttt{n2\_equiv} for $n=2$, and $n=1$ is trivial.
\item[(c)] For every $n\ge4$, the path $P_n$ is connected (\texttt{path\_reach\_from\_zero}) and has no hub (\texttt{path\_no\_universal}): the 2-cell is strict for \emph{every} $n\ge4$, not only for $n\in\{4,5\}$. For $n=4$ there are exactly $15$ connected labeled graphs without a hub (\texttt{n4\_count}): the $12$ labelings of $P_4$ and the $3$ of $C_4$, exactly the two types of the numerical example, while the star $K_{1,3}$ has a hub. For $n=5$ there are $472$: they are the $728$ connected labeled graphs minus the $256$ with a hub, by inclusion--exclusion. This count is not in Lean; it is checked by exhaustive enumeration of the $2^{10}$ labeled graphs on five units in \texttt{pistas\_numerico.py}, which also checks the counts $0$, $0$ and $15$ for $n=2,3,4$.
\item[(d)] Quantitative version, for $n\ge2$: if a unit $w$ satisfies $\max(|A_{vw}|,|A_{wv}|)\ge c$ for every $v\ne w$ (a hub of strength $c$), then $J\ge\tfrac12\log(1+c^2/\kappa_E)$. If the coupling has strength $c$ in both directions, $|A_{vw}|,|A_{wv}|\ge c$, the bound is $J\ge\tfrac12\log(1+2c^2/\kappa_E)$.
\end{itemize}
\end{theorem}
\begin{proof}[Proof of (d)]
Every cut separates the hub $w$ from some unit $v$, so $A_{vw}$ and $A_{wv}$ cross it. Hence $w_{\min}\ge c^2$, or $w_{\min}\ge2c^2$ if both have strength $c$. Theorem~\ref{thm:cheeger} applies.
\end{proof}
\emph{Scope.} The 2-cell holds when both projections are computed on \emph{the same} causal model and at the same grain. It is not a theorem about the theories as such. IIT and GNWT can also disagree when the grain or the substrate changes: for example, a workspace architecture implemented on a conventional computer would satisfy functional broadcast while having, according to IIT, a negligible $\Phi$~\cite{TononiKoch2015}. Moreover, $J$ is not the $\Phi$ of IIT~3.0.

\subsection{Route 3: selection in the fibres}\label{sec:via3}
\begin{theorem}[\texttt{section\_factor\_iff}]\label{thm:section}
Let $\Phi:D\to V$ and $u:D\to U$ be functions, and let $s:V\to D$ be a section of $\Phi$ on its image, that is, $\Phi(s(\Phi(x)))=\Phi(x)$ for every $x\in D$. Then the values at the representatives define a factor of $u$, that is, $u=u\circ s\circ\Phi$, if and only if $S_u$ already holds: $\Phi(x)=\Phi(y)\Rightarrow u(x)=u(y)$.
\end{theorem}
With $\Phi=\Phi_Q$ on $D=D_0$, choosing representatives adds nothing to $S_u$: the other realizations in the fibre remain in $D_0$ and must have the same degree. The version with content is a \emph{physical selection law}.
\begin{definition}[SEL]\label{def:SEL}
Let $\mathcal F$ be a real-valued functional on $D_0$. The \emph{physical selection law} SEL for $\mathcal F$ states that every realization in $D_0$ belongs to the unique class, up to isomorphism, that minimizes $\mathcal F$ in its fibre: for every $s\in D_0$, the minimizers of $\mathcal F$ on the fibre $\{t\in D_0:\Phi_Q(t)=\Phi_Q(s)\}$ form a single isomorphism class at grain $\gamma$, and $s$ belongs to it.
\end{definition}
Under SEL, any two realizations in $D_0$ with the same value of $\Phi_Q$ belong to the minimizing class of their common fibre, so every fibre of $\Phi_Q$ on $D_0$ is a single isomorphism class (fibre purity). With OI$_\gamma$ (Definition~\ref{hyp:OI}), which makes $u$ constant on isomorphism classes, purity gives $S_u$ (\texttt{purity\_gives\_Su}). But SEL contradicts Rob as soon as $\Phi_Q$ is continuous there. Assume Rob (Definition~\ref{hyp:Rob}), with the conscious realization $s_0$, the convex set $P$ of kernel parameters and the realizations $s_\theta$ as there, and suppose that $\theta\mapsto\Phi_Q(s_\theta)$ is continuous on $P$ and that $\Phi_Q(s_0)$ is interior to its range. Then Proposition~\ref{prop:extension}(i), which rests on Theorem~\ref{thm:continuum}, gives uncountably many pairwise non-isomorphic realizations in the fibre through $s_0$, and a single non-isomorphic pair in one fibre already destroys purity (\texttt{witness\_breaks\_purity}). Without continuity there is no contradiction: an admissible descriptor that is discontinuous and injective on the classes makes all fibres pure.

Three points about FEP in this role:
\begin{itemize}
\item variational free energy is defined for each system relative to its own generative model, so comparing different systems requires a common reference that is not given;
\item the non-equilibrium steady state (NESS) is a state of a system, not a selection among systems;
\item the uniqueness of the minimizer is not guaranteed.
\end{itemize}

\subsection{Route 4: the connected domain}\label{sec:via4}
On $\mathcal C_{J>0}=\{C_4,K_{1,3},P_4\}$ the integration presence predicate equals $1$ everywhere, so every presence 2-cell into it is trivial (\texttt{connected\_domain\_trivial}). The non-commutation with broadcast persists on $C_4$ and $P_4$. Restricting the domain does not remove it: it removes the case in which both theories say 0. Only the quantitative 2-cell (Theorem~\ref{thm:laxa}(d)) provides information on $\mathcal C_{J>0}$.

\subsection{Weak unification theorem (structural)}\label{sec:unif}
\begin{theorem}[Structural unification in $\mathbf{Open}_G$]\label{thm:unif}
In the symmetric monoidal groupoid of open linear-Gaussian systems, with structural broadcast and blankets:
\begin{enumerate}
\item[(i)] \textbf{Extensive plus defect.} For systems with at least two units, $I_{\rm pred}=I_{\rm cut}+J$, and $J>0$ if and only if the support is connected. $I_{\rm pred}$ is strongly monoidal, and $J$ and broadcast vanish on the tensor ideal $\mathcal Z$.
\item[(ii)] \textbf{Sandwich.} For systems with at least two units, $J$ is comparable to $w_{\min}$ in the sense of Theorem~\ref{thm:cheeger}.
\item[(iii)] \textbf{Lax 2-cell.} A hub of strength $c$ gives $J\ge\tfrac12\log(1+c^2/\kappa_E)$, and $J\ge\tfrac12\log(1+2c^2/\kappa_E)$ if the coupling has that strength in both directions. The 2-cell is strict for every $n\ge4$ and is an equality of predicates for $1\le n\le3$.
\item[(iv)] \textbf{Retained obstructions.}
  \begin{itemize}
  \item There is no strict commutation of the presence bridges.
  \item There is no translation between positive integration and broadcast.
  \item There are blind spots.
  \item No selection in the fibres gives $S_u$ without changing the domain, and the law SEL (Definition~\ref{def:SEL}) contradicts Rob if $\theta\mapsto\Phi_Q(s_\theta)$ is continuous on the convex set $P$ of Rob and $\Phi_Q(s_0)$ is interior to its range (Section~\ref{sec:via3}). This point concerns the original domain $D_0$, not $\mathbf{Open}_G$.
  \end{itemize}
\end{enumerate}
\end{theorem}
\emph{Status.} (i)--(iii) are theorems about invariants of the model, on the physical side, with numerical checks. Their analytic statements about $J$, that $J>0$ if and only if the support is connected, so that $J$ vanishes exactly on $\mathcal Z$, and the sandwich, are proved on paper and checked in Lean with Mathlib (Lemma~\ref{lem:J}, Theorem~\ref{thm:cheeger}) \stA\ \stM. The decomposition $I_{\rm pred}=I_{\rm cut}+J$, the strong monoidality of $I_{\rm pred}$ and the hub bound in (iii) are proved on paper only \stA, and the combinatorial parts are checked in Lean without Mathlib \stL. (iv) collects the compiled counterexamples. They order the projections into a graded family with a common defect structure. They relate two structural proxies computed on the same model, not IIT and GNWT as theories (Section~\ref{sec:via2}, \emph{Scope}). They establish no bridge law, and $S_u$ remains open.


{\sloppy Lean without Mathlib checks only the logical and combinatorial skeleton of this block, on Boolean predicates and abstract types \stL. \texttt{FiltracionCociente.lean} contains \texttt{noncomm\_locus}, \texttt{commute\_off\_band}, \texttt{zmul\_assoc}, \texttt{strong\_monoidal\_zero}, \texttt{factors\_through\_collapse}, \texttt{restriction\_vacuous}, \texttt{seg\_is\_Su}, \texttt{seg\_implies\_Su}, \texttt{segd\_iff} and \texttt{pullback\_determined}. \texttt{PistasUnificacion.lean} contains \texttt{universal\_connected}, \texttt{n3\_equiv}, \texttt{path\_reach\_from\_zero}, \texttt{path\_no\_universal}, \texttt{n4\_count}, \texttt{tensor\_no\_universal}, \texttt{section\_factor\_iff}, \texttt{purity\_gives\_Su}, \texttt{witness\_breaks\_purity} and \texttt{connected\_domain\_trivial}; \texttt{n2\_equiv} is in \texttt{GaussConexion.lean}. \texttt{noncomm\_locus} and \texttt{commute\_off\_band} take as a hypothesis that the hub implies connectivity. \texttt{strong\_monoidal\_zero} assumes strict unit laws for the tensor and that every object other than the unit is nonempty, and, on the projection, the three conditions stated in Section~\ref{sec:l2}: the unit goes to the adjoined unit $e$, every nonempty object receives a value other than $e$, and every product of two nonempty objects has the value $0$. Lean with Mathlib checks the following analytic statements about $J$, for $n\ge2$ \stM: Lemma~\ref{lem:J} and the vanishing of $J$ exactly on disconnected supports (\texttt{QAudit.lemmaJ\_stationary}, \texttt{QAudit.J\_eq\_zero\_iff}); the bounds on $J$ (Theorem~\ref{thm:cheeger}, \texttt{QAudit.sandwich}); and the identification of the presence $J>0$ with $w_{\min}>0$ (\texttt{QAuditBridge.J\_pos\_iff\_wmin\_pos}), which, with $J\ge0$, gives the vanishing of $J$ on the class $\mathcal I$ defined in Section~\ref{sec:l2}. The following remain on paper \stA: the decomposition $I_{\rm pred}=I_{\rm cut}+J$, the additivity of $I_{\rm pred}$, the hub bound of Theorem~\ref{thm:laxa}(d) and what is said about the restricted Gaussian family. The conditional unification proposition is not formalized.\par}

\subsection{Line 1: filtration by connectivity}\label{sec:l1}
Let $\mathcal C_{\rm Central}$ be the systems with a broadcast hub and $\mathcal C_{w_{\min}>0}$ the nonempty systems with connected support, that is, with $w_{\min}>0$: the one-unit systems, for which $w_{\min}=+\infty$ (Remark~\ref{rem:small}), and the systems with at least two units and $w_{\min}>0$. On nonempty systems, $w_{\min}>0$ is equivalent to $J>0$, by Lemma~\ref{lem:J} and Remark~\ref{rem:small}; for $n\ge2$ this equivalence is checked in Lean with Mathlib (\texttt{QAuditBridge.J\_pos\_iff\_wmin\_pos}) \stM. There is a filtration by full subcategories
\begin{equation}
\mathcal C_{\rm Central}\ \subseteq\ \mathcal C_{w_{\min}>0}\ \subseteq\ \mathbf{Open}_G^{+},
\end{equation}
whose first inclusion is \texttt{universal\_connected} (for one unit it is trivial). The complement of $\mathcal C_{w_{\min}>0}$ in $\mathbf{Open}_G^{+}$ is the class $\mathcal I$ of Section~\ref{sec:l2}, the systems with at least two units and disconnected support; the empty system lies outside $\mathbf{Open}_G^{+}$ (Remark~\ref{rem:small}). For every $n\ge4$ the first inclusion is strict: the path $P_n$ is connected (\texttt{path\_reach\_from\_zero}) and has no hub (\texttt{path\_no\_universal}).
\begin{theorem}[Exact locus of non-commutation; \texttt{noncomm\_locus}, \texttt{commute\_off\_band}]\label{thm:locus}
On $\mathbf{Open}_G^{+}$, the presence bridges of IIT type ($J>0$) and GNWT type (hub) disagree exactly on the band $\mathcal C_{w_{\min}>0}\setminus\mathcal C_{\rm Central}$, and commute outside it.
\end{theorem}
\begin{proof}
Apply \texttt{noncomm\_locus} and \texttt{commute\_off\_band} to the Boolean predicates ``has a hub'' and ``$J>0$''; on nonempty systems the latter is connectivity of the support (Lemma~\ref{lem:J} and Remark~\ref{rem:small}). Their hypothesis, that a hub implies connectivity, is Theorem~\ref{thm:laxa}(a).
\end{proof}
For $n=4$ the band has $15$ labeled graphs: the paths $P_4$ and the cycles $C_4$. The lax 2-cell is the inclusion of the filtration. At the central level it is quantitative: a hub of strength $c$ gives $J\ge\tfrac12\log(1+c^2/\kappa_E)$, or $\tfrac12\log(1+2c^2/\kappa_E)$ with coupling in both directions (Theorem~\ref{thm:laxa}(d)). On the band, whose systems have at least four units, only the minimum-cut bound holds, $J\ge\tfrac12\log(1+w_{\min}/\kappa_E)>0$. No global equivalence is possible: $P_n$ is in the band for every $n\ge4$.

\subsection{Line 2: quotient by the factorizable ideal}\label{sec:l2}
$\mathcal I=\{X:\ X\text{ has at least two units and }w_{\min}(X)=0\}$ is the class of systems with at least two units and disconnected support; by Lemma~\ref{lem:J} it is the ideal $\mathcal Z$ of Section~\ref{sec:via1}. A one-unit system is not in $\mathcal I$, since $w_{\min}=+\infty$ there, and neither is the empty system, on which $w_{\min}$ is not defined (Remark~\ref{rem:small}). $\mathcal I$ contains the decomposable objects $X\otimes Y$ ($X,Y$ nonempty), but it does not coincide with them if $E$ is not block diagonal along the cut: with $A=\mathrm{diag}(0.5,\,0.4)$ and $E=\bigl(\begin{smallmatrix}1&0.6\\0.6&1\end{smallmatrix}\bigr)$ one has $w_{\min}=J=0$, and the stationary mutual information between the two units is $0.2188$ nats, so the system is not a tensor product. ``Factorizable'' here means factorizable in the dynamics $A$, not in the joint law. It is a tensor ideal: $X\in\mathcal I\Rightarrow X\otimes Y\in\mathcal I$, and $X\otimes Y\in\mathcal I$ for nonempty $X,Y$. The collapse $q:\mathbf{Open}_G\to\mathbf{Open}_G/\mathcal I$ sends $\mathcal I$ to a point. Since $\mathbf{Open}_G$ is a groupoid, the collapse is not an abelian-category quotient.
\begin{theorem}[\texttt{factors\_through\_collapse}, \texttt{strong\_monoidal\_zero}]\label{thm:collapse}
\begin{itemize}
\item[(a)] A projection factors through $q$ if and only if it is constant on $\mathcal I$. The projections of IIT type (the value of $J$ and the presence $J>0$) and broadcast are constant on $\mathcal I$, since they take the value $0$ on $\mathcal I$. $I_{\rm pred}$ is not.
\item[(b)] Precise replacement of monoidality: let $(D\sqcup\{e\},\circledast)$ have unit $e$ and $a\circledast b=0$ for $a,b\ne e$, where $0\in D$ is a fixed zero value; it is a monoid (\texttt{zmul\_assoc}). If one adopts the convention that the empty system goes to the unit $e$ and no nonempty system goes to $e$, then $J$, with $D=[0,+\infty]$, and broadcast, with $D=\{0,1\}$ for the presence of a hub or $D=\mathbb N$ for the number of hubs (zero $0$), are \emph{strong monoidal} functors into that zero-product monoid: $F(\varnothing)=e$ and $F(X\otimes Y)=F(X)\circledast F(Y)$ for all $X,Y$. Without that convention they are not: if the empty system went to a value of $D$, such as $0$, then for $X=K_{1,3}$ one would have $F(X\otimes\varnothing)=F(X)\ne0=F(X)\circledast F(\varnothing)$. This is a reformulation of the fact that they vanish on the products of nonempty systems, not an additional property. The set of hubs itself is not an isomorphism invariant, so by Proposition~\ref{prop:funtor} it is not a functor into this discrete monoid; it is transported along isomorphisms, as a functor into finite sets and bijections (Section~\ref{sec:tesis}). $I_{\rm pred}$ is strong monoidal into $(\mathbb R_{\ge0},+,0)$.
\end{itemize}
\end{theorem}
\begin{proof}
(a) is \texttt{factors\_through\_collapse}, with the point to which $\mathcal I$ collapses sent to the constant value; the values on $\mathcal I$ are given by Lemma~\ref{lem:J} for $J$ and by Theorem~\ref{thm:laxa}(a) for broadcast, since a system with a hub has connected support. Some systems of $\mathcal I$ have different values of $I_{\rm pred}$, for instance $A=\mathrm{diag}(a,a)$ and $E=I$ with different values of $a$. (b) Lean proves the abstract form, \texttt{strong\_monoidal\_zero}, for a tensor, a unit and a predicate ``nonempty'' on an abstract type of objects, under six hypotheses: strict left and right unit laws, $\mathrm{unit}\otimes x=x$ and $x\otimes\mathrm{unit}=x$ (\texttt{hunitL}, \texttt{hunitR}); every object other than the unit is nonempty (\texttt{hne}); the projection sends the unit to $e$ (\texttt{hF0}); every nonempty object receives a value in $D$, not $e$ (\texttt{hFne}); and every product of two nonempty objects has the value $0$ (\texttt{hvan}). The first three concern the monoidal structure: they hold on isomorphism classes of systems, on which an isomorphism-invariant projection is defined and on which the class of the empty system is the only one without units. \texttt{hF0} and \texttt{hFne} are the convention; for $J$, \texttt{hFne} holds because $J=+\infty$ on one-unit systems (Remark~\ref{rem:small}). \texttt{hvan} is Proposition~\ref{prop:monoidal}: $J(X\otimes Y)=0$, and $X\otimes Y$ has no hub (\texttt{tensor\_no\_universal}). The additivity of $I_{\rm pred}$ is Proposition~\ref{prop:monoidal}.
\end{proof}
\emph{Volume and obstruction:} for systems with at least two units, $I_{\rm pred}=I_{\rm cut}+J$, where $I_{\rm pred}$ is the extensive term, which survives in $X\otimes Y$, and $J$ is the obstruction to partition. The correct bound is $J\ge\tfrac12\log(1+w_{\min}/\kappa_E)$, not $J\ge w_{\min}$. Calling $I_{\rm pred}$ the ``FEP term'' remains an interpretation (Section~\ref{sec:via1}).

\subsection{Line 3: variational selection in the fibres}\label{sec:l3}
\begin{theorem}[\texttt{section\_factor\_iff}, \texttt{restriction\_vacuous}, \texttt{segd\_iff}]\label{thm:selection}
\begin{itemize}
\item[(a)] No section of $\Phi$, not even the section that chooses a free-energy minimizer, gives $S_u$: the values at the representatives define a factor if and only if $S_u$ already holds.
\item[(b)] Restricting the domain to the image of a section leaves each fibre with a single element, and $S_u$ becomes trivially true for \emph{every} $u$. The obligation is not met: it is eliminated by changing the domain, which this audit excludes by method: it does not change the descriptor, the domain, the components or the meaning of Q (Section~\ref{sec:intro}).
\item[(c)] Requiring the fibre to ``satisfy SEG'' in order to collapse it is circular. Under $B_\psi$, SEG (Definition~\ref{def:SEG}) holds if and only if $S_u$ holds with $u=\psi\circ\mathfrak U$ nonconstant on $D_0$ (\texttt{segd\_iff}: for the right-to-left direction, take the level sets of $\psi$ as the groups), and already a split into at least two groups with $\psi$ constant on each group and condition~(ii) of Definition~\ref{def:SEG} implies $S_u$ (\texttt{seg\_implies\_Su}). \texttt{seg\_is\_Su} only records that ``equal $\Phi$, equal $\psi\circ\mathfrak U$'' is, by definition, $S_u$ (\texttt{Iff.rfl}).
\end{itemize}
\end{theorem}
\emph{The restricted Gaussian family} of Section~\ref{sec:gauss} is an auxiliary test bench. Its descriptor is $J$, which is not $\Phi_Q$, and its obligation is $S_J$, which is not $S_u$. In this family, at the level of laws, each fibre of $J$ is a single class: the normalization $Z_i=X_i/\sqrt{s_i}$ gives identity noise and coupling $\sqrt k$, with $k=c^2s_2/s_1$, so there is nothing to select. The multiplicity lies in the pointed state, $z\sim N(0,\mathrm{diag}(1+k,1))$, and in the histories $z_{0:T}$. If the free energy of a state is taken to be its surprisal $-\log p_\theta(z)$, choosing the state of minimum free energy, the mode $z=0$, discards the realizations: they have $z\ne0$ almost surely, and $S_J$ requires that they too have the degree of the mode. Finally, a \emph{physical selection} law (only the minimizers are realized) would give fibre purity. In the family it contradicts the model itself, in which $z\ne0$ almost surely; in the original domain it is the law SEL of Definition~\ref{def:SEL}, which contradicts Rob when $\theta\mapsto\Phi_Q(s_\theta)$ is continuous on the convex set $P$ of Rob and $\Phi_Q(s_0)$ is interior to its range (Section~\ref{sec:via3}).

\subsection{Line 4: the weak limit}\label{sec:l4}
The pullback of a section $s$ along $\Phi$ is the subcategory of canonical representatives, $\{X: X=s(\Phi(X))\}$. On it, every additional function of the representative is a function of $\Phi$ (\texttt{pullback\_determined}): the weak limit adds no information to the section. The cut spectrum $w_{\min}$ is one of those functions, but not an invariant of $\mathbf{Open}_G$. Positive rescalings, which are isomorphisms, change it: on $C_4$, $w_{\min}=0.36$ with $D=I$ and $0.765$ with $D=\mathrm{diag}(1,2,1,2)$, with $J=0.2213$ in both cases. Only its vanishing is invariant, and the sandwich of Theorem~\ref{thm:cheeger} holds for each representative. Linking the weak limit to $S_u$ or to SEG therefore presupposes $S_u$: by part~(a) of Theorem~\ref{thm:selection}, the values at the representatives define a factor only if $S_u$ already holds.

The only thing that can be asserted here is conditional and, given $S_u$, immediate:
\begin{proposition}[Conditional unification]\label{prop:conditional}
Assume $B_\psi$. If SEG holds on $D_0$ (that is, $S_u$ with $u_{\rm orig}$ nonconstant; Definition~\ref{def:SEG}), then:
\begin{itemize}
\item $u_{\rm orig}=h\circ\Phi_Q$ with $h$ unique on $\Phi_Q(D_0)$;
\item every section computes $h$;
\item if each realization in $D_0$ is represented by a nonempty system of $\mathbf{Open}_G$ at grain $\gamma$, the filtration by cuts stratifies $D_0$ into the central, band and disconnected-support levels, on which the presence bridges commute (central and disconnected) or disagree (band). This point does not depend on SEG.
\end{itemize}
\end{proposition}
\begin{proof}
SEG gives $S_u$ (Definition~\ref{def:SEG} and \texttt{seg\_implies\_Su}), and $S_u$ is equivalent to the factorization $u_{\rm orig}=h\circ\Phi_Q$ with $h$ unique on $\Phi_Q(D_0)$ (Definition~\ref{def:phiQ}(e)). For a section $s$, \texttt{section\_factor\_iff} gives $u_{\rm orig}=u_{\rm orig}\circ s\circ\Phi_Q$, so $h(v)=u_{\rm orig}(s(v))$ for $v\in\Phi_Q(D_0)$. The third point is Theorem~\ref{thm:locus} applied to the representing systems.
\end{proof}
The premise is the open obligation, not its substitute. Nothing in this block contradicts the compiled counterexamples: no global equivalence is asserted, and $P_n$ always remains in the band.


\part*{Act II.\quad Validation in a Gaussian laboratory}
\section{The Gaussian laboratory}\label{sec:gauss}
\providecommand{\stL}{\textsf{[L]}}\providecommand{\stA}{\textsf{[A]}}\providecommand{\stM}{\textsf{[M]}}\providecommand{\stS}{\textsf{[S]}}\providecommand{\stN}{\textsf{[N]}}
This section defines the test bench of the paper: a two-unit linear-Gaussian family $F_0$ and a family with a driver with memory that contains it; Section~\ref{sec:networks} studies numerically the four-unit networks of Section~\ref{sec:base}. On the two families the quantities studied are obtained in closed form and checked exactly; on the networks the values are numerical. The results are then connected with the theorems of Act~I. Each claim carries its status:
\begin{itemize}
\item \stL: compiled in Lean without Mathlib (exit~0; axioms in Appendix~\ref{sec:archivos});
\item \stM: compiled in Lean with Mathlib, over real matrices (exit~0; only the standard axioms);
\item \stA: proved on paper;
\item \stS: verified exactly with SymPy~1.14 in the Gaussian model;
\item \stN: numerical, without a certificate.
\end{itemize}
What holds on the test bench does not thereby hold in the original domain $D_0$ of Q (Definition~\ref{def:phiQ}). The functional $J$ of Lemma~\ref{lem:J} is an auxiliary functional of the law of the process, not a selected $\Phi_Q$ and not the $\Phi$ of IIT~3.0. The sufficiency $S_J$ defined below is not $S_u$, and $J$ is nowhere identified with $\Phi_Q$.

\subsection{The test bench}
\begin{definition}[The restricted Gaussian family $F_0$]\label{def:gauss-F0}
For $\theta=(c,s_1,s_2)$ with $c\ge0$ and $s_1,s_2>0$, let $A_\theta=\bigl(\begin{smallmatrix}0&c\\0&0\end{smallmatrix}\bigr)$ and $E_\theta=\mathrm{diag}(s_1,s_2)$. The \emph{system} $\theta$ is the two-unit process
\[
X_{t+1}=A_\theta X_t+\varepsilon_{t+1},\qquad (\varepsilon_t)_t\ \text{independent, each}\ N(0,E_\theta),
\]
that is, $X_{1,t+1}=cX_{2,t}+\varepsilon_{1,t+1}$ and $X_{2,t+1}=\varepsilon_{2,t+1}$: unit~2 is the \emph{driver} and unit~1 the \emph{receiver}. Since $A_\theta^2=0$, the spectral radius of $A_\theta$ is $0$, so $(A_\theta,E_\theta)$ is an object of $\mathbf{Open}_G$ (Section~\ref{sec:base}), and the process has a unique stationary law $L_\theta$, under which it is always taken. $F_0$ is the set of these systems, and we write $k=k(\theta)=c^2s_2/s_1$.
\begin{enumerate}
\item[(a)] A \emph{pointed model} of $F_0$ is a pair $(\theta,x)$ with $x\in\mathbb R^2$, a \emph{pointed state}, or a pair $(\theta,x_{0:T})$ with $x_{0:T}=(x_0,\dots,x_T)\in(\mathbb R^2)^{T+1}$, a \emph{pointed history} of length $T$.
\item[(b)] An \emph{isomorphism} of pointed models is a change of coordinates $X\mapsto DX$, with $D=\mathrm{diag}(\alpha,\beta)$ and $\alpha,\beta\ne0$, that keeps the system in $F_0$. It sends $\theta$ to $(\alpha c/\beta,\ \alpha^2s_1,\ \beta^2s_2)$, which requires $\alpha c/\beta\ge0$, and each state $x_t$ to $Dx_t$. It is \emph{positive} if $\alpha,\beta>0$. These are the positive rescalings and sign changes of $\mathbf{Open}_G$ that fix each unit. (A permutation of the two units leaves $F_0$ unless $c=0$; for $c=0$ it would also identify $z$ with $(z_2,z_1)$ in (G3) below, and $|z|$ would remain invariant.)
\item[(c)] On $F_0$ the functional $J$ of Lemma~\ref{lem:J} has a single bipartition, $\{1\}|\{2\}$, whose predictors are the diagonal matrices $B=\mathrm{diag}(b_1,b_2)$:
\[
J(\theta)=\min_{b_1,b_2\in\mathbb R}\ \tfrac12\log\frac{\det\bigl(E_\theta+(A_\theta-B)P_\theta(A_\theta-B)^\top\bigr)}{\det E_\theta},
\]
where $P_\theta$ is the stationary covariance. $J$ is a function of the law, and it is evaluated on a pointed model through its system. Logarithms are natural, and values are in nats.
\item[(d)] For a function $u$ on a set of systems or pointed models, $S_J$ is the sufficiency of $J$ for $u$: $J(m)=J(m')\Rightarrow u(m)=u(m')$. It is the obligation $S_u$ of Definition~\ref{def:phiQ}(e), with $J$ in place of $\Phi_Q$ and the test bench in place of $D_0$.
\end{enumerate}
\end{definition}

\begin{definition}[The family with a driver with memory]\label{def:gauss-Fmem}
For $c\ge0$, $s_1,s_2>0$ and $a,a_{22}\in(-1,1)$, consider the process $X_{t+1}=AX_t+\varepsilon_{t+1}$ with
\[
A=\begin{pmatrix}a&c\\0&a_{22}\end{pmatrix},\qquad E=\mathrm{diag}(s_1,s_2),
\]
that is, $X_{1,t+1}=aX_{1,t}+cX_{2,t}+\varepsilon_{1,t+1}$ and $X_{2,t+1}=a_{22}X_{2,t}+\varepsilon_{2,t+1}$. The coefficient $a_{22}$ is the self-coupling of the driver, its memory, and $a$ is that of the receiver. The spectral radius of $A$ is $\max(|a|,|a_{22}|)<1$, and the process is taken under its stationary law. The \emph{family with a driver with memory}, $F_{\rm mem}$, consists of the systems with $a=0$, indexed by $(c,s_1,s_2,a_{22})$; $F_0$ is its subfamily with $a_{22}=0$. We keep $k=c^2s_2/s_1$ and put
\[
p=\frac1{1-a_{22}^2}\ \ge1,\qquad r(k,p)=1+\frac{kp(k+1)}{kp+1};
\]
$p$ is the stationary variance of the driver in units of its innovation variance, and every $p\ge1$ is realized, by $a_{22}=\sqrt{1-1/p}$. Pointed models, isomorphisms, $J$ and $S_J$ are as in Definition~\ref{def:gauss-F0}, for the systems of this family. Systems with $a\ne0$ are used only at the end of (G7).
\end{definition}

\subsection{Exact results on the test bench}
There were $126$ exact checks in four scripts for the items (G1)--(G7) below, and $30$ more for the connections (i)--(v) of the next subsection, $156$ in all, with no failures; the scripts are listed in Appendix~\ref{sec:archivos}.

\paragraph{(G1) Exact value.} Since $A_\theta^2=0$, the stationary process is $X_t=\varepsilon_t+A_\theta\varepsilon_{t-1}$, and $P_\theta=E_\theta+A_\theta E_\theta A_\theta^\top=\mathrm{diag}(s_1+c^2s_2,\ s_2)$ is the unique symmetric solution of the Lyapunov equation $P=A_\theta PA_\theta^\top+E_\theta$. With $B=\mathrm{diag}(b_1,b_2)$ and $P_{11}=s_1+c^2s_2$,
\begin{equation}
\det\bigl(E_\theta+(A_\theta-B)P_\theta(A_\theta-B)^\top\bigr)-s_2P_{11}=b_1^2s_2P_{11}\,(1+b_2^2)+b_2^2s_1s_2 ,
\end{equation}
a sum of nonnegative terms that vanishes only at $b_1=b_2=0$. Hence $J=\tfrac12\log(P_{11}/s_1)=\tfrac12\log(1+k)$, with a unique minimizer, $B=0$. The objective is not jointly convex: with $c=s_1=s_2=1$, the Hessian of the residual determinant has determinant $-16$ at $(b_1,b_2)=(1,1)$, and that of $\tfrac12\log$ of the ratio, $-4/343$. Uniqueness follows from the identity, not from convexity. For example, with $c=s_1=1$, $J=\tfrac12\log2$ if $s_2=1$ and $J=\log2$ if $s_2=3$; with $(c,s_1,s_2)=(1.3,\,0.7,\,1.9)$, $J=0.86023$ nats, that is, $1.24106$ bits. \stS

\paragraph{(G2) Scale.} Under $X_i\mapsto\alpha_iX_i$ with $\alpha_i>0$, the system $(c,s_1,s_2)$ becomes $(\alpha_1c/\alpha_2,\ \alpha_1^2s_1,\ \alpha_2^2s_2)$. Then $k$ and $J$ are invariant, and so is the whole objective of (G1), whereas $c^2s_2$ is only covariant: it is multiplied by $\alpha_1^2$. The normalization $Z_i=X_i/\sqrt{s_i}$ gives identity noise, coupling $\sqrt k$ and stationary covariance $\mathrm{diag}(1+k,1)$. The correspondence $J\leftrightarrow k$ is bijective, with $k=e^{2J}-1$. \stS

\paragraph{(G3) Pointed models.} For positive isomorphisms, the pair $(k,z)$, with $z=(x_1/\sqrt{s_1},\,x_2/\sqrt{s_2})$, is a complete invariant of pointed states, and the action is free. With the sign changes, $z$ is identified with $-z$ if $c>0$, and with $(\pm z_1,\pm z_2)$ if $c=0$; $|z|$ remains invariant. Each fibre $\{J=J_0\}$ contains a continuum of non-isomorphic pointed classes, indexed by $|z|$ and, for histories, by $z_{0:T}=(z_0,\dots,z_T)$. At the level of laws, by contrast, each fibre is a single class, because $k$ is a complete invariant of the systems. \stS, with the following written proof \stA.
\begin{proof}[Proof of (G3)]
A positive isomorphism $D=\mathrm{diag}(\alpha,\beta)$ sends $(c,s_1,s_2)$ to $(\alpha c/\beta,\alpha^2s_1,\beta^2s_2)$ and $x$ to $(\alpha x_1,\beta x_2)$, so $k$ and $z$ are unchanged. Conversely, if $(\theta,x)$ and $(\theta',x')$ have the same $k$ and the same $z$, then $D=\mathrm{diag}\bigl(\sqrt{s_1'/s_1},\sqrt{s_2'/s_2}\bigr)$ sends $s_1,s_2$ to $s_1',s_2'$, sends $c$ to the unique $c'\ge0$ with $c'^2s_2'/s_1'=k$, and sends $x=(\sqrt{s_1}z_1,\sqrt{s_2}z_2)$ to $x'$. If $D$ fixes $\theta$, then $\alpha^2s_1=s_1$ and $\beta^2s_2=s_2$ with $\alpha,\beta>0$, so $\alpha=\beta=1$: the action is free. A sign change keeps $\alpha c/\beta\ge0$; if $c>0$ this requires $\alpha\beta>0$, and the only extra identification is $z\mapsto-z$; if $c=0$ all signs are allowed. In every case $|z|$ is invariant. By (G2), $\{J=J_0\}=\{k=k_0\}$ with $k_0=e^{2J_0}-1$, and the pointed states $\bigl((\sqrt{k_0},1,1),(\rho,0)\bigr)$, $\rho\ge0$, lie in that fibre and have $|z|=\rho$, so they are pairwise non-isomorphic. For histories the same $D$ acts at every time, so $z_{0:T}$ is complete for positive isomorphisms. Finally, two systems with the same $k$ are related by the $D$ above, so $k$ is a complete invariant of the systems, and each fibre of $J$ consists of one class of laws.
\end{proof}

\paragraph{(G4) Countermodel.} Let $u=|z|^2=x_1^2/s_1+x_2^2/s_2$, or the mean of $|z_t|^2$ for histories. It is invariant and constant on each class, and it violates $S_J$. With $\theta_0=(1,1,1)$, the states $x=(0,0)$ and $x=(1,0)$ have the same law and give $u=0$ and $u=1$. The violation does not depend on null sets: under every law in the fibre, $z\sim N(0,\mathrm{diag}(1+k,1))$ by (G2), so $|z|^2$ has a continuous law and two independent draws differ almost surely. It is a formal function and does not attribute phenomenal unity to synthetic states. \stS

\paragraph{(G5) IPS$_K$ is $S_J$.} By (G3), invariance under isomorphisms only yields $u=F(k,z)$, not $u=H(k)$. Dependence on $k$ alone is a separate premise.
\begin{definition}[Premise IPS$_K$]\label{def:IPS}
Let $u$ be a function on a set $M$ of systems or pointed models of $F_0$ or of $F_{\rm mem}$, and let $k(m)$ be the value of $k=c^2s_2/s_1$ for the system of $m\in M$. The premise \emph{IPS$_K$} for $u$ is that $u$ depends only on $k$: $u=H\circ k$ on $M$ for some function $H$, or equivalently $k(m)=k(m')\Rightarrow u(m)=u(m')$ for all $m,m'\in M$.
\end{definition}
IPS$_K$ is a premise of the form ``the degree depends only on $\kappa$'' of Section~\ref{sec:intento}, with $\kappa=k$; it is examined here, not adopted. On $F_0$, $J=\tfrac12\log(1+k)$ and $k\mapsto\tfrac12\log(1+k)$ is injective (G2), so $J(m)=J(m')$ if and only if $k(m)=k(m')$, and IPS$_K$ is equivalent to $S_J$. In this narrow family the premise coincides with $S_J$. Being equivalent is not being circular, but neither does it supply an independent justification: justifying IPS$_K$ is justifying $S_J$. Finite check: on the five pointed models $(\theta_0,(0,0))$, $(\theta_0,(1,0))$, $(\theta_3,(0,0))$, $(\theta_3,(1,0))$ and $(\theta_1,(0,0))$, with $\theta_3=(\sqrt3,1,1)$ and $\theta_1=(2,1,\tfrac14)$, of which only the first and the last are isomorphic, there are $32$ assignments of values in $\{0,1\}$; $16$ of them are invariant; exactly the same $4$ satisfy IPS$_K$ and $S_J$; and $12$ are invariant and violate $S_J$. \stS

\paragraph{(G6) Finite windows.} The realized window of a system does not determine $J$.
\begin{proposition}[Finite-window non-identifiability]\label{prop:windows}
Let $T\ge0$, and for $\theta\in F_0$ let $\mu_\theta^T$ be the law of the window $(X_0,\dots,X_T)$ under the stationary law $L_\theta$.
\begin{enumerate}
\item[(i)] $\mu_\theta^T=N(0,\Sigma_T)$ on $\mathbb R^{2(T+1)}$, with $\Sigma_T=\mathcal L_T\,\mathrm{diag}(P_\theta,E_\theta,\dots,E_\theta)\,\mathcal L_T^\top$, where $\mathcal L_T$ is the block lower-triangular matrix with blocks $(\mathcal L_T)_{t0}=A_\theta^t$, $(\mathcal L_T)_{tj}=A_\theta^{t-j}$ for $1\le j\le t\le T$ and $0$ for $j>t$; and $\det\Sigma_T=(s_1+c^2s_2)\,s_2\,(s_1s_2)^T>0$.
\item[(ii)] For all $\theta,\theta'\in F_0$, the laws $\mu_\theta^T$ and $\mu_{\theta'}^T$ are mutually absolutely continuous.
\item[(iii)] If $k(\theta)\ne k(\theta')$, no measurable function $f$ on $\mathbb R^{2(T+1)}$ satisfies both $f=J(\theta)$ $\mu_\theta^T$-almost surely and $f=J(\theta')$ $\mu_{\theta'}^T$-almost surely.
\end{enumerate}
\end{proposition}
\begin{proof}
\stA\ (i) Iterating the recursion gives $X_t=A_\theta^tX_0+\sum_{j=1}^tA_\theta^{t-j}\varepsilon_j$, so $(X_0,\dots,X_T)=\mathcal L_T\,(X_0,\varepsilon_1,\dots,\varepsilon_T)$. Under $L_\theta$, $X_0\sim N(0,P_\theta)$ by (G1), and $X_0$, a function of the innovations up to time $0$, is independent of $\varepsilon_1,\dots,\varepsilon_T$, which are independent $N(0,E_\theta)$. Hence the window is centred Gaussian with covariance $\Sigma_T$. The diagonal blocks of $\mathcal L_T$ are identities, so $\det\mathcal L_T=1$ and $\det\Sigma_T=\det P_\theta\,(\det E_\theta)^T$; since $P_\theta,E_\theta\succ0$ and $\mathcal L_T$ is invertible, $\Sigma_T\succ0$. (ii) A Gaussian law with positive definite covariance has an everywhere positive density with respect to Lebesgue measure, so its null sets are exactly the Lebesgue-null sets, and two such laws have the same null sets. (iii) By (G2), $k(\theta)\ne k(\theta')$ gives $J(\theta)\ne J(\theta')$. If $f$ were such a function, the set $\{f\ne J(\theta)\}$ would be $\mu_\theta^T$-null, hence $\mu_{\theta'}^T$-null by (ii). Then the disjoint events $\{f=J(\theta)\}$ and $\{f=J(\theta')\}$ would both have $\mu_{\theta'}^T$-probability $1$, which is impossible.
\end{proof}
The factorization of $\Sigma_T$, the equality $\det\mathcal L_T=1$ and the value of $\det\Sigma_T$ were checked exactly for $T=0,\dots,5$ \stS; the proof above is general in $T$. For example, $\theta=(1,1,1)$ and $\theta'=(1,1,3)$ have $J=\tfrac12\log2$ and $J=\log2$ (G1). Only the innovation variance of the driver, which feeds the receiver, differs; the coupling $c$ and the receiver's own noise $s_1$ are the same, and no measurable function of any finite window returns $J$ almost surely under both laws. In the terms of Definition~\ref{def:adm}(c), $J$ is a law-level descriptor that no finite realized window determines. The proposition holds only for finite windows. On infinite histories $k$ is recovered almost surely by the strong law of large numbers, since the pairs $(X_{1,t+1},X_{2,t})=(c\varepsilon_{2,t}+\varepsilon_{1,t+1},\ \varepsilon_{2,t})$ are independent and identically distributed in $t$; this standard argument is not carried out here.

\paragraph{(G7) Driver with memory.} On $F_{\rm mem}$ (Definition~\ref{def:gauss-Fmem}) the stationary covariance has $P_{22}=s_2p$, $P_{12}=ca_{22}s_2p$, $P_{11}=s_1+c^2s_2p$ and $\det P=s_2p\,(s_1+c^2s_2)$. With $B=\mathrm{diag}(b_1,b_2)$ and $b_1^*=cP_{12}/P_{11}$,
\begin{equation}
\det\bigl(E+(A-B)P(A-B)^\top\bigr)=s_1s_2\,r(k,p)+s_2P_{11}\,(b_1-b_1^*)^2+(a_{22}-b_2)^2\bigl(s_1P_{22}+b_1^2\det P\bigr).
\end{equation}
Hence $J=\tfrac12\log r(k,p)$, with a unique minimizer at $B^*=\mathrm{diag}(b_1^*,a_{22})$. Moreover $\partial r/\partial p=k(k+1)/(kp+1)^2>0$, $r(k,1)=1+k$ and $2+k-r(k,p)=(k+1)/(kp+1)>0$, so at fixed $k>0$, $r(k,\cdot)$ has exact range $[1+k,2+k)$ on $p\ge1$. For $0<k\le k'<k+1$, the value $p_*=k'/(k(k+1-k'))$ satisfies $p_*\ge1$ and $r(k,p_*)=1+k'=r(k',1)$. \stS
\begin{proposition}[Conditional constancy]\label{prop:gauss-chains}
Let $u$ be a function on the systems of $F_{\rm mem}$ with $c>0$. If $u$ satisfies IPS$_K$ (Definition~\ref{def:IPS}) and $S_J$, then $u$ is constant.
\end{proposition}
\begin{proof}
\stA\ By IPS$_K$, $u=H\circ k$. Let $0<k\le k'<k+1$. Take a system $\theta$ with $k(\theta)=k$ and $p=p_*$, and a system $\theta'$ with $k(\theta')=k'$ and $a_{22}=0$; both exist, because $c>0$, $s_1$, $s_2$ and $a_{22}=\sqrt{1-1/p_*}$ can be chosen freely. By (G7), $J(\theta)=\tfrac12\log r(k,p_*)=\tfrac12\log r(k',1)=J(\theta')$, so $S_J$ gives $H(k)=u(\theta)=u(\theta')=H(k')$. Given $0<k<K$, put $n=\lfloor K-k\rfloor+1$ and $k_i=k+i(K-k)/n$ for $i=0,\dots,n$; consecutive values differ by $(K-k)/n<1$, so $H(k)=H(k_1)=\dots=H(K)$. Hence $H$ is constant on $(0,\infty)$, and so is $u$.
\end{proof}
Writing $(k,p)$ for a system with these values of $k$ and $p$, each link uses the premise once, between $(k,1)$ and $(k,p_*)$, and $S_J$ once, between $(k,p_*)$ and $(k',1)$: finite chains of steps smaller than one connect every pair of positive values of $k$, and four such chains were checked exactly \stS. In words: if the whole positive-coupling family lay in the domain of the degree, with $u$ attributed to its members, and the premise IPS$_K$ made $u$ depend only on $k$, then $S_J$ would force $u$ to be constant on that family. It is a conditional result, not a refutation of Q: neither Q nor this paper supplies a requirement that the degree vary on that family (A9--U allows plateaus, Definition~\ref{def:orderB}). It warns against introducing a bridge premise as if it were neutral: together with $S_J$, IPS$_K$ forces the constancy, which $S_J$ alone does not force ($u=J$ satisfies $S_J$ and is not constant).

IPS$_K$ also discards the internal coefficients. With a white driver ($a_{22}=0$), $J=\tfrac12\log(1+k)$ for every receiver coefficient $a\in(-1,1)$. If $a\ne0$ and $a_{22}\ne0$, $J$ depends on $a$ at fixed $k$: with $c=s_1=s_2=1$ and $a_{22}=\tfrac12$, the determinant ratio $e^{2J}$ is $15/7$ with $a=0$ and $187/87$ with $a=\tfrac12$. IPS$_K$ makes $u$ ignore $a$ even though $J$ does not ignore it, and that content is added by the premise. \stS\ The script of (G7) also checks the reading of the premise with the stationary variance of the driver in place of its innovation variance, $\tilde k=c^2P_{22}/s_1=kp$: at fixed $\tilde k$, $r$ decreases strictly in $p$, with range $(1+\tilde k/(\tilde k+1),\,1+\tilde k]$ on $p\ge1$ \stS.

\subsection{Connections with the theorems of Act~I}
\paragraph{(i) Commutation on $\{K_{1,3},2K_2\}$.} The undirected support of $F_0$ is $K_2=K_{1,1}$, the one-leaf star, if $c>0$, and two isolated units, $2K_1$, if $c=0$. With the undirected reading of broadcast, which is the one that Lean formalizes, $F_0$ falls entirely in the central and disconnected-support levels of the filtration (Section~\ref{sec:l1}), and the bridge square commutes on $F_0$: $c>0$ gives $G_{\rm IIT}=G_{\rm GNWT}=1$, and $c=0$ gives $0$ for both. This is covered by \texttt{n2\_equiv} \stL, which equates, for two units, connectivity and having a universal vertex, and by $J>0\iff c>0$ \stS. It is the two-unit analog of \texttt{square\_commutes\_on\_sub}, which deals with the four-unit star and split system.

The tensor product of two members of $F_0$ with $c,c'>0$ has support $2K_2$, the split system of Table~\ref{tab:C-1}. Its $J$ is $0$: the predictor $B=A$ does not cross the cut $\{1,2\}|\{3,4\}$ and leaves the residual $E$ \stS. Nor does it have a hub (\texttt{tensor\_no\_universal}) \stL. Thus $F_0$ and $F_0\otimes F_0$ realize, with closed-form values, the two types of system on which the square commutes.

\emph{A clarification on broadcast.} In the systems of Table~\ref{tab:C-1} the couplings are symmetric, and the undirected reading coincides with the demanding one, which requires coupling in both directions. On $F_0$ with $c>0$ they differ: the undirected support has a universal vertex, but no unit is coupled in both directions (\texttt{f0\_hub\_readings}) \stL. With the demanding reading, $F_0$ with $c>0$ would already be in the band at $n=2$, and the square would fail on all of $F_0$ with $c>0$. With the directed sender reading, by contrast, the driver $X_2$ influences the other unit and is a hub, so $F_0$ is not in the band. The statements of this paper use the undirected reading (Section~\ref{sec:proyecciones}).

\paragraph{(ii) The band $P_n$, $n\ge4$, and the sandwich.} $F_0$ cannot witness the discrepancy. For $n=2$ and $n=3$, connectivity and having a universal vertex coincide (\texttt{n2\_equiv}, \texttt{n3\_equiv}) \stL; for $n=1$ it is trivial. So the band begins at $n=4$. It is occupied by $P_n$ for every $n\ge4$ (\texttt{path\_reach\_from\_zero}, \texttt{path\_no\_universal}) \stL, with $J>0$ by Lemma~\ref{lem:J} \stA\ \stM.

What $F_0$ does provide is an exact check of the min-cut sandwich (Theorem~\ref{thm:cheeger}), which is proved on paper and checked in Lean with Mathlib \stA\ \stM. $w_{\min}$ and the two bounds depend on the representative, not only on the class: $J$ is invariant and $w_{\min}$ is not. In the representative $\theta=(c,s_1,s_2)$, $w_{\min}=c^2$ and, for $c>0$,
\begin{equation}
\tfrac12\log\bigl(1+c^2/\kappa_E\bigr)\ \le\ J=\tfrac12\log(1+k)\ <\ \tfrac12\,\frac{\lambda_{\max}(P)}{\lambda_{\min}(E)}\,c^2 .
\end{equation}
For $c>0$, the lower bound is attained if and only if $s_2\le s_1$ (for $c=0$ both sides vanish). If $s_2>s_1$, the slack is $k-c^2/\kappa_E=c^2(s_2^2-s_1^2)/(s_1s_2)$ \stS. On the member $\theta=(\tfrac12,1,1)$ the exact value $J=\tfrac12\log\tfrac54$, and with it the attainment of the lower bound, is also checked in Lean with Mathlib (\texttt{QAudit.Audit.J\_ex\_eq}) \stM. Since equality holds on the whole subfamily $s_1/s_2=\kappa_E\ge1$, for each value of $\kappa_E$, the constant of the lower bound cannot be improved in general. In the normalized representative ($E=I$, coupling $\sqrt k$) the bound is exact for every $\theta$. The values of $J$ on $P_4$ and $C_4$ remain numerical \stN: there is no verified closed form in the band.

\paragraph{(iii) Predictive additivity versus the zero-product monoid (Line~2).} On $F_0$,
\begin{equation}
I_{\rm pred}=\tfrac12\log(\det P/\det E)=\tfrac12\log(1+k)=J,
\end{equation}
so $I_{\rm cut}=0$. The best disconnected predictor is $B=0$, which captures nothing, because $E+APA^\top=P$: all the predictive information is integration. On $F_{\rm mem}$, with memory in the driver, $I_{\rm pred}=\tfrac12\log\bigl(p(1+k)\bigr)$ and
\begin{equation}
I_{\rm cut}=\tfrac12\log\frac{p(1+k)}{r(k,p)},\qquad \frac{p(1+k)}{r(k,p)}-1=\frac{(p-1)\bigl(1+kp(1+k)\bigr)}{1+kp(k+2)} ,
\end{equation}
which vanishes if and only if $p=1$. The separable part is the driver's own memory \stS.

For $X,Y\in F_0$ with $c,c'>0$,
\begin{equation}
I_{\rm pred}(X\otimes Y)=J(X)+J(Y)>0,\qquad J(X\otimes Y)=0=J(X)\circledast J(Y),
\end{equation}
and moreover $I_{\rm cut}(X\otimes Y)-I_{\rm cut}(X)-I_{\rm cut}(Y)=J(X)+J(Y)$ \stS. This is an exact instance of the decomposition of Section~\ref{sec:via1}:
\begin{itemize}
\item $I_{\rm pred}$ is strong monoidal into $(\mathbb R_{\ge0},+,0)$;
\item $J$ is strong monoidal into the zero-product monoid, whose operation and whose abstract criterion are in Lean (\texttt{zmul\_assoc}, \texttt{strong\_monoidal\_zero}) \stL;
\item the cut defect is all the information of $X\otimes Y$.
\end{itemize}
Lean proves the abstract form. Besides strict unit laws for the tensor and the requirement that only the unit be empty, which hold at least on isomorphism classes, the hypotheses of \texttt{strong\_monoidal\_zero} on the projection are three: the empty system goes to the adjoined unit, every nonempty object receives a value (with a single unit, $J=+\infty$ by convention) and every product of two nonempty objects has value $0$. That $J$ satisfies the third, $J(X\otimes Y)=0$, is Proposition~\ref{prop:monoidal} \stA. It also follows from \texttt{QAudit.J\_eq\_zero\_iff} \stM, because a product of two nonempty systems has disconnected support; that step is proved on paper and not formalized with Mathlib, and Lean without Mathlib checks only the analogous fact for hubs, that a product has none (\texttt{tensor\_no\_universal}) \stL. It is checked here on $F_0\otimes F_0$ \stS.

\paragraph{(iv) Triviality or discarding in free-energy selection (Line~3).} There are three levels, and each one gives an answer.
\begin{itemize}
\item \emph{Law.} The stationary law of each $\theta$ is unique, and each fibre of $J$ is a single class of laws (G3). Every invariant $u$ that depends only on the law is a function of $k$, hence of $J$. At that level $S_J$ is trivial and there is nothing to select. But the degree $u_{\rm orig}$ of Q is assigned to realizations, each with its realized history over a window (Definitions~\ref{def:pcm}(a) and~\ref{def:unity}): a reference to parameters of the law must be compatible with that episode-level assignment, and the realized history and the law of the process are different objects (history-level and law-level descriptors, Definition~\ref{def:adm}(c)). Triviality at the level of laws does not reach $u_{\rm orig}$.
\item \emph{Pointed state.} Take the free energy of a state to be its surprisal $-\log p_\theta(z)$, which is the bound that variational free energy attains when the recognition density is exact. Its unique minimizer is the mode, $z=0$ \stS. That set is null under every law of the family, so selecting it discards all realizations almost surely. On isomorphism classes, or for every invariant $u$, each fibre of $J$ restricted to $\{z=0\}$ has a single element and $S_J$ holds for every $u$ (\texttt{restriction\_vacuous}) \stL, but at the cost of changing the domain. On the representatives $(\theta,0)$, which form a continuum in each fibre, a non-invariant $u$ can violate it: with $\theta_0=(1,1,1)$ and $\theta_1=(2,1,\tfrac14)$, both with $k=1$, $u=s_2$ takes the values $1$ and $\tfrac14$. The values at the section give the factor $h\equiv0$, while $u=|z|^2>0$ almost surely. By \texttt{section\_factor\_iff} \stL, the section does not give a factor, because $S_J$ fails (G4).
\item \emph{History.} The history of minimum surprisal, $x_{0:T}=0$, is null under every law, because $\Sigma_T\succ0$ (Proposition~\ref{prop:windows}): selecting it also discards all realizations almost surely.
\end{itemize}
At all three levels, free-energy selection is trivial or discarding, as Line~3 states (Section~\ref{sec:l3}).

\emph{A separate, conditional result.} A degree that supervenes on the finite window meets a law-level descriptor as follows.
\begin{proposition}[Plateau]\label{prop:gauss-plateau}
Let $\Lambda$ be a set of laws of a process, let $\mathcal H$ be the set of its windows (histories over a window of fixed finite length), write $\operatorname{supp}L\subseteq\mathcal H$ for the support of the law of the window under $L\in\Lambda$, and let $D=\{(h,L)\in\mathcal H\times\Lambda:\ h\in\operatorname{supp}L\}$, so that every history admissible under a law is realized with that law. Let $\Phi(h,L)=F(L)$ be a law-level descriptor and $u(h,L)=U(h)$ a degree that depends only on the history. Then $S$: $\Phi(a)=\Phi(b)\Rightarrow u(a)=u(b)$ for $a,b\in D$ holds if and only if, for every value $v$ of $F$, $U$ is constant on $H_v=\bigcup_{F(L)=v}\operatorname{supp}L$. In particular, if all the laws of $\Lambda$ have the same support, $S$ holds only if $u$ is constant on $D$.
\end{proposition}
\begin{proof}
\stA\ Two realizations $(h,L)$ and $(h',L')$ with $F(L)=F(L')=v$ have the same $\Phi$, so $S$ requires $U(h)=U(h')$ for all $h,h'\in H_v$. Conversely, if $U$ is constant on each $H_v$, two realizations with $\Phi=v$ have their histories in $H_v$ and hence the same $u$. If all the supports equal one set $H$, then $H_v=H$ for every value $v$, and every realization has its history in $H$.
\end{proof}
On $F_0$, with windows of length $T$, all the laws $\mu_\theta^T$ have the full support $\mathbb R^{2(T+1)}$ (Proposition~\ref{prop:windows}), so $D$ is the product $\mathbb R^{2(T+1)}\times F_0$ of histories and laws. If $u$ supervenes on the finite window, as LSS$_\Delta$ (Definition~\ref{hyp:LSS}) would give if the window were the intrinsic history of a bearer, and all laws are realized, then $S_J$ forces $u$ to be constant. \texttt{plateau\_product} \stL\ formalizes only that consequence, on the product domain and in the direction ``$S$ implies constancy'', not the biconditional of the proposition: two co-occurrence steps, first same law and then same window. It is a case of the co-occurrence dichotomy of Theorem~\ref{thm:dicotomia}, with $\kappa$ the window. It is not a universal dichotomy between history-level descriptors and plateaus: without a common support, $S$ links two histories only when both lie in one set $H_v$, and $U$ must be constant on each component of the equivalence generated by these links but may differ between components. If, for example, the supports are pairwise disjoint and $F$ is injective, the history determines the law, and $S$ requires $U$ to be constant on each support but imposes nothing across supports, so a law-level descriptor coexists with $S$ and with a degree that is not globally constant \stA. That the support is full is Proposition~\ref{prop:windows}: $\Sigma_T\succ0$, with the written proof \stA\ and the exact check for $T\le5$ \stS.

\paragraph{(v) The constancy proofs as instances of the Lean theorems.} The two constancy results of this section are instances of \texttt{connected\_forces\_constant} \stL\ (Theorem~\ref{thm:dicotomia}).
\begin{itemize}
\item In the chains of Proposition~\ref{prop:gauss-chains}, $\kappa=k$ and the descriptor is $J$ on $F_{\rm mem}$. Co-occurrence links $k$ with $k'$ whenever $k\le k'<k+1$.
\item In the plateau of Proposition~\ref{prop:gauss-plateau}, $\kappa$ is the finite window, and the common support makes co-occurrence connected.
\end{itemize}
IPS$_K$, aligned with $J$ by a bijection, is $S_J$ reformulated (\texttt{aligned\_premise\_iff}) \stL, just as SEG (Definition~\ref{def:SEG}) is non-degenerate $S_u$ reformulated (\texttt{segd\_iff}) \stL. Neither of these premises is an independent justification, and neither of the constancy results refutes Q. Table~\ref{tab:E-1} summarizes the laboratory.

\begin{table}[htbp]\centering\small
\begin{tabular}{>{\raggedright\arraybackslash}p{0.38\textwidth}>{\raggedright\arraybackslash}p{0.38\textwidth}>{\raggedright\arraybackslash}p{0.14\textwidth}}
\toprule
Fact in $F_0$ & Theorem it instantiates & Status\\
\midrule
$J=\tfrac12\log(1+k)$, with a unique minimizer & Lemma~\ref{lem:J}: $J>0$ if and only if the support is connected & \stS\\
Support $K_{1,1}$ or $2K_1$; $F_0\otimes F_0$ has support $2K_2$ & commutation on $\{K_{1,3},2K_2\}$ (two-unit analog) & \stL\ and \stS\\
No band with $n\le3$ (undirected reading of the hub) & band $P_n$, $n\ge4$ & \stL\\
Lower bound attained if $s_2\le s_1$ & sandwich by $w_{\min}$ & \stA, \stM\ and \stS\\
$I_{\rm pred}=J$, $I_{\rm cut}=0$; additivity and $J(X\otimes Y)=0$ & Line~2: zero-product monoid & \stS\ and \stL\\
Fibres of laws consist of one class; the mode and the zero history have measure zero & Line~3: trivial selection or discarding & \stS\ and \stL\\
IPS$_K$ is equivalent to $S_J$ & entailment criterion; SEG and non-degenerate $S_u$ & \stS\ and \stL\\
Chains and Plateau & co-occurrence dichotomy & \stS, \stA\ and \stL\\
\bottomrule
\end{tabular}
\caption{Summary of the Gaussian laboratory: each fact in $F_0$, the theorem it instantiates, and its status.}\label{tab:E-1}
\end{table}

\paragraph{Limits.} Three things are not checked here: closed forms of $J$ for families with three or more units, so that none is verified in the band (Section~\ref{sec:networks}); readings of the premise through dependence measures other than $k$ and $\tilde k$, such as a transfer entropy; and the recovery of $k$ from infinite histories mentioned after Proposition~\ref{prop:windows}.

\subsection{Four-unit coupled networks}\label{sec:networks}
The two-unit family cannot exhibit the band (connection (ii) above). The four-unit networks of Section~\ref{sec:base}, used in Sections~\ref{sec:categorica} and~\ref{sec:debil}, can, and they test the bounds of Act~I numerically. For a graph $\Gamma\in\{C_4,K_{1,3},P_4,2K_2\}$ on four units, with adjacency matrix $M_\Gamma$, the network is the system of $\mathbf{Open}_G$ with
\[
X_{t+1}=AX_t+\varepsilon_{t+1},\qquad A=0.2\,I+0.3\,M_\Gamma,\qquad E=I:
\]
self-coupling $a=0.2$, symmetric couplings $c=0.3$ on the edges of $\Gamma$, and identity noise. The eigenvalues of $A$ are $0.2+0.3\lambda$ with $\lambda$ an eigenvalue of $M_\Gamma$; for these four graphs $|\lambda|\le2$, so the spectral radius of $A$ is at most $0.8<1$. $J$ is the functional of Lemma~\ref{lem:J}: for each of the seven bipartitions of the four units, the objective is minimized over the predictors that do not cross it by local optimization (BFGS) with restarts, and the least value found is reported. Table~\ref{tab:networks} reports $J$, the two sides of the min-cut sandwich (Theorem~\ref{thm:cheeger}), and the decomposition $I_{\rm pred}=I_{\rm cut}+J$ of Section~\ref{sec:via1}.
\begin{table}[htbp]\centering\small
\begin{tabular}{lccccc}
\toprule
network & $J$ & lower bound & upper bound & $I_{\rm pred}$ & $I_{\rm cut}$\\
\midrule
ring $C_4$ & $0.221293$ & $0.153742$ & $0.500000$ & $0.638824$ & $0.417531$\\
star $K_{1,3}$ & $0.113499$ & $0.082757$ & $0.186662$ & $0.459446$ & $0.345948$\\
path $P_4$ & $0.105871$ & $0.082757$ & $0.169743$ & $0.440231$ & $0.334360$\\
split $2K_2$ & $0$ & $0$ & $0$ & $0.297732$ & $0.297732$\\
\bottomrule
\end{tabular}
\caption{Four-unit networks ($a=0.2$, $c=0.3$, $E=I$): the integration functional $J$, the lower and upper bounds of the min-cut sandwich, and the predictive and cut information, in nats. Values from local optimization with restarts (\texttt{pistas\_numerico.py}); a second, independent optimizer, which is not archived, reproduced $J$ to seven digits. \stN}
\label{tab:networks}
\end{table}

What the table shows, and with which status:
\begin{itemize}
\item The sandwich holds on all four networks, and on twelve random four-unit systems with random noise covariance, eleven of them with connected support (in the remaining one, the first drawn, the random mask removed every off-diagonal coupling, so $J=0$ and both bounds vanish) (sparse Gaussian couplings and $E=GG^\top+\tfrac12I$ with Gaussian $G$, drawn from a fixed seed) \stN; it is proved on paper and checked in Lean with Mathlib (Theorem~\ref{thm:cheeger}; \texttt{QAudit.sandwich}) \stA\ \stM.
\item The star has a broadcast hub of strength $c$ in both directions, so $J\ge\tfrac12\log(1+2c^2/\kappa_E)$ (Theorem~\ref{thm:laxa}(d)). For the star with $E=I$ and $c\in\{0.1,0.2,0.3,0.4\}$, $J=0.010603$, $0.044930$, $0.113499$ and $0.252894$, above the bounds $0.009901$, $0.038481$, $0.082757$ and $0.138816$ \stN.
\item The split system $2K_2$ is a tensor product: $J=0$ while $I_{\rm pred}$ is additive, as in the exact two-unit computation of connection (iii) above. In the tensor product of a two-unit system ($\Gamma=K_2$, $c=0.4$) and a three-unit path ($\Gamma=P_3$, $c=0.3$), built in the same way with $a=0.2$ and $E=I$, $J(X\otimes Y)=0$ while $J(X)=0.181508$ and $J(Y)=0.104778$, and $I_{\rm pred}(X\otimes Y)=I_{\rm pred}(X)+I_{\rm pred}(Y)$ to twelve digits \stN.
\item The ring and the path are in the band: connected, without a hub, with $J>0$. This is where the presence bridges of IIT and GNWT type disagree (\texttt{square\_fails}, \texttt{values}) \stL, while the star and the split system are where they agree (\texttt{square\_commutes\_on\_sub}) \stL.
\end{itemize}


\part*{Act III.\quad A positive route and a methodological guide}
\section{A positive route: intervention monotonicity}\label{sec:via-futura}
Theorem~\ref{thm:dicotomia} closes one door. If equality of degree is required on every pair of $\ker\Phi_Q$, and the defended feature $\kappa$ co-occurs connectedly with the levels of $\Phi_Q$, the degree is constant. This section asks which reformulation of $S_u$ escapes that collapse without circularity. Two natural candidates do not escape (Section~\ref{sec:collapse}). One does (Section~\ref{sec:im}), at a price that we state exactly; we have not shown that it is the only reformulation that does. The surviving route can be instantiated in the band of Act~I through a monotonicity theorem for $J$, proved on paper and checked in Lean with Mathlib (Section~\ref{sec:scaling}). We close with a methodological guide and with the statements that should be formalized next. The Lean statements about degrees cited in this section are in \texttt{ViaFutura.lean} (the graph lemmas cited in Corollary~\ref{cor:band-scaling} are in \texttt{PistasUnificacion.lean}); they take degrees with integer values, and their arguments use only the structure of a totally ordered abelian group, so the statements carry over verbatim to real values, with the same arguments; that transcription is not formalized, and the Lean proofs, which use integer tactics such as \texttt{omega} and \texttt{decide}, would have to be rewritten. The statements of Proposition~\ref{prop:scaling} are checked with Mathlib, over real matrices, in \texttt{Scaling.lean}, \texttt{Interval.lean} and \texttt{Bridge.lean}.

\subsection{Two reformulations that collapse}\label{sec:collapse}
\paragraph{Exact covariance under perturbations $\Delta\Phi_Q$.} Let $\mathcal T$ be a class of perturbation steps, and suppose that along every step $u$ changes exactly as a function of $\Phi_Q$: $u(t)-u(s)=H(\Phi_Q(t))-H(\Phi_Q(s))$ for $(s,t)\in\mathcal T$. Then $u-H\circ\Phi_Q$ is constant along chains of steps, so the two ends of any directed chain with equal $\Phi_Q$ have equal $u$ (\texttt{exact\_covariance\_gives\_Su}) \stL. Exact covariance is symmetric, so the same holds along the symmetric closure of $\mathcal T$, and $S_u$ holds on every connected component of the undirected step graph (\texttt{exact\_covariance\_components}) \stL. Exact covariance is therefore $S_u$ on components, not an escape from it.

\paragraph{The order clause B on all pairs.} Suppose $u=f\circ\kappa$ and B holds on all pairs: $\Phi_Q(s)<\Phi_Q(t)$ implies $f(\kappa(s))\le f(\kappa(t))$. If a class $k$ occurs strictly below a class $k'$ and $k'$ occurs strictly below $k$, then $f(k)=f(k')$ (\texttt{order\_all\_pairs\_general}) \stL. With a transversal $\kappa$ whose image is a product of at least two distinct levels of $\Phi_Q$ and the realized classes, every pair of realized classes is in this situation, and $f$ is constant on them; \texttt{order\_all\_pairs\_collapse} \stL\ is the instance with two levels and two classes.

\paragraph{A cohomological obstruction $H^1$ to gluing bands of SEG is ill-posed.} Consider the presheaf of local factorizations, $\mathcal F(U)=\{h:U\cap\operatorname{im}\Phi_Q\to\mathbb R \mid u=h\circ\Phi_Q \text{ on } \Phi_Q^{-1}(U)\}$. It is a sheaf, and its sections over any $U$ are either empty or unique. SEG, and with it $S_u$, is a condition on stalks (local existence of a factorization), not a gluing condition, so $H^1$ detects nothing. With coefficients in calibrations, such as monotone rescalings of $u$, the parameter space of stable linear-Gaussian systems is star-shaped ($tA$ is stable for $t\in[0,1]$ whenever $A$ is), hence contractible, and the constant-coefficient $H^1$ vanishes as well. \stA

\subsection{The route that survives}\label{sec:im}
Let $\mathcal T$ be a declared class of interventions on one and the same bearer, for instance deformations of its couplings. \emph{Intervention monotonicity} is the requirement
\begin{equation}\label{eq:im}
\mathrm{IM}_{\mathcal T}:\qquad (s,t)\in\mathcal T,\ \ \Phi_Q(s)<\Phi_Q(t)\ \Longrightarrow\ u(s)\le u(t).
\end{equation}
For $u=f\circ\kappa$, the \emph{intervention digraph} has the realized coexperience classes, $\operatorname{im}\kappa$, as vertices, and an edge $\kappa(s)\to\kappa(t)$ for every $(s,t)\in\mathcal T$ with $\Phi_Q(s)<\Phi_Q(t)$.
\begin{theorem}[Exact frontier; \texttt{im\_iff\_edges}, \texttt{constant\_on\_scc}, \texttt{nonconstant\_if\_not\_scc}]\label{thm:frontier}
Let $u=f\circ\kappa$.
\begin{enumerate}
\item[(a)] $\mathrm{IM}_{\mathcal T}$ holds if and only if $f$ is monotone along the edges of the intervention digraph.
\item[(b)] Every such monotone $f$ is constant on each strongly connected component of the digraph.
\item[(c)] If $k=\kappa(s)$ and $k'=\kappa(s')$ are realized classes and $k$ is not reachable from $k'$, there is an $f$, monotone along the edges, with $f(k)\ne f(k')$, and hence $u(s)\ne u(s')$.
\end{enumerate}
All three parts are compiled in Lean. \stL
\end{theorem}
The route therefore escapes the collapse exactly when the intervention digraph on the realized classes is not strongly connected: when the interventions that raise $\Phi_Q$ cannot carry every realized coexperience class to every other one in both directions. A finite model escapes with a connected product image (\texttt{intervention\_monotonicity\_escapes}) \stL. The frontier is sharp in the other direction too: if equality is also required on $\Phi_Q$-preserving steps, and such steps chain together the realizations of a fibre, $S_u$ returns on that fibre (\texttt{equality\_on\_fibre\_steps\_collapses}) \stL.

\emph{Status.} $\mathrm{IM}_{\mathcal T}$ is \emph{not} $S_u$, and it is not original Q. It compares only states of one bearer that are related by declared interventions, and it compares them ordinally. It is a declared weakening of the order clause B to a class of interventions. In exchange, it needs neither a selection of representatives nor SEG, and it asks only for ordinal comparisons within a subject, not for absolute comparisons between different architectures.

\subsection{An instance in the band: monotonicity of \texorpdfstring{$J$}{J} under scaling}\label{sec:scaling}
\begin{proposition}[Monotonicity under scaling]\label{prop:scaling}
Let $n\ge2$ and $E\succ0$.
\begin{enumerate}
\item[(a)] \emph{Uniform scaling.} Let $\operatorname{spr}(A)<1$ and $t\in[0,1/\operatorname{spr}(A))$ (all $t\ge0$ if $\operatorname{spr}(A)=0$). Then $t\mapsto J(tA,E)$ is nondecreasing, with $J(0,E)=0$, and it is strictly increasing if the undirected support of $A$ is connected.
\item[(b)] \emph{Fixed scalar self-coupling.} Let $A(c)=aI+cW$ with $|a|<1$ and $W$ of zero diagonal, and let $c$ range over $\mathcal C=\{c\ge0:\operatorname{spr}(A(c))<1\}$, an interval containing $0$ (each eigenvalue $a+c\mu$ of $A(c)$ has $|a+c\mu|$ convex in $c$). Then $c\mapsto J(A(c),E)$ is nondecreasing on $\mathcal C$, and it is strictly increasing if the undirected support of $W$ is connected.
\end{enumerate}
In both cases, for $t>0$ and $c>0$ the support is that of $A$ (respectively $W$), and in a fixed representative the minimum cut grows as $t^2$ (respectively $c^2$). \stA\ \stM
\end{proposition}
\begin{proof}
(a) (1) $P_t=\sum_{j\ge0}t^{2j}A^jE(A^j)^\top$ solves $P_t=(tA)P_t(tA)^\top+E$, and $P_{t_2}-P_{t_1}=\sum_j(t_2^{2j}-t_1^{2j})A^jE(A^j)^\top\succeq0$ for $t_1\le t_2$: $P_t$ is nondecreasing in the Loewner order.
(2) Fix a bipartition $\pi$ and $t>0$. The map $B'\mapsto tB'$ is a bijection of $\mathfrak B_\pi$, and $E+(tA-tB')P_t(tA-tB')^\top=E+t^2(A-B')P_t(A-B')^\top=:R_t(B')$. For fixed $B'$ and $0<t_1<t_2$, with $C=A-B'$,
\begin{equation}\label{eq:resid-incr}
R_{t_2}(B')-R_{t_1}(B')=t_2^2\,C(P_{t_2}-P_{t_1})C^\top+(t_2^2-t_1^2)\,CP_{t_1}C^\top\succeq0,
\end{equation}
by (1) and because $X\mapsto CXC^\top$ preserves the Loewner order. Since $\log\det$ is monotone in the Loewner order, $f_t(B')=\tfrac12\log\det R_t(B')-\tfrac12\log\det E$ is nondecreasing in $t$, and so are its infimum over $B'$ and the minimum over $\pi$. At $t=0$, $B=0$ gives $J=0$, and $J\ge0$ always.
(3) If the support is connected, $A\notin\mathfrak B_\pi$ for every $\pi$, so $A-B'\ne0$ for every $B'\in\mathfrak B_\pi$. Let $0<t_1<t_2$, and let $B'_*$ minimize $f_{t_2}$ on $\mathfrak B_\pi$. It exists because $f_{t_2}$ is continuous and coercive: $f_t(B')\ge\tfrac12\log\bigl(1+t^2\lambda_{\min}(P_t)\lambda_{\max}(E)^{-1}\lVert A-B'\rVert_F^2\bigr)$, using $\det(I+X)\ge1+\operatorname{tr}X$ for $X\succeq0$. Put $M=t_1^2(A-B'_*)P_{t_1}(A-B'_*)^\top$, which is positive semidefinite and nonzero because $P_{t_1}\succ0$. By (1), $R_{t_2}(B'_*)\succeq E+(t_2/t_1)^2M$, and $\lambda\mapsto\log\det(E+\lambda M)$ is strictly increasing on $\lambda\ge0$, with derivative $\operatorname{tr}((E+\lambda M)^{-1}M)>0$. Hence, writing $J_\pi(tA)=\min_{B'\in\mathfrak B_\pi}f_t(B')$ for the inner minimum, $J_\pi(t_2A)=f_{t_2}(B'_*)>f_{t_1}(B'_*)\ge J_\pi(t_1A)$ for every $\pi$; taking $\pi$ minimizing at $t_2$ gives $J(t_2A)>J(t_1A)$. For $t_1=0$, $J(0)=0<J(t_2A)$ by Lemma~\ref{lem:J}.
(4) For $t>0$ the nonzero coefficients of $tA$ are those of $A$ multiplied by $t$; the support is unchanged and the crossing coefficients scale by $t$.

(b) Let $P(c)$ be the stationary covariance of $A(c)$ and $Q(c)=c^2P(c)$ (a matrix, unrelated to the theory Q), so that $Q=AQA^\top+c^2E$. Since $aI\in\mathfrak B_\pi$, for $c>0$ the map $B'\mapsto aI+cB'$ is a bijection of $\mathfrak B_\pi$, and $A(c)-(aI+cB')=c(W-B')$; the residual is $E+(W-B')Q(c)(W-B')^\top$. Differentiating $Q=AQA^\top+c^2E$ gives $Q'=AQ'A^\top+S$ with $S=WQA^\top+AQW^\top+2cE$, and the Lyapunov equation for $P$ turns this into
\begin{equation}\label{eq:S-identity}
S=c\bigl[(1-a^2)P+c^2WPW^\top+E\bigr]\succ0\qquad(c>0).
\end{equation}
Hence $Q'=\sum_{j\ge0}A^jS(A^j)^\top\succ0$, and $Q(c)$ is strictly increasing in the Loewner order. The rest of the argument is that of (a), with $Q(c)$ in place of $t^2P_t$ and $W-B'$ in place of $A-B'$: for fixed $B'$ the residual is nondecreasing in $c$, and for $0<c_1<c_2$ the residual at $c_2$ exceeds the residual at $c_1$ by $(W-B')\bigl(Q(c_2)-Q(c_1)\bigr)(W-B')^\top$, which is positive semidefinite and, when the support of $W$ is connected (so that $W-B'\ne0$ for every $B'\in\mathfrak B_\pi$), nonzero. The minimizer exists by coercivity, since $Q(c)\succ0$ for $c>0$. At $c=0$ the support is empty and $J=0$.
\end{proof}
{\sloppy \emph{Formal check} \stM. Both parts of the proposition are checked in Lean with Mathlib, over real matrices, as \texttt{QAuditScaling.prop\_scaling\_a} and \texttt{QAuditScaling.prop\_scaling\_b}; the other theorems named in this paragraph are in the same namespace, \texttt{QAuditScaling}, unless another one is given. The last sentence of the proposition is \texttt{supportGraph\_smul} and \texttt{supportGraph\_Acoup} for the support, and \texttt{wmin\_smul} and \texttt{wmin\_Acoup} for the minimum cut. The stationary covariance of $M=tA$ or $M=A(c)$ is the series $\sum_{j\ge0}M^jE(M^j)^\top$, which solves $P=MPM^\top+E$ (\texttt{Pt\_lyap}, \texttt{Pser\_lyap}) and whose convergence is proved from $\operatorname{spr}(M)<1$ by Gelfand's formula (\texttt{serSummable\_of\_spr\_lt\_one}). Here $\operatorname{spr}$ is Mathlib's spectral radius of $M$ regarded as a complex matrix; that it is the largest modulus of an eigenvalue is standard and is not formalized. The range of $t$ in (a) is written as $t\operatorname{spr}(A)<1$, using $\operatorname{spr}(tA)=t\operatorname{spr}(A)$ (\texttt{spr\_smul}). The inner minimum is an infimum, and it is attained (\texttt{exists\_minimizer}). The checked statement is slightly more general: (a) does not assume $\operatorname{spr}(A)<1$, and (b) does not assume that $W$ has zero diagonal. Neither $Q'=AQ'A^\top+S$ nor the identity~\eqref{eq:S-identity} is stated in Lean: the Lean proof of (b) replaces the derivative by the inequality $Q(c_2)-Q(c_1)\succeq(c_2-c_1)^2E$ for $c_1\le c_2$ in $\mathcal C$ (\texttt{Qb\_mono}). That $\mathcal C$ is an interval containing $0$ is checked with the argument given in (b), by \texttt{C\_ordConnected} and \texttt{C\_interval\_containing\_zero} in the namespace \texttt{QAuditScaling.AuditInterval}. \texttt{Scaling.lean} defines $\mathfrak B_\pi$, $J$, $w_{\min}$ and the support graph independently of the Lean statements of Lemma~\ref{lem:J}, and \texttt{QAuditBridge.J\_eq}, \texttt{QAuditBridge.wmin\_eq} and \texttt{QAuditBridge.supportGraph\_eq} prove that the two sets of definitions agree, those of $J$ for $E\succ0$ and $P\succeq0$.\par}

The identity~\eqref{eq:S-identity} is an exact polynomial identity modulo the Lyapunov equation, checked with generic symbolic $3\times3$ matrices (for general $n$ it is the computation in the proof \stA), and $Q'=AQ'A^\top+S$ holds exactly on a rational instance with symbolic $c$ (\texttt{escalado\_diagonal\_escalar.py}) \stS.

\begin{remark}[The diagonal must scale, or be scalar]\label{rem:fixed-diagonal}
Proposition~\ref{prop:scaling} does not extend to strengthening the couplings while a \emph{non-scalar} diagonal of self-couplings is held fixed; then $S$ can be indefinite. With $n=2$, $A(c)=D+cW$, $D=\mathrm{diag}(-0.23,\,0.31)$, $W=\bigl(\begin{smallmatrix}0&-0.19\\-0.83&0\end{smallmatrix}\bigr)$ and $E=\bigl(\begin{smallmatrix}4.61&2.59\\2.59&1.46\end{smallmatrix}\bigr)\succ0$, one has $J(A(0.33))\approx0.104083>J(A(0.38))\approx0.099702$ (\texttt{contraejemplo\_diagonal\_fija.py}: exact elimination of one variable and real roots to 50 digits) \stN. The four-unit band systems of Table~\ref{tab:networks}, with scalar self-coupling $a=0.2$ and all couplings equal to $c$, are covered by part~(b).
\end{remark}
\begin{corollary}[The band under scaling]\label{cor:band-scaling}
For every $n\ge4$, a stable dynamics whose undirected off-diagonal support is exactly the path $P_n$ (in~(b), $A=aI+cW$ with the undirected support of $W$ exactly $P_n$) stays in the band along the deformations of Proposition~\ref{prop:scaling} with $t>0$ (respectively $c>0$): its support is unchanged, so it remains connected without a universal vertex (\texttt{path\_reach\_from\_zero}, \texttt{path\_no\_universal}) \stL. Along these deformations $J$ increases strictly and the minimum cut grows quadratically \stA\ \stM. The same holds when the support is exactly the cycle $C_n$, $n\ge4$, whose maximum degree $2$ is smaller than $n-1$.
\end{corollary}
The corollary combines two checks, the graph part without Mathlib, on the Boolean adjacency of $P_n$, and the part on $J$ with Mathlib, through Proposition~\ref{prop:scaling}; it is not a single Lean statement.

Numerically, $J$ increases strictly on grids of $t$ (uniform scaling) for $P_4$, $P_5$, $P_6$, $C_4$, $C_5$ and five random connected systems with random $E$, with $w_{\min}(tA)=t^2w_{\min}(A)$ to machine precision (\texttt{monotonia\_escalado.py}); and on grids of $c$ (scalar self-coupling $a=0.2$) for $P_4$, $C_4$, $P_5$, $C_5$ and four random connected systems (\texttt{escalado\_diagonal\_escalar.py}, \texttt{banda\_monotonia.py}) \stN.

\emph{What $\mathrm{IM}_{\mathcal T}$ would then predict.} Take $\mathcal T$ to be one of these deformations of a band system, increasing in $t$ or in $c$. If $J$ were taken as the descriptor in $\mathrm{IM}_{\mathcal T}$, which this paper does not justify because $J$ is not $\Phi_Q$ (Section~\ref{sec:gauss}), $\mathrm{IM}_{\mathcal T}$ would predict that $u$ does not decrease as the couplings of the system are strengthened in that way. With the original descriptor the prediction also requires that $\Phi_Q$ increase along the deformation, which has not been shown. Every edge of the corresponding intervention digraph goes from a smaller to a larger parameter. The digraph on the realized classes is strongly connected only if coexperience classes recur in both directions along the deformation, and whether they do is an empirical, intra-subject question. Testing the prediction needs an ordinal intra-subject measure of $u$ along such interventions, that is, a calibrated measurement channel for $u$, which is not available: the stipulation A14--U of Definition~\ref{def:orderB} fixes a rule assigning observables to the degree, not a calibration.

\subsection{A methodological guide}\label{sec:guide}
The results of the three acts suggest a way of comparing theories of consciousness that neither presupposes $S_u$ nor collapses it.
\begin{enumerate}
\item \emph{Do not equate degrees across different architectures.} Static equality across systems with equal descriptor is exactly the obligation that the principles examined in Act~I cannot discharge without circularity, collapse, or a change of the meaning of the degree or of the domain.
\item \emph{Locate the disagreement first.} The structural presence bridges of IIT type and GNWT type commute on the central and disconnected levels and disagree exactly on the band (Section~\ref{sec:l1}). Systems in the band, such as $P_n$ and $C_n$ for $n\ge4$, are where the two kinds of proxy make different predictions.
\item \emph{Intervene within one system.} Use declared interventions, such as the deformations of Proposition~\ref{prop:scaling}, and ask only for ordinal comparisons of $u$ along them.
\item \emph{Check the digraph.} By Theorem~\ref{thm:frontier}, the comparison is informative exactly outside the strongly connected components of the intervention digraph on the realized classes; within a component, $u$ is forced to be constant.
\end{enumerate}

\subsection{What should be formalized next}\label{sec:next}
\begin{enumerate}
\item[(i)] In Lean~4 with Mathlib, the analytic statements that are not yet formalized: the hub bound of Theorem~\ref{thm:laxa}(d), the decomposition $I_{\rm pred}=I_{\rm cut}+J$ and the additivity of $I_{\rm pred}$ under $\otimes$ (Proposition~\ref{prop:monoidal}), which are paper proofs, and the statements about the restricted Gaussian family of Section~\ref{sec:gauss}, which are exact computations in SymPy or paper proofs (only the value $J=\tfrac12\log\tfrac54$ on the member $\theta=(\tfrac12,1,1)$ is checked with Mathlib, \texttt{QAudit.Audit.J\_ex\_eq}). Lemma~\ref{lem:J}, the min-cut sandwich (Theorem~\ref{thm:cheeger}) and Proposition~\ref{prop:scaling} are already checked with Mathlib \stM; what links them to the graph facts checked without Mathlib, such as the connectivity of $P_n$ in Corollary~\ref{cor:band-scaling} and the disconnected support of a tensor product, is still on paper, and so are the conventions of Remark~\ref{rem:small} for fewer than two units.
\item[(ii)] For a physically specified class $\mathcal T$ of deformations on $\mathbf{Open}_G$, the structure of strongly connected components of its intervention digraph on coexperience classes. It is the exact condition under which the route does not degenerate.
\item[(iii)] Not a theorem but its empirical prerequisite: a calibrated, ordinal, intra-subject measurement link for $u$, which the stipulation A14--U does not supply (Definition~\ref{def:orderB}(c)).
\end{enumerate}
Combining the zero-product monoid and the additive monoid of Section~\ref{sec:l2} into one structure with two operations only reorganizes the bookkeeping, and using $P_n$ as a test bench requires data on $u$ that do not exist; neither of these alternatives addresses $S_u$.


\section*{Use of generative AI and verification}\phantomsection\addcontentsline{toc}{section}{Use of generative AI and verification}\label{sec:ai}
This paper was drafted, and its Lean files and scripts were written and run, by an AI assistant (Claude Code, by Anthropic, with Anthropic's Claude models) at the request and under the direction of the author, who chose the questions, the scope and the structure of the argument and who is responsible for the content. Agents of the same assistant (separate sessions of it) read primary sources, verified the Gaussian family, translated the text of version~1.0 from its Spanish master version (the changes of version~1.1 were written in English), formalized the analytic results, and Theorem~\ref{thm:frontier} with real-valued degrees, in Lean with Mathlib, and reviewed the mathematics, the citations, the scope of the claims, the reproducibility of the archive, the formalization and the new material adversarially; their agreement is not an independent validation, and there has been no independent human review. The pre-submission review that led to version~1.1, namely its adversarial verification, the conditional theorem of Section~\ref{sec:conditional} and its Lean file, the Lean module with real-valued degrees for Theorem~\ref{thm:frontier} and the archived second check of Table~\ref{tab:networks}, was also done by agents of the same assistant, with no independent human review. During the research the author also used another AI system in exploratory conversations; that material was treated as unverified input, and no result of this paper depends on it.

Every statement labelled \stL\ was checked by the kernel of Lean~4~\cite{Lean4}: each of the seven Lean files compiles with exit code~0, contains no \texttt{sorry}, \texttt{admit}, \texttt{native\_decide}, \texttt{axiom} or \texttt{opaque} declarations, and depends at most on the standard axioms listed in Appendix~\ref{sec:archivos}. Every statement labelled \stM\ was checked by the kernel of Lean~4 with the Mathlib library~\cite{Mathlib2020} (v4.24.0), in the scope stated where it is cited: the eight modules of the Lean project \path{lean-mathlib} (\path{Core.lean}, \path{LemmaJ.lean}, \path{Scaling.lean}, \path{Bridge.lean}, \path{Interval.lean}, \path{ExamplesLemmaJ.lean}, \path{ExamplesScaling.lean} and \path{FrontierReal.lean}) compile with exit code~0 and contain no \texttt{sorry}, \texttt{admit}, \texttt{native\_decide}, \texttt{axiom} or \texttt{opaque} declarations, and none of their $104$ \texttt{\#print axioms} reports lists an axiom other than the standard axioms \texttt{propext}, \texttt{Classical.choice} and \texttt{Quot.sound} (Appendix~\ref{sec:archivos}). The exact computations labelled \stS\ were checked with SymPy~\cite{SymPy2017} for the Gaussian model, or by exact integer enumeration in Python (the band counts). Statements labelled \stA\ alone are paper proofs that have not been formalized, while \stA\ \stM\ marks a paper proof whose statement is also checked with Mathlib; statements labelled \stN\ are numerical illustrations, not certificates.

\appendix
\section{Correspondence between results, Lean theorems and scripts}\label{sec:correspondencia}
This appendix lists every numbered result of the paper (theorems, propositions, lemmas, corollaries, and the remarks that carry a claim) with its status labels (Section~\ref{sec:status}) and the Lean theorems or scripts that check it (Tables~\ref{tab:corr-1}, \ref{tab:corr-2} and~\ref{tab:corr-3} for Sections~\ref{sec:obligacion}, \ref{sec:categorica} and~\ref{sec:debil}, and Table~\ref{tab:corr-4} for Sections~\ref{sec:gauss} and~\ref{sec:via-futura}); ``paper proof'' means that no Lean file and no script checks it. Table~\ref{tab:corr-5} does the same for the statements that are checked by Lean or by a script but are not numbered results, so that each of the $57$ theorems of the Lean files without Mathlib whose axioms are printed in Appendix~\ref{sec:archivos}, each of the seven modules of the Lean project with Mathlib, through its main theorems, and each of the eleven scripts appears in at least one of these tables. Table~\ref{tab:corr-6} lists the definitions and hypotheses of Section~\ref{sec:obligacion}, and the two premises SEL and IPS$_K$ defined in later sections, with the places where each is used.

Lean theorems are named as in the six Lean files without Mathlib and the seven modules of the Lean project with Mathlib (Table~\ref{tab:files}). The theorems of the Mathlib project are written with their namespaces: \texttt{QAudit} (Lemma~\ref{lem:J} and Theorem~\ref{thm:cheeger}), \texttt{QAuditScaling} (Proposition~\ref{prop:scaling}), \texttt{QAuditBridge} and \texttt{QAuditScaling.\allowbreak AuditInterval}, with the examples that check non-vacuity in \texttt{QAudit.\allowbreak Audit} and \texttt{QAuditScaling.\allowbreak Audit}. The developments \texttt{QAudit} and \texttt{QAuditScaling} define $\mathfrak B_\pi$, $J$, the support graph and $w_{\min}$ independently; \texttt{QAuditBridge.\allowbreak J\_eq}, \texttt{QAuditBridge.\allowbreak supportGraph\_eq} and \texttt{QAuditBridge.\allowbreak wmin\_eq} prove that the two definitions agree (for $E\succ0$ and $P\succeq0$), so that both are about the one $J$ of Lemma~\ref{lem:J}. Scripts are named without their directory (\path{python/exact}, \path{python/high_precision} or \path{python/numeric}; Table~\ref{tab:files}). In Sections~\ref{sec:obligacion}--\ref{sec:debil} the status labels are implicit (Section~\ref{sec:status}); the tables state them explicitly by the rule given there: a statement that cites a Lean name of the files without Mathlib is \stL\ in the abstract form stated next to it, its link to $J$ is Lemma~\ref{lem:J}, \stA\ \stM\ for at least two units, and the analytic statements are \stA, with \stM\ where a theorem of the Mathlib project is cited. A label \stL\ means that Lean checks the logical form of a statement on abstract types or finite graphs, and a label \stM\ that Lean with Mathlib checks it over real matrices, in the scope stated in the row; no compilation decides whether a physical or phenomenal premise holds.

\begin{table}[htbp]\centering\small
\begin{tabular}{>{\raggedright\arraybackslash}p{0.27\textwidth}>{\raggedright\arraybackslash}p{0.09\textwidth}>{\raggedright\arraybackslash}p{0.54\textwidth}}
\toprule
Result & Status & Checked by\\
\midrule
Proposition~\ref{prop:provable} (provable descriptor equalities) & \stA & paper proof\\
Theorem~\ref{thm:coverage} (relational coverage) & \stA & paper proof; it is the case of Theorem~\ref{thm:impl} in which $R$ is an equivalence relation contained in $\ker\Phi$\\
Theorem~\ref{thm:continuum} (uncountably many classes in a fibre) & \stA & paper proof\\
Proposition~\ref{prop:extension} (form of any extension) & \stA & paper proof\\
Theorem~\ref{thm:threshold} (threshold reduction) & \stA & paper proof\\
Theorem~\ref{thm:impl} (entailment criterion) & \stL & \texttt{entails\_Su\_iff} (\texttt{IntentoSu.lean}, like the next four rows)\\
Theorem~\ref{thm:holgura} (slack and saturation) & \stL & \texttt{nonconstant\_iff}, \texttt{saturated}\\
Theorem~\ref{thm:dicotomia} (co-occurrence dichotomy) & \stL & \texttt{connected\_forces\_constant}; converse: \texttt{constant\_suff}\\
Theorem~\ref{thm:residuo} (residual under OI) & \stL & \texttt{oi\_residual}\\
Theorem~\ref{thm:robusta} (robust form) & \stL & \texttt{robust\_iff\_relational}\\
Proposition~\ref{prop:seg} (SEG and $S_u$) & \stL & (a) \texttt{seg\_implies\_Su}; (b) \texttt{segd\_iff} (\texttt{FiltracionCociente.lean}), where the predicate \texttt{SEGd} is Definition~\ref{def:SEG} and \texttt{SEG} is the split of part~(a)\\
\bottomrule
\end{tabular}
\caption{Numbered results of Section~\ref{sec:obligacion}, with their status and the Lean theorems that check them.}\label{tab:corr-1}
\end{table}

\begin{table}[htbp]\centering\small
\begin{tabular}{>{\raggedright\arraybackslash}p{0.27\textwidth}>{\raggedright\arraybackslash}p{0.09\textwidth}>{\raggedright\arraybackslash}p{0.54\textwidth}}
\toprule
Result & Status & Checked by\\
\midrule
Proposition~\ref{prop:funtor} (functors into a discrete category) & \stA & paper proof\\
Lemma~\ref{lem:J} (Gaussian integration and connectivity) & \stA, \stM & paper proof; for $n\ge2$, \texttt{QAudit.\allowbreak lemmaJ\_stationary}, with $P\succ0$ and $P=APA^\top+E$ given as data, and \texttt{QAudit.\allowbreak lemmaJ} and \texttt{QAudit.\allowbreak J\_eq\_zero\_iff} for every $P\succ0$ (together they give $J\ge0$); with $P$ constructed as the series $\sum_jA^jE(A^j)^\top$ when $\operatorname{spr}(A)<1$, \texttt{QAuditBridge.\allowbreak stationary\_of\_spr} and \texttt{QAuditBridge.\allowbreak lemmaJ\_spr}; the inner minimum is an infimum, attained by \texttt{QAudit.\allowbreak exists\_minimizer} and \texttt{QAudit.\allowbreak exists\_J\_minimizer}; $\det(I+Y)\ge1+\operatorname{tr}Y$ for $Y\succeq0$ is \texttt{QAudit.\allowbreak one\_add\_trace\_le\_det\_one\_add}; the convention $J=+\infty$ for one unit is not formalized; non-vacuity, $J=\tfrac12\log\tfrac54$ on a two-unit system, \texttt{QAudit.\allowbreak Audit.\allowbreak J\_ex\_eq} \stM; numerical illustration on the four networks, \path{J_conectividad.py} \stN; exact instance on $F_0$, $J=\tfrac12\log(1+k)$, \path{v12_familia_escala.py} \stS\\
Remark~\ref{rem:small} (systems with fewer than two units) & \stA & paper proof (conventions, and the invariance of $J$ under isomorphism); not formalized: in Lean an infimum over the empty set of bipartitions is $0$, so the Mathlib statements of Lemma~\ref{lem:J}, of Theorem~\ref{thm:cheeger} and of the strict parts of Proposition~\ref{prop:scaling} assume $n\ge2$\\
Proposition~\ref{prop:monoidal} (monoidal behavior) & \stA, \stL & \texttt{tensor\_no\_universal} (no hub on $X\otimes Y$); the rest is a paper proof; instances: $F_0\otimes F_0$, \path{gauss_conexiones.py} \stS, and a two-unit system tensored with a three-unit path, \path{pistas_numerico.py} \stN\\
Theorem~\ref{thm:translation} (translation criterion) & \stL & \texttt{nat\_iff\_functional} (\texttt{UnificacionCat.lean}), with $D_j$ nonempty (\texttt{Inhabited})\\
Theorem~\ref{thm:reconstruccion} (reconstruction) & \stL, \stA & criterion: \texttt{reconstruction\_iff\_injective}; (i): \texttt{blind\_spot}, \texttt{blind\_spot\_multiset}; paper: (ii), and the invariance of the degree multiset under isomorphism\\
Theorem~\ref{thm:bridge} (bridge consistency) & \stL, \stA & \texttt{bridge\_iff\_agree}; the uniqueness, which is immediate, is not part of the Lean statement\\
\bottomrule
\end{tabular}
\caption{Numbered results of Section~\ref{sec:categorica}, with their status and the Lean theorems or scripts that check them.}\label{tab:corr-2}
\end{table}

\begin{table}[htbp]\centering\small
\begin{tabular}{>{\raggedright\arraybackslash}p{0.27\textwidth}>{\raggedright\arraybackslash}p{0.09\textwidth}>{\raggedright\arraybackslash}p{0.54\textwidth}}
\toprule
Result & Status & Checked by\\
\midrule
Proposition~\ref{prop:decomposition} (extensive decomposition plus defect) & \stA & paper proof; exact instances on $F_0$, $F_{\rm mem}$ and $F_0\otimes F_0$, \path{gauss_conexiones.py} \stS; values on the networks of Table~\ref{tab:networks}, \path{pistas_numerico.py} \stN\\
Theorem~\ref{thm:cheeger} (min-cut sandwich) & \stA, \stM & paper proof; for $n\ge2$, exactly as stated, with the extreme eigenvalues and with $P\succ0$, $P=APA^\top+E$ given as data, \texttt{QAudit.\allowbreak sandwich}; with Loewner bounds $e_1I\preceq E\preceq e_2I$ and $P\preceq p_2I$ in place of the eigenvalues, \texttt{QAudit.\allowbreak sandwich\_lower} and \texttt{QAudit.\allowbreak sandwich\_upper}; with $P$ constructed as a series when $\operatorname{spr}(A)<1$, \texttt{QAuditBridge.\allowbreak sandwich\_spr}, and along the scaling $tA$, \texttt{QAuditBridge.\allowbreak sandwich\_along\_scaling}; on a two-unit system, where the lower bound is attained, \texttt{QAudit.\allowbreak Audit.\allowbreak sandwich\_ex} \stM; numerical check on the four networks and on twelve random systems (eleven non-degenerate), \path{pistas_numerico.py} \stN; exact check on $F_0$, where, for $c>0$, the lower bound is attained if and only if $s_2\le s_1$, \path{gauss_conexiones.py} \stS\\
Theorem~\ref{thm:laxa} (lax 2-cell) & \stL, \stA, \stM, \stS, \stN & (a) \texttt{universal\_connected}, with Lemma~\ref{lem:J} \stA\ \stM\ for $n\ge2$ and Remark~\ref{rem:small} \stA\ for $n=1$; (b) \texttt{n2\_equiv} (\texttt{GaussConexion.lean}), \texttt{n3\_equiv}; (c) \texttt{path\_reach\_from\_zero}, \texttt{path\_no\_universal}, \texttt{n4\_count}; the counts $0$, $0$, $15$ and $472$ for $n=2,\dots,5$ by exhaustive enumeration, \path{pistas_numerico.py} \stS; (d) paper proof \stA, with a numerical illustration on the star, \path{pistas_numerico.py} \stN\\
Theorem~\ref{thm:section} (factor through a section) & \stL & \texttt{section\_factor\_iff} (\texttt{PistasUnificacion.lean})\\
Theorem~\ref{thm:unif} (structural unification) & \stA, \stL, \stM, \stN & (i)--(iii): as the results they collect (Proposition~\ref{prop:decomposition}, Lemma~\ref{lem:J}, Proposition~\ref{prop:monoidal}, Theorems~\ref{thm:cheeger} and~\ref{thm:laxa}), with their combinatorial parts \stL, their parts checked with Mathlib \stM\ and numerical checks \stN; (iv): \texttt{square\_fails}, \texttt{values}, \texttt{no\_translation}, \texttt{blind\_spot}, \texttt{section\_factor\_iff}, \texttt{restriction\_vacuous}, \texttt{witness\_breaks\_purity} \stL, and Proposition~\ref{prop:extension} \stA\ for SEL against Rob\\
Theorem~\ref{thm:locus} (exact locus of non-commutation) & \stL, \stA, \stM & \texttt{noncomm\_locus}, \texttt{commute\_off\_band} (\texttt{FiltracionCociente.lean}), on Boolean predicates; their hypothesis, that a hub implies connectivity, is Theorem~\ref{thm:laxa}(a); that $J>0$ is connectivity on nonempty systems is Lemma~\ref{lem:J}, \stA\ \stM\ for at least two units, and Remark~\ref{rem:small} \stA\ for one unit\\
Theorem~\ref{thm:collapse} (collapse and zero-product monoid) & \stL, \stA, \stM & abstract form: \texttt{factors\_through\_collapse}, \texttt{zmul\_assoc} (with \texttt{zmul\_unit\_left}, \texttt{zmul\_unit\_right}) and \texttt{strong\_monoidal\_zero}, under its hypotheses \texttt{hunitL}, \texttt{hunitR}, \texttt{hne}, \texttt{hF0}, \texttt{hFne}, \texttt{hvan}; application to $J$ and to broadcast: Lemma~\ref{lem:J} \stA\ \stM, Theorem~\ref{thm:laxa}(a) and Proposition~\ref{prop:monoidal} \stA; instance on $F_0\otimes F_0$, \path{gauss_conexiones.py} \stS\\
Theorem~\ref{thm:selection} (variational selection in the fibres) & \stL & (a) \texttt{section\_factor\_iff}; (b) \texttt{restriction\_vacuous}; (c) \texttt{segd\_iff}, \texttt{seg\_implies\_Su}, \texttt{seg\_is\_Su}\\
Proposition~\ref{prop:conditional} (conditional unification) & \stA & paper proof, not formalized; its steps use \texttt{seg\_implies\_Su}, \texttt{section\_factor\_iff} and Theorem~\ref{thm:locus}\\
\bottomrule
\end{tabular}
\caption{Numbered results of Section~\ref{sec:debil}, with their status and the Lean theorems or scripts that check them.}\label{tab:corr-3}
\end{table}

\begin{table}[htbp]\centering\small
\begin{tabular}{>{\raggedright\arraybackslash}p{0.27\textwidth}>{\raggedright\arraybackslash}p{0.09\textwidth}>{\raggedright\arraybackslash}p{0.54\textwidth}}
\toprule
Result & Status & Checked by\\
\midrule
Proposition~\ref{prop:windows} (finite-window non-identifiability) & \stA, \stS & paper proof, general in $T$; the factorization of $\Sigma_T$, $\det\mathcal L_T=1$ and the value of $\det\Sigma_T$ for $T=0,\dots,5$, \path{v4_historias_ac.py} \stS\\
Proposition~\ref{prop:gauss-chains} (conditional constancy) & \stA, \stS, \stL & paper proof; the identities of (G7) and four exact chains, \path{v5_rkp_cadenas.py} \stS; an instance of \texttt{connected\_forces\_constant} \stL\\
Proposition~\ref{prop:gauss-plateau} (Plateau) & \stA, \stL & paper proof; only the consequence on the product domain, in the direction ``$S$ implies constancy'', is \texttt{plateau\_product} (\texttt{GaussConexion.lean}); the full support of the window laws is Proposition~\ref{prop:windows}\\
Theorem~\ref{thm:frontier} (exact frontier) & \stL & (a) \texttt{im\_iff\_edges}; (b) \texttt{constant\_on\_scc}, with \texttt{monotone\_on\_reach}; (c) \texttt{nonconstant\_if\_not\_scc} (\texttt{ViaFutura.lean}); in Lean the degrees take integer values, and the transcription to real values is not formalized (Section~\ref{sec:via-futura})\\
Proposition~\ref{prop:scaling} (monotonicity under scaling) & \stA, \stM, \stS, \stN & paper proof; (a) \texttt{QAuditScaling.\allowbreak prop\_scaling\_a}, on $\{t\ge0:t\operatorname{spr}(A)<1\}$ and without assuming $\operatorname{spr}(A)<1$; (b) \texttt{QAuditScaling.\allowbreak prop\_scaling\_b}, on $\mathcal C$ and without assuming that $W$ has zero diagonal; that $\mathcal C$ is an interval containing $0$, \texttt{QAuditScaling.\allowbreak AuditInterval.\allowbreak C\_interval\_containing\_zero} (with \texttt{QAuditScaling.\allowbreak AuditInterval.\allowbreak C\_ordConnected}); the support and the minimum cut, \texttt{QAuditScaling.\allowbreak supportGraph\_smul}, \texttt{QAuditScaling.\allowbreak supportGraph\_Acoup}, \texttt{QAuditScaling.\allowbreak wmin\_smul} and \texttt{QAuditScaling.\allowbreak wmin\_Acoup}; the covariance is the series $\sum_jM^jE(M^j)^\top$, convergent by \texttt{QAuditScaling.\allowbreak serSummable\_of\_spr\_lt\_one}, with $\operatorname{spr}$ Mathlib's spectral radius of the complexified matrix, and the inner minimum is attained, \texttt{QAuditScaling.\allowbreak exists\_minimizer}; in~(b) the Lean proof uses $Q(c_2)-Q(c_1)\succeq(c_2-c_1)^2E$, \texttt{QAuditScaling.\allowbreak Qb\_mono}, and states neither $Q'=AQ'A^\top+S$ nor identity~\eqref{eq:S-identity}; non-vacuity, \texttt{QAuditScaling.\allowbreak Audit.\allowbreak partA\_nonvacuous} and \texttt{QAuditScaling.\allowbreak Audit.\allowbreak partB\_nonvacuous} \stM; identity~\eqref{eq:S-identity}, with generic symbolic $3\times3$ matrices \stS\ (for general $n$, the computation in the proof \stA), and $Q'=AQ'A^\top+S$ on a rational instance, \path{escalado_diagonal_escalar.py} \stS; grids for (a), \path{monotonia_escalado.py}, and for (b), \path{escalado_diagonal_escalar.py} and \path{banda_monotonia.py} \stN\\
Remark~\ref{rem:fixed-diagonal} (the diagonal must scale, or be scalar) & \stN & \path{contraejemplo_diagonal_fija.py}: exact elimination of one variable and real roots to $50$ digits; not a certificate\\
Corollary~\ref{cor:band-scaling} (the band under scaling) & \stL, \stA, \stM, \stN & graph part: \texttt{path\_reach\_from\_zero}, \texttt{path\_no\_universal} \stL; the unchanged support, the strict increase of $J$ and the quadratic growth of the minimum cut, by Proposition~\ref{prop:scaling} \stA\ \stM\ (\texttt{QAuditScaling.\allowbreak supportGraph\_smul}, \texttt{QAuditScaling.\allowbreak supportGraph\_Acoup}, \texttt{QAuditScaling.\allowbreak prop\_scaling\_a}, \texttt{QAuditScaling.\allowbreak prop\_scaling\_b}, \texttt{QAuditScaling.\allowbreak wmin\_smul}, \texttt{QAuditScaling.\allowbreak wmin\_Acoup}); the corollary as one statement, which combines results of the files without Mathlib and of the Mathlib project, and the case of $C_n$ are \stA; grids, and $w_{\min}(tA)=t^2w_{\min}(A)$ to machine precision, \path{monotonia_escalado.py} \stN\\
\bottomrule
\end{tabular}
\caption{Numbered results of Sections~\ref{sec:gauss} and~\ref{sec:via-futura}, with their status and the Lean theorems or scripts that check them.}\label{tab:corr-4}
\end{table}

\begin{table}[htbp]\centering\small
\begin{tabular}{>{\raggedright\arraybackslash}p{0.36\textwidth}>{\raggedright\arraybackslash}p{0.09\textwidth}>{\raggedright\arraybackslash}p{0.45\textwidth}}
\toprule
Statement & Status & Checked by\\
\midrule
Countermodels CM1--CM4 (Section~\ref{sec:intento}) & \stL & \texttt{cm1\_not\_Su} (with \texttt{cm1\_OI}, \texttt{cm1\_Adm}), \texttt{cm2\_not\_Su} (with \texttt{cm2\_factor}), \texttt{cm3\_Su}, \texttt{cm3\_nonconstant}, \texttt{cm4\_not\_rel}\\
The bare blanket predicate, constant on $\mathbf{Open}$, differs from broadcast on the ring (Section~\ref{sec:proyecciones}) & \stL & \texttt{blanket\_vs\_broadcast}\\
No translation from positive integration to broadcast; none in the reverse direction (Section~\ref{sec:traduccion}) & \stL; \stA & \texttt{no\_translation}; the reverse direction follows from \texttt{values}, \texttt{nat\_iff\_functional} and Lemma~\ref{lem:J}, and is not a separate compiled statement\\
The cone of projections factors uniquely through the product (Section~\ref{sec:limite}) & \stL & \texttt{pairing\_unique}\\
The bridge square commutes on $\{K_{1,3},2K_2\}$ and not on $C_4$ or $P_4$; the values of $J$ in Table~\ref{tab:C-1} (Section~\ref{sec:puentes}) & \stL; \stN & \texttt{square\_commutes\_on\_sub}, \texttt{square\_fails}, \texttt{values}; \path{J_conectividad.py}, \path{pistas_numerico.py}\\
With OI$_\gamma$, fibre purity gives $S_u$; one non-isomorphic pair in a fibre destroys purity (Section~\ref{sec:via3}) & \stL & \texttt{purity\_gives\_Su}, \texttt{witness\_breaks\_purity}\\
Presence 2-cells into $\mathcal C_{J>0}$ are trivial (Section~\ref{sec:via4}) & \stL & \texttt{connected\_domain\_trivial}\\
On nonempty systems, $w_{\min}>0$ if and only if $J>0$ (Section~\ref{sec:l1}) & \stA, \stM & $n\ge2$: \texttt{QAuditBridge.\allowbreak J\_pos\_iff\_wmin\_pos}; $n=1$: Remark~\ref{rem:small} \stA\\
On the canonical representatives every function of the representative is a function of $\Phi$ (Section~\ref{sec:l4}) & \stL & \texttt{pullback\_determined}\\
Items (G1)--(G2) of Section~\ref{sec:gauss} & \stS & \path{v12_familia_escala.py}\\
Items (G3)--(G5) & \stS, \stA & \path{v3_modelos_apuntados.py}; written proof of (G3) \stA; for (G5) see connection (v)\\
Item (G6) & \stA, \stS & Proposition~\ref{prop:windows}; \path{v4_historias_ac.py}\\
Item (G7) & \stS & \path{v5_rkp_cadenas.py}\\
Connections (i)--(iv) of Section~\ref{sec:gauss} & \stS, \stL & \path{gauss_conexiones.py} (the analytic parts are Theorem~\ref{thm:cheeger} and Proposition~\ref{prop:gauss-plateau}); (i) \texttt{n2\_equiv}, \texttt{tensor\_no\_universal}, \texttt{f0\_hub\_readings}; (ii) \texttt{n2\_equiv}, \texttt{n3\_equiv}, \texttt{path\_reach\_from\_zero}, \texttt{path\_no\_universal}; (iii) \texttt{zmul\_assoc}, \texttt{strong\_monoidal\_zero}; (iv) \texttt{restriction\_vacuous}, \texttt{section\_factor\_iff}\\
Connection (v): the constancy results as instances; IPS$_K$ and SEG as reformulations & \stL & \texttt{connected\_forces\_constant}, \texttt{aligned\_premise\_iff}, \texttt{segd\_iff}\\
The four-unit networks: Table~\ref{tab:networks}, the hub bounds on the star and the tensor product (Section~\ref{sec:networks}) & \stN & \path{pistas_numerico.py}, \path{J_conectividad.py}\\
Exact covariance gives $S_u$ along chains and on components; under the order clause on all pairs, two classes that each occur strictly below the other have equal degrees; the $H^1$ obstruction is ill-posed (Section~\ref{sec:collapse}) & \stL; \stA & \texttt{exact\_covariance\_gives\_Su}, \texttt{exact\_covariance\_components}, \texttt{order\_all\_pairs\_general}, \texttt{order\_all\_pairs\_collapse} (integer-valued degrees in Lean); the paragraph on $H^1$ is a paper argument \stA\\
A finite model escapes the collapse; equality required on $\Phi_Q$-preserving steps that chain a fibre restores $S_u$ on it (Section~\ref{sec:im}) & \stL & \texttt{intervention\_monotonicity\_escapes}, \texttt{equality\_on\_fibre\_steps\_collapses}\\
\bottomrule
\end{tabular}
\caption{Statements checked by Lean or by a script that are not numbered results, with their location and status.}\label{tab:corr-5}
\end{table}

\begin{table}[htbp]\centering\small
\begin{tabular}{>{\raggedright\arraybackslash}p{0.30\textwidth}>{\raggedright\arraybackslash}p{0.62\textwidth}}
\toprule
Definition or hypothesis & Where it is used\\
\midrule
Definition~\ref{def:pcm} (realizations, pointed causal models, isomorphism) & Definitions~\ref{def:adm}, \ref{hyp:OI} and~\ref{hyp:Rob}; Proposition~\ref{prop:provable}; Proposition~\ref{prop:extension}, through part~(e); connection (iv) of Section~\ref{sec:gauss}\\
Definition~\ref{def:adm} (descriptors, Adm, dependence base) & Proposition~\ref{prop:provable}; the remarks after Theorem~\ref{thm:coverage} and after Definition~\ref{hyp:LSS}; Sections~\ref{sec:laguna} and~\ref{sec:intento}; Proposition~\ref{prop:funtor} (the remark after it); Proposition~\ref{prop:windows} (law-level descriptors) and connection (iv) of Section~\ref{sec:gauss}\\
Definition~\ref{def:phiQ} (Q, $D_0$, $\Phi_Q$, components, $S_x$) & every statement about $S_u$ \eqref{eq:Su}; part~(e), factorization, in Section~\ref{sec:laguna} and in Proposition~\ref{prop:conditional} (conditional unification); the architecture of Q in Section~\ref{sec:puentes}; $S_J$ in Definition~\ref{def:gauss-F0}(d); the other obligations of Q in Section~\ref{sec:mejor}\\
Definition~\ref{def:unity} (episodes, coexperience, $\mathfrak U$, $u$, $\psi$, $B_\psi$) & Definitions~\ref{hyp:UT}, \ref{hyp:pres}, \ref{hyp:N} and~\ref{def:SEG}; Theorem~\ref{thm:robusta}; Section~\ref{sec:intento}; Table~\ref{tab:B-1}; Proposition~\ref{prop:seg}; Sections~\ref{sec:l3} and~\ref{sec:l4}; connection (iv) of Section~\ref{sec:gauss}\\
Definition~\ref{def:orderB} (order clause B, A9--U, A14--U) & Theorem~\ref{thm:threshold} and the paragraph after it; Section~\ref{sec:mejor}; the discussion of Proposition~\ref{prop:gauss-chains}; Sections~\ref{sec:collapse}, \ref{sec:im} and~\ref{sec:scaling}\\
Definition~\ref{hyp:LSS} (LSS$_\Delta$) & Theorem~\ref{thm:coverage} (the remark after it); Proposition~\ref{prop:extension}(ii); Sections~\ref{sec:laguna}, \ref{sec:intento} and~\ref{sec:mejor}; the discussion of Proposition~\ref{prop:gauss-plateau}\\
Definition~\ref{hyp:OI} (OI$_\gamma$) & Theorem~\ref{thm:coverage} (the remark after it); Proposition~\ref{prop:extension}(ii); Section~\ref{sec:laguna}; the application of Theorem~\ref{thm:residuo} in Section~\ref{sec:intento}; Table~\ref{tab:B-1}; Sections~\ref{sec:mejor} and~\ref{sec:via3}; Table~\ref{tab:F-3}\\
Definition~\ref{hyp:Rob} (Rob) & Proposition~\ref{prop:extension}; Section~\ref{sec:laguna}; Table~\ref{tab:B-1}; Section~\ref{sec:mejor}; Definition~\ref{def:SEL} and Section~\ref{sec:via3}; Theorem~\ref{thm:unif}(iv); Section~\ref{sec:l3}\\
Definition~\ref{hyp:UT} (UT) & the paragraph after Theorem~\ref{thm:threshold}; Section~\ref{sec:intento}; Table~\ref{tab:B-1}; Section~\ref{sec:mejor}\\
Definition~\ref{hyp:pres} (B$_\psi^{\rm pres}$) & the remark after Definition~\ref{hyp:OI}; the paragraph after Theorem~\ref{thm:threshold}; Section~\ref{sec:intento}; Table~\ref{tab:B-1}; Section~\ref{sec:mejor}\\
Definition~\ref{hyp:N} (N) & the paragraph after Theorem~\ref{thm:threshold}; Section~\ref{sec:intento}; Table~\ref{tab:B-1}; Section~\ref{sec:mejor}\\
Definition~\ref{hyp:anchor} (absence anchor) & Theorem~\ref{thm:threshold}(v) and the paragraph after it; Section~\ref{sec:mejor}\\
Definition~\ref{def:TL} (threshold law TL) & Theorem~\ref{thm:threshold}(iii) and~(v) and the paragraph after it; Section~\ref{sec:mejor}\\
Definition~\ref{def:SEG} (SEG) & Proposition~\ref{prop:seg}; Figure~\ref{fig:graph}; Sections~\ref{sec:l3} and~\ref{sec:l4}; connection (v) of Section~\ref{sec:gauss}; Sections~\ref{sec:collapse} and~\ref{sec:im}\\
\midrule
Definition~\ref{def:SEL} (SEL, Section~\ref{sec:via3}) & Section~\ref{sec:via3}; Theorem~\ref{thm:unif}(iv); Section~\ref{sec:l3}\\
Definition~\ref{def:IPS} (IPS$_K$, Section~\ref{sec:gauss}) & item (G5); Proposition~\ref{prop:gauss-chains}; connection (v) of Section~\ref{sec:gauss}\\
\bottomrule
\end{tabular}
\caption{Definitions and hypotheses of Section~\ref{sec:obligacion}, and the two further premises defined in later sections, with the places where each is used. None of the hypotheses is adopted in this paper.}\label{tab:corr-6}
\end{table}

\clearpage
\subsection{Five candidate consequences of the categorical thesis}\label{sec:consecuencias}
One might draw five consequences from the categorical thesis of Section~\ref{sec:categorica}. Table~\ref{tab:F-3} states each one and what the results of this paper support; each entry rests only on the results and sections that it cites.
\begin{table}[htbp]\centering\small
\begin{tabular}{>{\raggedright\arraybackslash}p{0.22\textwidth}>{\raggedright\arraybackslash}p{0.70\textwidth}}
\toprule
Candidate consequence & What the results of the paper support\\
\midrule
1. The disputes between the theories come to an end & It does not follow. For the proxies studied, the framework locates the disagreement between the presence bridges of IIT type and GNWT type in squares that do not commute, exactly on the band of connected systems without a hub, such as $C_4$ and $P_4$ (Section~\ref{sec:puentes} and Theorem~\ref{thm:locus}). Whether the theories themselves disagree there is an empirical question (Section~\ref{sec:tesis}).\\
2. A dictionary between the theories given by natural transformations & Only where there is functional dependence: a translation exists if and only if the target projection is constant on the fibres of the source projection (\texttt{nat\_iff\_functional}, Section~\ref{sec:traduccion}), and there is none between positive integration and broadcast, in either direction. In general the translation is the relation $R_{ij}=\mathrm{im}\langle F_i,F_j\rangle$. The lax 2-cell of Theorem~\ref{thm:laxa} is an inequality between the presence predicates, strict for every $n\ge4$, together with lower bounds on $J$: it gives bounds, not a dictionary.\\
3. Criteria for consciousness in artificial systems & The projections give isomorphism invariants of a model (Section~\ref{sec:tesis}); which invariants are necessary for consciousness is a question about bridge laws that this paper does not settle, and Q does not fix whether $D_0$ contains artificial systems (Definition~\ref{def:phiQ}(a)). IIT and GNWT can also disagree when the grain or the substrate changes, for example on a workspace architecture implemented on a conventional computer (Section~\ref{sec:via2}, \emph{Scope}). Under OI$_\gamma$, which is not adopted, a realization isomorphic at grain $\gamma$ to a conscious one has the same total experience (Definition~\ref{hyp:OI}); so a criterion that denies experience to every artificial system is incompatible with OI$_\gamma$ as soon as some artificial system is isomorphic at grain $\gamma$ to a conscious realization. Choosing between such a criterion and OI$_\gamma$ is a decision, not a theorem.\\
4. A single clinical measure & It does not follow. The joint projection takes values in the product of the codomains of the projections, not on a single scale, and it has blind spots (Section~\ref{sec:limite} and Theorem~\ref{thm:reconstruccion}). A common scalar would need bridge laws, which this paper does not provide (strong thesis, Section~\ref{sec:tesis}), and a measurement of $u$ would need a calibrated channel, which is not available (Definition~\ref{def:orderB}(c) and Section~\ref{sec:mejor}).\\
5. Falsifiability & In principle, but not with what is available. Refuting SEG, the testable form of non-degenerate $S_u$, would need two original realizations with the same computed value of $\Phi_Q$, coexperience structures determined by a calibrated channel and a specified $\psi$, and none of the three is available (Section~\ref{sec:mejor}). The band is where the two kinds of proxy make different predictions (Section~\ref{sec:guide}). The prediction of Section~\ref{sec:scaling}, that $u$ does not decrease along the deformations of Proposition~\ref{prop:scaling}, is conditional on $\mathrm{IM}_{\mathcal T}$ with $J$ in place of $\Phi_Q$, which is not justified, and testing it needs an ordinal intra-subject measure of $u$, a calibrated measurement channel for $u$, which is not available (A14--U is a stipulated rule, not a calibration).\\
\bottomrule
\end{tabular}
\caption{Five consequences that one might draw from the categorical thesis, and what the results of this paper support.}\label{tab:F-3}
\end{table}

\paragraph{Three tempting formalizations that are vacuous.} Three ways of formalizing these consequences come to mind and prove nothing.
\begin{itemize}
\item Defining the GNWT-type predicate as $J>\theta$, for a threshold $\theta\ge0$, makes the 2-cell into the IIT-type predicate $J>0$ true by definition, since $J>\theta$ implies $J>0$. It is a higher threshold on the same integration and no longer involves broadcast, whereas the 2-cell of Theorem~\ref{thm:laxa} relates a broadcast hub to $J$.
\item A canonical section of $\Phi$ obtained from a choice function that does not use the free energy is an arbitrary choice of representatives, and no section, whatever its rule of choice, gives $S_u$: by \texttt{section\_factor\_iff} (Section~\ref{sec:via3}), the values at the representatives define a factor of $u$ if and only if $S_u$ already holds.
\item A uniqueness statement of the form ``there is exactly one $x$ with $x=c$'' holds for every $c$; it says nothing about the uniqueness of a minimizer of the free energy in a fibre, which is not guaranteed (Section~\ref{sec:via3}).
\end{itemize}
None of the three removes the uncountably many isomorphism classes that the fibre of $\Phi_Q$ through $s_0$ contains under Rob when $\theta\mapsto\Phi_Q(s_\theta)$ is continuous on the convex set $P$ of Rob and $\Phi_Q(s_0)$ is interior to its range (Theorem~\ref{thm:continuum} and Proposition~\ref{prop:extension}(i)). The formulation with content is the physical selection law SEL (Definition~\ref{def:SEL}), which contradicts Rob under the same conditions (Section~\ref{sec:via3}). Instead of these three, the paper uses the lax 2-cell of Theorem~\ref{thm:laxa}, whose combinatorial parts are compiled in Lean (Table~\ref{tab:corr-3}), and three theorems of \texttt{PistasUnificacion.lean}, \texttt{section\_factor\_iff}, \texttt{purity\_gives\_Su} and \texttt{witness\_breaks\_purity}, which state what a selection in the fibres gives and what it does not give (Section~\ref{sec:via3}).
\clearpage

\section{Files, execution log and axioms}\label{sec:archivos}
The archive that accompanies this paper contains the source of the paper, the seven Lean files without Mathlib, the Lean project with Mathlib, the scripts and their expected outputs, and three scripts that rerun everything and compare the results (\path{lean/verify_lean.sh}, \path{lean-mathlib/verify_mathlib.sh} and \path{python/verify_python.sh}). A master script at its root, \path{verify_all.sh} (also run by \texttt{make verify}), first checks the SHA-256 sums of all the files of the archive (\path{SHA256SUMS}) and then runs the three in turn; the Lean project with Mathlib, which downloads about 1~GB and needs about 7~GB of disk, can be skipped with \texttt{SKIP\_MATHLIB=1}. Its \texttt{README.md} explains how to use them, and Table~\ref{tab:files} lists what each file checks. File and script names are those of the original working files, partly in Spanish, and so are some comments and printed messages in the scripts.

\paragraph{Tools.} Lean~4.24.0 (commit \texttt{797c613e}), without Mathlib for the seven Lean files, and with Mathlib~v4.24.0 for the eight modules of the Lean project, built with Lake (\path{lean-mathlib/lake-manifest.json} pins Mathlib at commit \texttt{f897ebcf}); Python~3.13 with NumPy~2.5~\cite{NumPy2020}, SciPy~1.18~\cite{SciPy2020}, SymPy~1.14~\cite{SymPy2017} and mpmath~1.3~\cite{mpmath}. Python is run as \texttt{python3 -I} (isolated mode).

\begin{table}[htbp]\centering\footnotesize
\begin{tabular}{>{\raggedright\arraybackslash}p{0.36\textwidth}>{\raggedright\arraybackslash}p{0.56\textwidth}}
\toprule
File & What it checks, and how\\
\midrule
\path{lean/IntentoSu.lean} & Theorems~\ref{thm:impl}--\ref{thm:robusta} and the countermodels CM1--CM4. \stL\\
\path{lean/UnificacionCat.lean} & Pairing, translation, bridges, reconstruction and the four-unit example (including \texttt{blind\_spot\_multiset}). \stL\\
\path{lean/PistasUnificacion.lean} & Combinatorial and logical parts of the four routes: hub implies connectivity, $n=3$ and $n=4$ counts, paths without hub, tensor products, sections and purity. \stL\\
\path{lean/FiltracionCociente.lean} & The four lines: locus of non-commutation, zero-product monoid, collapse of the ideal, vacuity of restriction, SEG (\texttt{segd\_iff}) and pullback. \stL\\
\path{lean/ViaFutura.lean} & Collapse of exact covariance and of the global order clause; intervention monotonicity; Theorem~\ref{thm:frontier} with integer degrees, (b) and~(c) for an abstract digraph. \stL\\
\path{lean/GaussConexion.lean} & Logical skeleton of the Gaussian connections: Plateau on a product domain, aligned premise, $n=2$, the two readings of the hub in $F_0$. \stL\\
\path{lean/CondicionalQ.lean} & Theorem~\ref{thm:conditional-discharge} and the paragraph on $S_q$ after it, in abstract form; non-vacuity. \stL\\
\path{lean-mathlib/QAuditAnalytic/Core.lean} & Matrix inequalities over $\mathbb R$: $\det(I+X)\ge1+\operatorname{tr}X$ and $\log\det(I+X)\le\operatorname{tr}X$ for $X\succeq0$, trace bounds under Loewner bounds, and the extreme eigenvalues $\lambda_{\min}$, $\lambda_{\max}$. \stM\\
\path{lean-mathlib/QAuditAnalytic/LemmaJ.lean} & Lemma~\ref{lem:J} and Theorem~\ref{thm:cheeger} for $n\ge2$, with $P\succ0$ and $P=APA^\top+E$ given as data, also with Loewner bounds in place of the eigenvalues; attainment of the minima. \stM\\
\path{lean-mathlib/QAuditAnalytic/Scaling.lean} & Proposition~\ref{prop:scaling}(a) and~(b), with the covariance as a series that converges when $\operatorname{spr}<1$; support and minimum cut under scaling. \stM\\
\path{lean-mathlib/QAuditAnalytic/Bridge.lean} & The two developments define the same $J$, support graph and $w_{\min}$; Lemma~\ref{lem:J} and Theorem~\ref{thm:cheeger} with the covariance constructed from $\operatorname{spr}(A)<1$; $J>0$ if and only if $w_{\min}>0$. \stM\\
\path{lean-mathlib/QAuditAnalytic/Interval.lean} & The set $\mathcal C$ of Proposition~\ref{prop:scaling}(b) is an interval containing $0$. \stM\\
\path{lean-mathlib/QAuditAnalytic/ExamplesLemmaJ.lean}, \path{ExamplesScaling.lean} & Non-vacuity: $J=\tfrac12\log\tfrac54$ on a two-unit system, where the eigenvalue form of the sandwich applies and its lower bound is attained; $\operatorname{spr}(I)=1$, and the range $[0,\tfrac12)$ for $2I$; the strict parts of Proposition~\ref{prop:scaling} on instances; outside the domain where the covariance series converges Lean gives $J$ the default value $0$, so monotonicity on $[0,\infty)$ fails there and the restriction of the domain is needed. \stM\\
\path{lean-mathlib/QAuditAnalytic/FrontierReal.lean} & Theorem~\ref{thm:frontier} and the sentence after it with real-valued degrees, as instances for degrees in an arbitrary order, with the digraph on $\operatorname{im}\kappa$; which order hypotheses (b) and (c) use; reachability and the digraph copied from \path{ViaFutura.lean}; non-vacuity. \stM\\
\path{python/exact/v12_familia_escala.py}, \path{v3_modelos_apuntados.py}, \path{v4_historias_ac.py}, \path{v5_rkp_cadenas.py} & The restricted Gaussian family: items (G1)--(G2), (G3)--(G5), (G6) and (G7) respectively, which the scripts number (1)--(2), (3), (4) and (5); 31, 30, 33 and 32 exact checks, all passed. \stS\\
\path{python/exact/gauss_conexiones.py} & Connections of $F_0$ with the theorems: 30 exact checks, all passed. \stS\\
\path{python/exact/escalado_diagonal_escalar.py} & Identity~\eqref{eq:S-identity} and $Q'=AQ'A^\top+S$ \stS; monotonicity in $c$ on grids \stN.\\
\path{python/high_precision/contraejemplo_diagonal_fija.py} & Remark~\ref{rem:fixed-diagonal}: exact elimination and roots to 50 digits \stN; $\det S<0$ \stS.\\
\path{python/numeric/J_conectividad.py}, \path{pistas_numerico.py} & $J$ of the four networks; sandwich, hub bound and tensor products (Table~\ref{tab:networks}) \stN; exact counts of the band for $n=2,\dots,5$ \stS.\\
\path{python/numeric/segundo_optimizador.py} & The second check of Table~\ref{tab:networks}: BFGS with restarts and an alternating weighted least-squares scheme, from $B=A-A_{\rm cross}(\pi)$; the value kept and each method's. \stN\\
\path{python/numeric/banda_monotonia.py}, \path{monotonia_escalado.py} & Monotonicity of $J$ along couplings with scalar self-coupling, and under uniform scaling. \stN\\
\path{lean/verify_lean.sh}, \path{lean-mathlib/verify_mathlib.sh}, \path{python/verify_python.sh}, \path{verify_all.sh}, \path{Makefile} & Rerun the checks: the Lean scripts refuse forbidden constructs, compile (with Lake for Mathlib) and compare the \texttt{\#print axioms} output with \path{expected/}; the Python script compares the outputs with \path{python/expected/} (\path{compare_outputs.py}); \path{verify_all.sh} (\texttt{make verify}) checks \path{SHA256SUMS} first and runs the three.\\
\bottomrule
\end{tabular}
\caption{Files of the archive and what each one checks.}
\label{tab:files}
\end{table}

\paragraph{Execution log.} Every Lean file without Mathlib and every module of the Lean project with Mathlib compiles with exit code~0; none contains \texttt{sorry}, \texttt{native\_decide}, \texttt{admit}, an axiom declaration or an \texttt{opaque} declaration. Every script ends with a line stating whether all its checks passed, and all of them pass. One run is worth recording: the first run of \path{gauss_conexiones.py} failed on a rational grid because of an error in the script, not in the mathematics (a Python integer division produced floating-point numbers); the script was corrected and rerun.

\paragraph{Axioms per theorem} (output of \texttt{\#print axioms} for the 85 theorems that the seven Lean files without Mathlib print: 57 in the first six files of Table~\ref{tab:files}, listed first, and 28 in \path{CondicionalQ.lean}, listed after them). The five auxiliary lemmas that Appendix~\ref{sec:correspondencia} cites without a printed report, \texttt{cm1\_OI}, \texttt{cm1\_Adm}, \texttt{cm2\_factor}, \texttt{zmul\_unit\_left} and \texttt{zmul\_unit\_right}, depend on no axioms (checked by adding \texttt{\#print axioms} for them; these reports are not among the expected outputs). No theorem depends on any axiom outside the three standard ones.
\begin{itemize}\raggedright
\item \emph{No axioms} (33):
\texttt{saturated}, \texttt{connected\_forces\_constant}, \texttt{constant\_suff}, \texttt{robust\_iff\_relational}, \texttt{cm3\_Su}, \texttt{pairing\_unique}, \texttt{bridge\_iff\_agree}, \texttt{values}, \texttt{square\_fails}, \texttt{no\_translation}, \texttt{square\_commutes\_on\_sub}, \texttt{blind\_spot}, \texttt{blind\_spot\_multiset}, \texttt{blanket\_vs\_broadcast}, \texttt{universal\_connected}, \texttt{n3\_equiv}, \texttt{n4\_count}, \texttt{section\_factor\_iff}, \texttt{purity\_gives\_Su}, \texttt{witness\_breaks\_purity}, \texttt{connected\_domain\_trivial}, \texttt{zmul\_assoc}, \texttt{restriction\_vacuous}, \texttt{seg\_implies\_Su}, \texttt{segd\_iff}, \texttt{seg\_is\_Su}, \texttt{pullback\_determined}, \texttt{equality\_on\_fibre\_steps\_collapses}, \texttt{im\_iff\_edges}, \texttt{plateau\_product}, \texttt{aligned\_premise\_iff}, \texttt{n2\_equiv}, \texttt{f0\_hub\_readings}.
\item \emph{Only \texttt{propext}} (13):
\texttt{entails\_Su\_iff}, \texttt{cm1\_not\_Su}, \texttt{cm2\_not\_Su}, \texttt{cm3\_nonconstant}, \texttt{cm4\_not\_rel}, \texttt{tensor\_no\_universal}, \texttt{noncomm\_locus}, \texttt{commute\_off\_band}, \texttt{factors\_through\_collapse}, \texttt{order\_all\_pairs\_general}, \texttt{intervention\_monotonicity\_escapes}, \texttt{monotone\_on\_reach}, \texttt{constant\_on\_scc}.
\item \emph{\texttt{propext} and \texttt{Quot.sound}} (3):
\texttt{exact\_covariance\_gives\_Su}, \texttt{exact\_covariance\_components}, \texttt{order\_all\_pairs\_collapse}.
\item \emph{\texttt{propext}, \texttt{Classical.choice} and \texttt{Quot.sound}} (8):
\texttt{nonconstant\_iff}, \texttt{oi\_residual}, \texttt{nat\_iff\_functional}, \texttt{reconstruction\_iff\_injective}, \texttt{path\_reach\_from\_zero}, \texttt{path\_no\_universal}, \texttt{strong\_monoidal\_zero}, \texttt{nonconstant\_if\_not\_scc}.
\end{itemize}
In \path{CondicionalQ.lean} (namespace \texttt{CondicionalQ}):
\begin{itemize}\raggedright
\item \emph{No axioms} (12):
\texttt{suff\_of\_cover}, \texttt{aligned\_cover\_iff}, \texttt{pointed\_reading\_kernel}, \texttt{Sq\_pointed}, \texttt{rel\_suff\_all\_psi}, \texttt{cover\_structures}, \texttt{robust\_iff\_relational}, \texttt{chain\_eqv\_iff}, \texttt{suff\_and\_B\_iff\_nonstrict}, \texttt{factor\_unique\_on\_image}, \texttt{orderB\_iff\_monotone\_factor}, \texttt{descent\_componentwise}.
\item \emph{Only \texttt{propext}} (2):
\texttt{single\_psi\_example}, \texttt{conditional\_discharge\_nonvacuous}.
\item \emph{Only \texttt{Quot.sound}} (2):
\texttt{descent\_of\_suff}, \texttt{descent\_iff\_factor}.
\item \emph{\texttt{propext}, \texttt{Classical.choice} and \texttt{Quot.sound}} (12):
\texttt{coverage\_iff}, \texttt{order\_cover\_B\_iff}, \texttt{order\_cover\_SuB\_iff}, \texttt{order\_cover\_Su\_iff}, \texttt{theorem19\_from\_order}, \texttt{suff\_iff\_factor}, \texttt{suffB\_iff\_monotone\_factor}, \texttt{im\_Su\_iff\_injective}, \texttt{im\_B\_iff}, \texttt{descent\_needs\_invariance}, \texttt{descent\_exact\_iff}, \texttt{conditional\_discharge}.
\end{itemize}


\paragraph{Axioms of the Lean project with Mathlib.} The eight modules print $104$ reports of \texttt{\#print axioms}: $14$ in \path{Core.lean}, $9$ in \path{LemmaJ.lean}, $27$ in \path{Scaling.lean}, $9$ in \path{Bridge.lean}, $2$ in \path{Interval.lean}, $5$ in \path{ExamplesLemmaJ.lean}, $7$ in \path{ExamplesScaling.lean} and $31$ in \path{FrontierReal.lean}. Of these, $92$ are \texttt{[propext, Classical.choice, Quot.sound]}, the three standard axioms, and the other twelve, all in \path{FrontierReal.lean}, depend on no axioms, except \texttt{QAuditFrontier.\allowbreak b\_needs\_transitivity}, which depends on \texttt{propext} alone. All $104$ are listed, with the names of their theorems, in \path{lean-mathlib/expected/}, against which \path{verify_mathlib.sh} compares the output of each module. The main ones, by module, are the following.
\begin{itemize}\raggedright
\item \path{Core.lean}: \texttt{QAudit.one\_add\_trace\_le\_det\_one\_add}, \texttt{QAudit.log\_det\_one\_add\_le\_trace}, \texttt{QAudit.lamMin\_isLeast\_spectrum}, \texttt{QAudit.lamMax\_isGreatest\_spectrum}.
\item \path{LemmaJ.lean} (all nine): \texttt{QAudit.connected\_iff\_forall\_cross}, \texttt{QAudit.lemmaJ}, \texttt{QAudit.J\_eq\_zero\_iff}, \texttt{QAudit.lemmaJ\_stationary}, \texttt{QAudit.sandwich\_lower}, \texttt{QAudit.sandwich\_upper}, \texttt{QAudit.sandwich}, \texttt{QAudit.exists\_minimizer}, \texttt{QAudit.exists\_J\_minimizer}.
\item \path{Scaling.lean}: \texttt{QAuditScaling.exists\_minimizer}, \texttt{QAuditScaling.supportGraph\_smul}, \texttt{QAuditScaling.supportGraph\_Acoup}, \texttt{QAuditScaling.Qb\_mono}, \texttt{QAuditScaling.serSummable\_of\_spr\_lt\_one}, \texttt{QAuditScaling.prop\_scaling\_a}, \texttt{QAuditScaling.prop\_scaling\_b}, \texttt{QAuditScaling.wmin\_smul}, \texttt{QAuditScaling.wmin\_Acoup}.
\item \path{Bridge.lean} (all nine): \texttt{QAuditBridge.J\_eq}, \texttt{QAuditBridge.supportGraph\_eq}, \texttt{QAuditBridge.wmin\_eq}, \texttt{QAuditBridge.Jt\_pos\_of\_connected}, \texttt{QAuditBridge.sandwich\_along\_scaling}, \texttt{QAuditBridge.stationary\_of\_spr}, \texttt{QAuditBridge.lemmaJ\_spr}, \texttt{QAuditBridge.sandwich\_spr}, \texttt{QAuditBridge.J\_pos\_iff\_wmin\_pos}.
\item \path{Interval.lean} (both): \texttt{QAuditScaling.AuditInterval.C\_ordConnected}, \texttt{QAuditScaling.AuditInterval.C\_interval\_containing\_zero}.
\item \path{ExamplesLemmaJ.lean} and \path{ExamplesScaling.lean}: \texttt{QAudit.Audit.J\_ex\_eq}, \texttt{QAudit.Audit.sandwich\_ex}, \texttt{QAuditScaling.Audit.partA\_nonvacuous}, \texttt{QAuditScaling.Audit.partB\_nonvacuous}, \texttt{QAuditScaling.Audit.not\_monotone\_on\_Ici}.
\item \path{FrontierReal.lean}: \texttt{QAuditFrontier.frontier\_real}, \texttt{QAuditFrontier.frontier\_c\_uniform\_real}, \texttt{QAuditFrontier.frontier\_paper\_form}, \texttt{QAuditFrontier.escape\_iff\_real}, \texttt{QAuditFrontier.escape\_paper\_form}, \texttt{QAuditFrontier.empty\_domain\_convention}, \texttt{QAuditFrontier.im\_iff\_edges}, \texttt{QAuditFrontier.constant\_on\_scc}, \texttt{QAuditFrontier.nonconstant\_if\_not\_scc}, \texttt{QAuditFrontier.nonconstant\_uniform}, \texttt{QAuditFrontier.chain\_iff\_reflTransGen}, \texttt{QAuditFrontier.b\_transitive\_antisymmetric}, \texttt{QAuditFrontier.b\_needs\_antisymmetry}, \texttt{QAuditFrontier.b\_needs\_transitivity}, \texttt{QAuditFrontier.c\_ne\_of\_two\_values}, \texttt{QAuditFrontier.c\_ne\_without\_strict\_pair}, \texttt{QAuditFrontier.c\_fails\_discrete\_order}, \texttt{QAuditFrontier.Example.nonvacuous}.
\end{itemize}


\begin{thebibliography}{99}\raggedright

\bibitem{Bayne2010}
T.~Bayne, \emph{The Unity of Consciousness}, Oxford University Press, Oxford (2010).
\href{https://doi.org/10.1093/acprof:oso/9780199215386.001.0001}{https://doi.org/10.1093/acprof:oso/9780199215386.001.0001}

\bibitem{BayneChalmers2003}
T.~Bayne and D.~J.~Chalmers, ``What is the unity of consciousness?'', in A.~Cleeremans (ed.), \emph{The Unity of Consciousness: Binding, Integration, and Dissociation}, Oxford University Press, Oxford (2003), pp.~23--58.
\href{https://doi.org/10.1093/acprof:oso/9780198508571.003.0002}{https://doi.org/10.1093/acprof:oso/9780198508571.003.0002}

\bibitem{Capucci2022}
M.~Capucci, B.~Gavranovi\'c, J.~Hedges and E.~F.~Rischel, ``Towards foundations of categorical cybernetics'', \emph{Electronic Proceedings in Theoretical Computer Science} \textbf{372}, 235--248 (2022).
\href{https://doi.org/10.4204/EPTCS.372.17}{https://doi.org/10.4204/EPTCS.372.17}

\bibitem{Chalmers1995}
D.~J.~Chalmers, ``Absent qualia, fading qualia, dancing qualia'', in T.~Metzinger (ed.), \emph{Conscious Experience}, Imprint Academic / Sch\"oningh, Thorverton and Paderborn (1995), pp.~309--328.

\bibitem{Cogitate2025}
Cogitate Consortium, O.~Ferrante, U.~Gorska-Klimowska, S.~Henin et al., ``Adversarial testing of global neuronal workspace and integrated information theories of consciousness'', \emph{Nature} \textbf{642}(8066), 133--142 (2025).
\href{https://doi.org/10.1038/s41586-025-08888-1}{https://doi.org/10.1038/s41586-025-08888-1}

\bibitem{DehaeneChangeux2011}
S.~Dehaene and J.-P.~Changeux, ``Experimental and theoretical approaches to conscious processing'', \emph{Neuron} \textbf{70}(2), 200--227 (2011).
\href{https://doi.org/10.1016/j.neuron.2011.03.018}{https://doi.org/10.1016/j.neuron.2011.03.018}

\bibitem{Friston2013}
K.~Friston, ``Life as we know it'', \emph{Journal of the Royal Society Interface} \textbf{10}(86), 20130475 (2013).
\href{https://doi.org/10.1098/rsif.2013.0475}{https://doi.org/10.1098/rsif.2013.0475}

\bibitem{Fritz2020}
T.~Fritz, ``A synthetic approach to Markov kernels, conditional independence and theorems on sufficient statistics'', \emph{Advances in Mathematics} \textbf{370}, 107239 (2020).
\href{https://doi.org/10.1016/j.aim.2020.107239}{https://doi.org/10.1016/j.aim.2020.107239}

\bibitem{GrazianoKastner2011}
M.~S.~A.~Graziano and S.~Kastner, ``Human consciousness and its relationship to social neuroscience: a novel hypothesis'', \emph{Cognitive Neuroscience} \textbf{2}(2), 98--113 (2011).
\href{https://doi.org/10.1080/17588928.2011.565121}{https://doi.org/10.1080/17588928.2011.565121}

\bibitem{HameroffPenrose2014}
S.~Hameroff and R.~Penrose, ``Consciousness in the universe: a review of the `Orch OR' theory'', \emph{Physics of Life Reviews} \textbf{11}(1), 39--78 (2014).
\href{https://doi.org/10.1016/j.plrev.2013.08.002}{https://doi.org/10.1016/j.plrev.2013.08.002}

\bibitem{NumPy2020}
C.~R.~Harris, K.~J.~Millman, S.~J.~van~der~Walt et al., ``Array programming with NumPy'', \emph{Nature} \textbf{585}(7825), 357--362 (2020).
\href{https://doi.org/10.1038/s41586-020-2649-2}{https://doi.org/10.1038/s41586-020-2649-2}

\bibitem{HurewiczWallman1941}
W.~Hurewicz and H.~Wallman, \emph{Dimension Theory}, Princeton Mathematical Series \textbf{4}, Princeton University Press, Princeton (1941).
\href{https://doi.org/10.1515/9781400875665}{https://doi.org/10.1515/9781400875665}

\bibitem{Kitazono2018}
J.~Kitazono, R.~Kanai and M.~Oizumi, ``Efficient algorithms for searching the minimum information partition in integrated information theory'', \emph{Entropy} \textbf{20}(3), 173 (2018).
\href{https://doi.org/10.3390/e20030173}{https://doi.org/10.3390/e20030173}. Preprint \href{https://arxiv.org/abs/1712.06745}{arXiv:1712.06745}

\bibitem{KleinerTull2021}
J.~Kleiner and S.~Tull, ``The mathematical structure of integrated information theory'', \emph{Frontiers in Applied Mathematics and Statistics} \textbf{6}, 602973 (2021).
\href{https://doi.org/10.3389/fams.2020.602973}{https://doi.org/10.3389/fams.2020.602973}

\bibitem{Mathlib2020}
The mathlib Community, ``The Lean mathematical library'', in \emph{Proceedings of the 9th ACM SIGPLAN International Conference on Certified Programs and Proofs (CPP 2020)}, ACM, New York (2020), pp.~367--381.
\href{https://doi.org/10.1145/3372885.3373824}{https://doi.org/10.1145/3372885.3373824}

\bibitem{Mendola2025}
J.~Mendola, \emph{The Neural Structure of Consciousness}, Cambridge Studies in Philosophy, Cambridge University Press, Cambridge (2025).
\href{https://doi.org/10.1017/9781009514774}{https://doi.org/10.1017/9781009514774}

\bibitem{SymPy2017}
A.~Meurer, C.~P.~Smith, M.~Paprocki et al., ``SymPy: symbolic computing in Python'', \emph{PeerJ Computer Science} \textbf{3}, e103 (2017).
\href{https://doi.org/10.7717/peerj-cs.103}{https://doi.org/10.7717/peerj-cs.103}

\bibitem{Lean4}
L.~de~Moura and S.~Ullrich, ``The Lean~4 theorem prover and programming language'', in A.~Platzer and G.~Sutcliffe (eds.), \emph{Automated Deduction -- CADE 28}, Lecture Notes in Computer Science \textbf{12699}, Springer, Cham (2021), pp.~625--635.
\href{https://doi.org/10.1007/978-3-030-79876-5_37}{https://doi.org/10.1007/978-3-030-79876-5\_37}

\bibitem{mpmath}
The mpmath development team, \emph{mpmath: a Python library for arbitrary-precision floating-point arithmetic}, version~1.3.0 (2023).
\href{https://mpmath.org/}{https://mpmath.org/}

\bibitem{NiuSpivak2023}
N.~Niu and D.~I.~Spivak, \emph{Polynomial Functors: A Mathematical Theory of Interaction}, London Mathematical Society Lecture Note Series \textbf{498}, Cambridge University Press, Cambridge (2025).
\href{https://doi.org/10.1017/9781009576734}{https://doi.org/10.1017/9781009576734}

\bibitem{Northoff2019}
G.~Northoff, N.~Tsuchiya and H.~Saigo, ``Mathematics and the brain: a category theoretical approach to go beyond the neural correlates of consciousness'', \emph{Entropy} \textbf{21}(12), 1234 (2019).
\href{https://doi.org/10.3390/e21121234}{https://doi.org/10.3390/e21121234}

\bibitem{Oizumi2014}
M.~Oizumi, L.~Albantakis and G.~Tononi, ``From the phenomenology to the mechanisms of consciousness: Integrated Information Theory 3.0'', \emph{PLoS Computational Biology} \textbf{10}(5), e1003588 (2014).
\href{https://doi.org/10.1371/journal.pcbi.1003588}{https://doi.org/10.1371/journal.pcbi.1003588}

\bibitem{Oizumi2016PNAS}
M.~Oizumi, N.~Tsuchiya and S.~Amari, ``Unified framework for information integration based on information geometry'', \emph{Proceedings of the National Academy of Sciences} \textbf{113}(51), 14817--14822 (2016).
\href{https://doi.org/10.1073/pnas.1603583113}{https://doi.org/10.1073/pnas.1603583113}

\bibitem{Schechter2018}
E.~Schechter, \emph{Self-Consciousness and ``Split'' Brains: The Minds' I}, Oxford University Press, Oxford (2018).
\href{https://doi.org/10.1093/oso/9780198809654.001.0001}{https://doi.org/10.1093/oso/9780198809654.001.0001}

\bibitem{SignorelliWangCoecke2021}
C.~M.~Signorelli, Q.~Wang and B.~Coecke, ``Reasoning about conscious experience with axiomatic and graphical mathematics'', \emph{Consciousness and Cognition} \textbf{95}, 103168 (2021).
\href{https://doi.org/10.1016/j.concog.2021.103168}{https://doi.org/10.1016/j.concog.2021.103168}

\bibitem{TokerSommer2019}
D.~Toker and F.~T.~Sommer, ``Information integration in large brain networks'', \emph{PLoS Computational Biology} \textbf{15}(2), e1006807 (2019).
\href{https://doi.org/10.1371/journal.pcbi.1006807}{https://doi.org/10.1371/journal.pcbi.1006807}

\bibitem{TononiKoch2015}
G.~Tononi and C.~Koch, ``Consciousness: here, there and everywhere?'', \emph{Philosophical Transactions of the Royal Society B} \textbf{370}(1668), 20140167 (2015).
\href{https://doi.org/10.1098/rstb.2014.0167}{https://doi.org/10.1098/rstb.2014.0167}

\bibitem{Tsuchiya2016}
N.~Tsuchiya, S.~Taguchi and H.~Saigo, ``Using category theory to assess the relationship between consciousness and integrated information theory'', \emph{Neuroscience Research} \textbf{107}, 1--7 (2016).
\href{https://doi.org/10.1016/j.neures.2015.12.007}{https://doi.org/10.1016/j.neures.2015.12.007}

\bibitem{TullKleiner2021}
S.~Tull and J.~Kleiner, ``Integrated information in process theories: towards categorical IIT'', \emph{Journal of Cognitive Science} \textbf{22}(2), 92--123 (2021).
\href{https://arxiv.org/abs/2002.07654}{arXiv:2002.07654}

\bibitem{SciPy2020}
P.~Virtanen, R.~Gommers, T.~E.~Oliphant et al., ``SciPy 1.0: fundamental algorithms for scientific computing in Python'', \emph{Nature Methods} \textbf{17}(3), 261--272 (2020).
\href{https://doi.org/10.1038/s41592-019-0686-2}{https://doi.org/10.1038/s41592-019-0686-2}

\bibitem{Zurek2009}
W.~H.~Zurek, ``Quantum Darwinism'', \emph{Nature Physics} \textbf{5}(3), 181--188 (2009).
\href{https://doi.org/10.1038/nphys1202}{https://doi.org/10.1038/nphys1202}

\end{thebibliography}

\end{document}
